% Perturbations du système immunitaire : Le VIH/Sida — SVTEEHB 3ème — Cours signé OnBuch+
% Programme officiel MINESEC. Compiler avec : tectonic N07-perturbations-systeme-immunitaire-vih-sida.tex
\newcommand{\DOCMATIERE}{SVTEEHB}
\newcommand{\DOCNIVEAU}{3ème}
\newcommand{\DOCMODULE}{SVTEEHB – classe de 3ème}
\newcommand{\DOCLECON}{N07}
\newcommand{\DOCTITRE}{Perturbations du système immunitaire : Le VIH/Sida}
% OnBuch+ — préambule « cours signés » ENRICHI (compilé avec tectonic / XeLaTeX)
% Système visuel « Pop » d'OnBuch+ : fond crème (jamais blanc), bordures encre
% épaisses, ombres « dures » décalées, titres Archivo Black en capitales.
% Version MATHÉMATIQUES : graphes (pgfplots/tikz), boîtes pédagogiques
% (définition, propriété, méthode, attention, le savais-tu, exercice résolu,
% exercice, TP, bilan), page de garde et signature OnBuch+ ancrée en bas.
%
% Macros attendues dans main.tex AVANT \input{preamble} :
%   \DOCMATIERE  \DOCNIVEAU  \DOCMODULE  \DOCLECON (numéro)  \DOCTITRE
\documentclass[11pt]{article}

\usepackage{fontspec}
\usepackage[french]{babel}
\usepackage[a4paper,margin=2cm,top=3.4cm,bottom=2.8cm,headheight=1.4cm,footskip=1.3cm]{geometry}
\usepackage{amsmath,amssymb,amsfonts}
\usepackage{mathtools} % \xmapsto, \coloneqq... souvent utilisés par les modèles
\usepackage{cancel} % simplifications barrées (bilans de forces)
% \abs{z} (module, valeur absolue) et \norm{u} (norme), utilisés par les modèles
\providecommand{\abs}{}\let\abs\relax\DeclarePairedDelimiter{\abs}{\lvert}{\rvert}
\providecommand{\norm}{}\let\norm\relax\DeclarePairedDelimiter{\norm}{\lVert}{\rVert}
\usepackage[table]{xcolor} % [table] : fournit aussi \rowcolors (alternance de lignes)
\usepackage{pagecolor}
\usepackage{tikz}
\usetikzlibrary{arrows.meta,positioning,calc,shapes.geometric,shapes.misc,shapes.arrows,decorations.pathmorphing,decorations.pathreplacing,decorations.markings,patterns,shadows,mindmap,trees,fit,backgrounds,matrix,intersections,angles,quotes}
\usepackage{pgfplots}
\usepackage[version=4]{mhchem} % équations nucléaires \ce{^{235}_{92}U}
\usepackage[european,siunitx]{circuitikz} % schémas électriques
\pgfplotsset{compat=1.18}
\usepgfplotslibrary{fillbetween,groupplots} % aires entre courbes (calcul intégral)
\usepackage{enumitem}
\usepackage{fancyhdr}
% [most] charge toutes les bibliothèques tcolorbox (listings, minted, raster,
% documentation...) et gonfle inutilement la mémoire TeX sur un document long
% (~300 boîtes) : on ne charge que ce qui est réellement utilisé.
\usepackage[skins,breakable]{tcolorbox}
\usepackage{graphicx}
\usepackage{colortbl}
\usepackage{booktabs}
\usepackage{tabularx}
\usepackage{multirow}
\usepackage{diagbox} % case d'en-tête diagonale des tableaux à double entrée
% Colonnes extensibles centrée / gauche / droite (C, L, R), souvent utilisées par les modèles
\newcolumntype{C}{>{\centering\arraybackslash}X}
\newcolumntype{L}{>{\raggedright\arraybackslash}X}
\newcolumntype{R}{>{\raggedleft\arraybackslash}X}
\usepackage{array}
\usepackage{multicol}
\usepackage{siunitx}
% parse-numbers=false : accepte aussi \SI{2,5 \times 10^{-3}}{...}
\sisetup{locale=FR,output-decimal-marker={,},group-minimum-digits=4,parse-numbers=false}
\DeclareSIUnit\ohm{\ensuremath{\symup{\Omega}}} % U+2126 absent de la police
\DeclareSIUnit\division{div} % réglages d'oscilloscope (ms/div, V/div)
\DeclareSIUnit\tr{tr} % tours (vitesse de rotation, ex. \SI{1500}{\tr\per\minute})
\DeclareSIUnit\molar{M} % concentration molaire (ex. \SI{3}{\molar}, \SI{1}{\micro\molar})
\DeclareSIUnit\an{an} % année (le modèle confond parfois avec \year, un compteur TeX) — ex. \SI{5}{\kilo\meter\per\an}
\DeclareSIUnit\FCFA{FCFA} % franc CFA (ex. \SI{500}{\FCFA\per\kilo\gram})
\DeclareSIUnit\degreeBrix{\textdegree Brix} % teneur en sucre d'un sirop/jus (réfractométrie)
\DeclareSIUnit\month{mois} % le modèle utilise parfois \month, un compteur TeX, comme unité de durée
\DeclareSIUnit\microliter{\micro\liter} % le modèle écrit parfois \microliter en un seul token
\DeclareSIUnit\million{millions} % dénombrement (ex. \SI{15}{\million\per\mL} spermatozoïdes)
\usepackage{qrcode}
% \degree : commande fréquemment utilisée par les modèles pour un angle ou une
% température en dehors de \SI{}{} (non fournie par les paquets chargés).
\providecommand{\degree}{^{\circ}}
% \qed (fin de démonstration), utilisé par les modèles sans amsthm
\providecommand{\qed}{\hfill$\square$}
% \barre : double barre des valeurs interdites dans un tableau de variations (tkz-tab)
\providecommand{\barre}{\ensuremath{\|}}
% environnement proof (amsthm non chargé) utilisé par les modèles
\ifdefined\proof\else\newenvironment{proof}[1][Démonstration]{\par\noindent\textit{#1.}\ }{\qed\par}\fi
\usepackage{lastpage}
\usepackage{hyperref}
\hypersetup{hidelinks}

% ============================================================
% Palette « Pop » OnBuch+ — mêmes hex que la classe P (plus_ui.dart)
% ============================================================
\definecolor{poporange}{HTML}{F59321}
\definecolor{popdark}{HTML}{E07A0C}
\definecolor{popcream}{HTML}{FFF6E8}
\definecolor{popink}{HTML}{1C1714}
\definecolor{poppurple}{HTML}{7A5AE0}
\definecolor{popgreen}{HTML}{2E9E6A}
\definecolor{poppink}{HTML}{E0457B}
\definecolor{popblue}{HTML}{2F7DE1}
\definecolor{popgold}{HTML}{C98A12}
\definecolor{popmuted}{HTML}{8A7F76}
\definecolor{popline}{HTML}{EADFD2}
\definecolor{popgray}{HTML}{8A7F76}
\colorlet{popgrey}{popgray}
% Alias tolérants (le modèle nomme parfois les couleurs différemment)
\colorlet{popwhite}{white}
\colorlet{popblack}{popink}
\colorlet{popred}{poppink}
\colorlet{popyellow}{popgold}
% Teintes claires (fonds de tableaux/encadrés) — le modèle nomme parfois
% les couleurs avec un suffixe L (ex. popblueL) sans les avoir définies.
\colorlet{poporangeL}{poporange!20}
\colorlet{popdarkL}{popdark!20}
\colorlet{poppurpleL}{poppurple!20}
\colorlet{popgreenL}{popgreen!20}
\colorlet{poppinkL}{poppink!20}
\colorlet{popblueL}{popblue!20}
\colorlet{popgoldL}{popgold!20}
\colorlet{popredL}{popred!20}
\colorlet{popgrayL}{popgray!20}
% Teintes claires pour fonds de tableaux / zones
\colorlet{popblueL}{popblue!10!white}
\colorlet{poporangeL}{poporange!14!white}
\colorlet{popgreenL}{popgreen!12!white}
\colorlet{poppurpleL}{poppurple!12!white}
\colorlet{poppinkL}{poppink!12!white}
\colorlet{popgoldL}{popgold!14!white}

\pagecolor{popcream}

% ============================================================
% Typographie
% ============================================================
\defaultfontfeatures{Ligatures=TeX}
\setmainfont{PlusJakartaSans-Medium.ttf}[Path=../../../fonts/,BoldFont=PlusJakartaSans-Bold.ttf]
% Maths en sans-serif (Fira Math) avec chiffres et lettres droites de Plus
% Jakarta, pour que les formules se fondent dans le texte.
\usepackage[math-style=ISO,bold-style=ISO]{unicode-math}
\setmathfont{FiraMath-Regular.otf}[Path=../../../fonts/]
\setmathfont{PlusJakartaSans-Medium.ttf}[Path=../../../fonts/,range={up/num,up/{Latin,latin},\mathup}]
\usepackage{icomma} % virgule décimale sans espace en mode math
% Caractères mathématiques Unicode absents des polices de texte (Plus Jakarta,
% Archivo Black) : rendus par leur équivalent mathématique, en texte comme en
% formule — sinon ils s'affichent en carré vide (ex. « Arithmétique dans ℤ »).
\usepackage{newunicodechar}
\newunicodechar{ℝ}{\ensuremath{\mathbb{R}}}
\newunicodechar{ℤ}{\ensuremath{\mathbb{Z}}}
\newunicodechar{ℂ}{\ensuremath{\mathbb{C}}}
\newunicodechar{ℕ}{\ensuremath{\mathbb{N}}}
\newunicodechar{ℚ}{\ensuremath{\mathbb{Q}}}
\newunicodechar{𝔻}{\ensuremath{\mathbb{D}}}
\newunicodechar{α}{\ensuremath{\alpha}}
\newunicodechar{β}{\ensuremath{\beta}}
\newunicodechar{γ}{\ensuremath{\gamma}}
\newunicodechar{δ}{\ensuremath{\delta}}
\newunicodechar{ε}{\ensuremath{\varepsilon}}
\newunicodechar{θ}{\ensuremath{\theta}}
\newunicodechar{λ}{\ensuremath{\lambda}}
\newunicodechar{μ}{\ensuremath{\mu}}
\newunicodechar{σ}{\ensuremath{\sigma}}
\newunicodechar{τ}{\ensuremath{\tau}}
\newunicodechar{φ}{\ensuremath{\varphi}}
\newunicodechar{ψ}{\ensuremath{\psi}}
\newunicodechar{ω}{\ensuremath{\omega}}
\newunicodechar{Σ}{\ensuremath{\Sigma}}
% majuscules produites par \MakeUppercase dans les titres (α-aminés -> Α)
\newunicodechar{Α}{\ensuremath{\alpha}}
\newunicodechar{Β}{\ensuremath{\beta}}
\newunicodechar{Γ}{\ensuremath{\Gamma}}
\newunicodechar{Δ}{\ensuremath{\Delta}}
\newunicodechar{Ε}{\ensuremath{\varepsilon}}
\newunicodechar{Θ}{\ensuremath{\Theta}}
\newunicodechar{Λ}{\ensuremath{\Lambda}}
\newunicodechar{Μ}{\ensuremath{\mu}}
\newunicodechar{Τ}{\ensuremath{\tau}}
\newunicodechar{Φ}{\ensuremath{\Phi}}
\newunicodechar{Ψ}{\ensuremath{\Psi}}
\newunicodechar{Ω}{\ensuremath{\Omega}}
\newunicodechar{↦}{\ensuremath{\mapsto}}
\newunicodechar{∈}{\ensuremath{\in}}
\newunicodechar{∉}{\ensuremath{\notin}}
\newunicodechar{∧}{\ensuremath{\wedge}}
\newunicodechar{⇔}{\ensuremath{\Leftrightarrow}}
\newunicodechar{⟺}{\ensuremath{\Longleftrightarrow}}
\newunicodechar{⇒}{\ensuremath{\Rightarrow}}
\newunicodechar{⟹}{\ensuremath{\Longrightarrow}}
\newunicodechar{∘}{\ensuremath{\circ}}
\newunicodechar{∪}{\ensuremath{\cup}}
\newunicodechar{∩}{\ensuremath{\cap}}
\newunicodechar{≡}{\ensuremath{\equiv}}
\newunicodechar{⊂}{\ensuremath{\subset}}
\newunicodechar{⊥}{\ensuremath{\perp}}
\newunicodechar{∃}{\ensuremath{\exists}}
\newunicodechar{∀}{\ensuremath{\forall}}
\newunicodechar{∛}{\ensuremath{\sqrt[3]{\,}}}
\newunicodechar{∜}{\ensuremath{\sqrt[4]{\,}}}
\newunicodechar{⃗}{\ensuremath{{}^{\to}}}
\newunicodechar{ⁿ}{\ifmmode^{n}\else\textsuperscript{\ensuremath{n}}\fi}
\newunicodechar{ˣ}{\ifmmode^{x}\else\textsuperscript{\ensuremath{x}}\fi}
\newunicodechar{ᵇ}{\ifmmode^{b}\else\textsuperscript{\ensuremath{b}}\fi}
\newunicodechar{ʳ}{\ifmmode^{r}\else\textsuperscript{\ensuremath{r}}\fi}
\newunicodechar{ᵘ}{\ifmmode^{u}\else\textsuperscript{\ensuremath{u}}\fi}
\newunicodechar{ʸ}{\ifmmode^{y}\else\textsuperscript{\ensuremath{y}}\fi}
\newunicodechar{ᵃ}{\ifmmode^{a}\else\textsuperscript{\ensuremath{a}}\fi}
\newunicodechar{ᵉ}{\ifmmode^{e}\else\textsuperscript{\ensuremath{e}}\fi}
\newunicodechar{ᵏ}{\ifmmode^{k}\else\textsuperscript{\ensuremath{k}}\fi}
\newunicodechar{ᵖ}{\ifmmode^{p}\else\textsuperscript{\ensuremath{p}}\fi}
\newunicodechar{ᴺ}{\ifmmode^{N}\else\textsuperscript{\ensuremath{N}}\fi}
\newunicodechar{ᶜ}{\ifmmode^{c}\else\textsuperscript{\ensuremath{c}}\fi}
\newunicodechar{ʰ}{\ifmmode^{h}\else\textsuperscript{\ensuremath{h}}\fi}
\newunicodechar{ᵠ}{\ifmmode^{q}\else\textsuperscript{\ensuremath{q}}\fi}
\newunicodechar{ₙ}{\ifmmode_{n}\else\textsubscript{\ensuremath{n}}\fi}
\newunicodechar{ₐ}{\ifmmode_{a}\else\textsubscript{\ensuremath{a}}\fi}
\newunicodechar{ᵢ}{\ifmmode_{i}\else\textsubscript{\ensuremath{i}}\fi}
\newunicodechar{ᵣ}{\ifmmode_{r}\else\textsubscript{\ensuremath{r}}\fi}
\newunicodechar{ₖ}{\ifmmode_{k}\else\textsubscript{\ensuremath{k}}\fi}
\newunicodechar{ᵦ}{\ifmmode_{\beta}\else\textsubscript{\ensuremath{\beta}}\fi}
\newunicodechar{ₓ}{\ifmmode_{x}\else\textsubscript{\ensuremath{x}}\fi}
\newfontfamily\popdisplay{ArchivoBlack-Regular.ttf}[Path=../../../fonts/]
\newfontfamily\poplogo{SpaceGrotesk-Bold.ttf}[Path=../../../fonts/]
\newfontfamily\popsemi{PlusJakartaSans-SemiBold.ttf}[Path=../../../fonts/]
\newfontfamily\popxbold{PlusJakartaSans-ExtraBold.ttf}[Path=../../../fonts/]
\newfontfamily\popxboldls{PlusJakartaSans-ExtraBold.ttf}[Path=../../../fonts/,LetterSpace=3.0]

\setlist[enumerate]{leftmargin=1.7em,itemsep=5pt,topsep=4pt}
\setlist[itemize]{leftmargin=1.7em,itemsep=4pt,topsep=4pt,label={\color{poporange}\textbullet}}
\setlength{\parindent}{0pt}
\setlength{\parskip}{5pt}
\renewcommand{\arraystretch}{1.25}

% pgfplots : style Pop par défaut
\pgfplotsset{
  every axis/.append style={axis line style={popink,line width=1pt},tick style={popink},
    grid style={popline,line width=0.5pt},label style={font=\small},tick label style={font=\footnotesize}},
  popcurve/.style={poporange,line width=1.8pt,no markers,smooth},
}
% Même style disponible dans un \draw TikZ simple (hors pgfplots)
\tikzset{popcurve/.style={poporange,line width=1.8pt}}
\tikzset{popgrid/.style={popline,very thin}}
\colorlet{popcurve}{poporange} % le nom du style est parfois utilisé comme couleur

% ============================================================
% Pill / squiggle
% ============================================================
\newcommand{\poppill}[3]{%
  \tikz[baseline=-0.55ex]\node[draw=popink,line width=1.2pt,rounded corners=2.6pt,%
    fill=#2,inner xsep=6.5pt,inner ysep=3pt]{%
    \popxboldls\fontsize{7}{7}\selectfont\color{#3}\MakeUppercase{#1}};%
}
\newcommand{\popsquiggle}[1]{%
  \tikz[baseline=-0.4ex]{
    \draw[#1,line width=2.6pt,line cap=round]
      plot [smooth] coordinates {(0,0.09) (0.34,0.01) (0.68,0.21) (1.02,0.01) (1.36,0.10)};
  }%
}
% Mot-clé surligné (marqueur orange sous le texte)
\newcommand{\cle}[1]{{\popxbold\color{popdark}#1}}

% ============================================================
% En-tête / pied
% ============================================================
\pagestyle{fancy}
\fancyhf{}
\renewcommand{\headrulewidth}{0pt}
\renewcommand{\footrulewidth}{0pt}
\lhead{\raisebox{-0.15ex}{\poplogo\fontsize{13}{13}\selectfont\color{popink}OnBuch\color{poporange}+}}
\rhead{\poppill{\DOCMATIERE\ \textbullet\ \DOCNIVEAU\ \textbullet\ Leçon \DOCLECON}{white}{popink}}
\lfoot{\popxbold\fontsize{7.6}{7.6}\selectfont\color{popmuted}\MakeUppercase{Cours signé OnBuch+}\ \textbullet\ {\popsemi\DOCTITRE}}
\rfoot{\tikz[baseline=-0.6ex]\node[rounded corners=8.5pt,fill=popink,minimum height=17pt,minimum width=17pt,inner xsep=5pt,inner sep=0]{\popxbold\fontsize{8}{8}\selectfont\color{popcream}\thepage\,/\,\pageref*{LastPage}};}

% ============================================================
% Titres
% ============================================================
\newcommand{\chaptitre}[2]{%
  \begin{tcolorbox}[enhanced,colback=poporange,colframe=popink,boxrule=2.6pt,arc=11pt,
      drop shadow=popink,left=15pt,right=15pt,top=12pt,bottom=11pt,before skip=3pt,after skip=18pt]
    \poppill{#2}{popink}{popcream}\par\vspace{9pt}
    {\raggedright\hyphenpenalty=10000\exhyphenpenalty=10000\popdisplay\fontsize{22}{24}\selectfont\color{popcream}\MakeUppercase{#1}\par}
  \end{tcolorbox}
}
% Section numérotée : pastille orange avec le numéro + titre Archivo
\newcounter{popsec}
\newcommand{\coursec}[1]{%
  \stepcounter{popsec}\setcounter{popsub}{0}%
  \par\vspace{16pt}\noindent
  \begin{minipage}{\linewidth}
  \tikz[baseline=-0.7ex]\node[draw=popink,line width=1.6pt,fill=poporange,rounded corners=4pt,minimum size=22pt,inner sep=0]{\popdisplay\fontsize{12}{12}\selectfont\color{popcream}\thepopsec};%
  \hspace{8pt}{\popdisplay\fontsize{15}{17}\selectfont\color{popink}\MakeUppercase{#1}}\par
  \vspace{4pt}\noindent{\color{poporange}\rule{36pt}{3.6pt}}
  \end{minipage}\par\nopagebreak\vspace{8pt}
}
\newcounter{popsub}
\newcommand{\courssub}[1]{%
  \stepcounter{popsub}%
  \par\vspace{9pt}\noindent{\popxbold\fontsize{11.5}{13}\selectfont\color{popink}{\color{poporange}\thepopsec.\thepopsub}\hspace{6pt}#1}\par\nopagebreak\vspace{4pt}
}

% ============================================================
% Boîtes pédagogiques — même recette : fond blanc, bordure épaisse colorée,
% ombre dure encre, badge plein. Titre optionnel [#1].
% ============================================================
\tcbset{popbase/.style={breakable,enhanced,colback=white,boxrule=2.3pt,arc=9pt,
  shadow={3pt}{-3pt}{0pt}{popink},left=14pt,right=14pt,top=11pt,bottom=11pt,before skip=11pt,after skip=15pt,
  fontupper=\popsemi\fontsize{10.6}{14.5}\selectfont\color{popink},
  fontlower=\popsemi\fontsize{10.6}{14.5}\selectfont\color{popink}}}
% Coche dessinée (le glyphe ✓ n'existe pas dans les polices du document)
\newcommand{\popcheck}[1]{\tikz[baseline=-0.5ex]\draw[#1,line width=1.6pt,line cap=round,line join=round](-0.13,0.0)--(-0.04,-0.09)--(0.13,0.1);}
\providecommand{\coche}{\popcheck{popgreen}}
\providecommand{\resultat}[1]{\par\smallskip\noindent\fcolorbox{popgreen}{popgreenL}{\parbox{\dimexpr\linewidth-2\fboxsep-2\fboxrule\relax}{\textbf{Résultat.} #1}}\par\smallskip}
\newcommand{\popboxhead}[3]{\poppill{#1}{#2}{white}\ifx\relax#3\relax\else\hspace{6pt}{\popxbold\fontsize{10.6}{12}\selectfont\color{popink}#3}\fi\par\vspace{7pt}}

\newtcolorbox{aretenir}[1][]{popbase,colframe=popgreen,before upper={\popboxhead{À retenir}{popgreen}{#1}}}
\newtcolorbox{exemplebox}[1][]{popbase,colframe=poporange,before upper={\popboxhead{Exemple}{poporange}{#1}}}
\newtcolorbox{definition}[1][]{popbase,colframe=popblue,before upper={\popboxhead{Définition}{popblue}{#1}}}
\newtcolorbox{propriete}[1][]{popbase,colframe=poppurple,before upper={\popboxhead{Propriété}{poppurple}{#1}}}
\newtcolorbox{methode}[1][]{popbase,colframe=popgold,before upper={\popboxhead{Méthode}{popgold}{#1}}}
\newtcolorbox{attention}[1][]{popbase,colframe=poppink,before upper={\popboxhead{Attention}{poppink}{#1}}}
\newtcolorbox{experience}[1][]{popbase,colframe=popblue,colback=popblueL,before upper={\popboxhead{Expérience}{popblue}{#1}}}
% « Le savais-tu ? » : carte violette pleine (contexte, culture, Cameroun)
\newtcolorbox{savaistu}[1][]{popbase,colback=poppurple,colframe=popink,
  fontupper=\popsemi\fontsize{10.4}{14.2}\selectfont\color{white},
  before upper={\poppill{Le savais-tu\,?}{white}{poppurple}\ifx\relax#1\relax\else\hspace{6pt}{\popxbold\color{white}#1}\fi\par\vspace{7pt}}}
% Exercice résolu : énoncé en haut, \tcblower puis la solution sur fond crème
\newcounter{popexo}\newcounter{popexr}
\newtcolorbox{exoresolu}[1][]{popbase,colframe=poporange,colbacklower=popcream,
  segmentation style={popink,line width=1.2pt,dashed},
  before upper={\stepcounter{popexr}\popboxhead{Exercice résolu \thepopexr}{poporange}{#1}},
  before lower={\poppill{Solution}{popink}{popcream}\par\vspace{6pt}}}
% Exercice (non résolu en place) — la correction est dans la partie Corrigés
\newtcolorbox{exercice}[1][]{popbase,colframe=popink,
  before upper={\stepcounter{popexo}\popboxhead{Exercice \thepopexo}{popink}{#1}}}
% Corrigé
\newtcolorbox{corrige}[1][]{popbase,colframe=popgreen,colback=popgreenL,
  before upper={\popboxhead{Corrigé}{popgreen}{#1}}}
% Objectifs / prérequis / bilan
\newtcolorbox{objectifs}{popbase,colframe=popink,colback=poporangeL,
  before upper={\popboxhead{Objectifs de la leçon}{popink}{}}}
\newtcolorbox{prerequis}{popbase,colframe=popink,colback=popblueL,
  before upper={\popboxhead{Prérequis}{popblue}{}}}
\newtcolorbox{bilan}{popbase,colframe=popink,colback=white,boxrule=2.8pt,
  before upper={\popboxhead{Fiche bilan}{poporange}{L'essentiel à connaître pour l'examen}}}
% Figure encadrée avec légende
\newtcolorbox{popfigure}[1][]{enhanced,breakable=false,colback=white,colframe=popline,boxrule=1.4pt,arc=8pt,
  left=10pt,right=10pt,top=10pt,bottom=8pt,before skip=10pt,after skip=12pt,halign=center,
  lower separated=false,halign lower=center,
  fontlower=\popsemi\fontsize{9}{11.5}\selectfont\color{popmuted},
  #1}
\newcommand{\legende}[1]{\par\vspace{6pt}{\popsemi\fontsize{9}{11.5}\selectfont\color{popmuted}#1}}

% ============================================================
% Signature OnBuch+
% ============================================================
\newcommand{\poplogoinline}[1]{{\poplogo\fontsize{#1}{#1}\selectfont\color{popink}OnBuch\color{poporange}+}}
\newcommand{\signatureonbuch}{%
  \begin{tcolorbox}[enhanced,colback=popink,colframe=popink,boxrule=2.2pt,arc=10pt,
      drop shadow=poporange,left=14pt,right=14pt,top=10pt,bottom=10pt,before skip=0pt,after skip=0pt]
    \tikz[baseline=-0.7ex]\node[circle,fill=popgreen,draw=popcream,line width=1.3pt,minimum size=22pt,inner sep=0]{\popcheck{white}};%
    \hspace{9pt}\begin{minipage}[c]{0.62\linewidth}
      {\popxbold\fontsize{12}{13}\selectfont\color{popcream}Signé\ \textbullet\ {\poplogo OnBuch\color{poporange}+}}\par\vspace{2pt}
      {\raggedright\popsemi\fontsize{8}{10.5}\selectfont\color{popline}\MakeUppercase{Équipe pédagogique OnBuch+}\\\MakeUppercase{Conforme au programme officiel MINESEC}\par}
    \end{minipage}\hfill
    {\poplogo\fontsize{22}{22}\selectfont\color{popcream}OnBuch\color{poporange}+}
  \end{tcolorbox}%
}

% ============================================================
% Page de garde
% #1 titre · #2 matière · #3 niveau · #4 module · #5 n° leçon · #6 durée ·
% #7 description / objectifs (texte ou itemize)
% ============================================================
\newcommand{\pagegarde}[7]{%
  \thispagestyle{empty}
  \newgeometry{a4paper,margin=1.7cm,top=1.6cm,bottom=1.4cm}%
  \begin{tikzpicture}[remember picture,overlay]
    \fill[poporange!18] ([xshift=-3.2cm,yshift=-3.6cm]current page.north east) circle (4.6cm);
    \fill[poppurple!14] ([xshift=2.4cm,yshift=7.6cm]current page.south west) circle (3.8cm);
  \end{tikzpicture}%
  {\poplogo\fontsize{30}{30}\selectfont\color{popink}OnBuch\color{poporange}+}\hfill
  \raisebox{0.6ex}{\poppill{Cours signé \textbullet\ Premium}{popink}{popcream}}\par
  \vspace{1pt}\popsquiggle{poppurple}\par
  \vspace{8pt}
  \begin{tcolorbox}[enhanced,colback=poporange,colframe=popink,boxrule=2.8pt,arc=13pt,
      drop shadow=popink,left=18pt,right=18pt,top=12pt,bottom=13pt,after skip=10pt]
    \poppill{#2 \textbullet\ #3}{popink}{popcream}\hspace{5pt}\poppill{#4}{white}{popink}\par\vspace{8pt}
    {\popdisplay\fontsize{14}{14}\selectfont\color{popink}LEÇON #5}\par\vspace{4pt}
    {\raggedright\hyphenpenalty=10000\exhyphenpenalty=10000\popdisplay\fontsize{25}{27}\selectfont\color{popcream}\MakeUppercase{#1}\par}
    \vspace{8pt}
    {\popxbold\fontsize{9.5}{9.5}\selectfont\color{popink}\MakeUppercase{Durée conseillée : #6}}
  \end{tcolorbox}
  \begin{tcolorbox}[enhanced,colback=white,colframe=popink,boxrule=2.2pt,arc=10pt,
      drop shadow=popink,left=16pt,right=16pt,top=10pt,bottom=10pt,after skip=8pt]
    \poppill{Dans cette leçon}{poppurple}{white}\par\vspace{5pt}
    \popsemi\fontsize{9.3}{12.6}\selectfont\color{popink}#7
  \end{tcolorbox}
  \noindent\begin{minipage}[t]{0.487\linewidth}
    \begin{tcolorbox}[enhanced,colback=popgreen,colframe=popink,boxrule=2.2pt,arc=10pt,
        drop shadow=popink,left=11pt,right=11pt,top=10pt,bottom=10pt,height=3.1cm,valign=center]
      \tcbox[colback=white,colframe=popink,boxrule=1.3pt,arc=5pt,boxsep=0pt,nobeforeafter,
        left=3pt,right=3pt,top=3pt,bottom=3pt,valign=center]{\qrcode[height=1.9cm]{https://play.google.com/store/apps/details?id=com.luvvix.onbuch&hl=fr}}%
      \hspace{8pt}\begin{minipage}[c]{0.5\linewidth}
        {\popdisplay\fontsize{11}{12.5}\selectfont\color{white}\MakeUppercase{Télécharge OnBuch}}\par\vspace{4pt}
        {\raggedright\popsemi\fontsize{8.4}{11}\selectfont\color{white}Gratuit \textbullet\ Google Play\\Tuteur IA, annales, concours, résultats\par}
      \end{minipage}
    \end{tcolorbox}
  \end{minipage}\hfill
  \begin{minipage}[t]{0.487\linewidth}
    \begin{tcolorbox}[enhanced,colback=poppurple,colframe=popink,boxrule=2.2pt,arc=10pt,
        drop shadow=popink,left=11pt,right=11pt,top=10pt,bottom=10pt,height=3.1cm,valign=center]
      \tikz[baseline=-0.7ex]\node[circle,fill=white,draw=popink,line width=1.3pt,minimum size=1.6cm,inner sep=0]{\popdisplay\fontsize{24}{24}\selectfont\color{poppurple}+};%
      \hspace{8pt}\begin{minipage}[c]{0.6\linewidth}
        {\popdisplay\fontsize{11}{12.5}\selectfont\color{white}\MakeUppercase{Passe à OnBuch+}}\par\vspace{4pt}
        {\raggedright\popsemi\fontsize{8.4}{11}\selectfont\color{white}Tous les cours signés, Tuteur IA illimité, fiches PDF — onglet OnBuch+ de l'app\par}
      \end{minipage}
    \end{tcolorbox}
  \end{minipage}
  \vspace{8pt}
  \popcontenu
  \vfill
  \signatureonbuch
  \restoregeometry
  \newpage
}

% Rangée « Ce cours contient » (page de garde)
\newcommand{\poptile}[3]{% #1 couleur #2 gros texte #3 légende
  \begin{tcolorbox}[enhanced,colback=#1,colframe=popink,boxrule=2pt,arc=9pt,shadow={2.5pt}{-2.5pt}{0pt}{popink},
      left=6pt,right=6pt,top=9pt,bottom=9pt,halign=center,valign=center,height=1.75cm,width=\linewidth,nobeforeafter]
    {\popdisplay\fontsize{12.5}{13}\selectfont\color{white}\MakeUppercase{#2}}\par\vspace{3pt}
    {\centering\popsemi\fontsize{7.8}{9.5}\selectfont\color{white}#3\par}
  \end{tcolorbox}}
\newcommand{\popcontenu}{%
  \par\noindent{\popxboldls\fontsize{7.5}{7.5}\selectfont\color{popmuted}CE COURS SIGNÉ CONTIENT}\par\vspace{7pt}
  \noindent\begin{minipage}[t]{0.235\linewidth}\poptile{popblue}{Cours}{complet, illustré, pas à pas}\end{minipage}\hfill
  \begin{minipage}[t]{0.235\linewidth}\poptile{poporange}{Méthodes}{et exercices résolus}\end{minipage}\hfill
  \begin{minipage}[t]{0.235\linewidth}\poptile{poppink}{Exercices}{type examen + corrigés}\end{minipage}\hfill
  \begin{minipage}[t]{0.235\linewidth}\poptile{popgreen}{Fiche bilan}{l'essentiel à retenir}\end{minipage}\par}

% Sommaire
\newtcolorbox{sommaire}{enhanced,colback=white,colframe=popink,boxrule=2.2pt,arc=10pt,
  drop shadow=popink,left=18pt,right=18pt,top=13pt,bottom=13pt,before skip=6pt,after skip=20pt,
  before upper={\poppill{Sommaire}{poppurple}{white}\par\vspace{9pt}\popsemi\fontsize{10.6}{14.5}\selectfont\color{popink}}}

% Signature de fin de cours — poussée tout en bas de la dernière page
\newcommand{\findecours}{%
  \par\vfill\nopagebreak
  \begin{center}\popsquiggle{poporange}\end{center}\vspace{4pt}
  \signatureonbuch
}
\AtBeginDocument{\def\llbracket{\mathopen{[\mkern-3.5mu[}}\def\rrbracket{\mathclose{]\mkern-3.5mu]}}} % intervalles d'entiers (glyphe ⟦ absent de la police)
% Glyphes absents de Fira Math : redéfinis à partir de glyphes disponibles
\AtBeginDocument{%
  \def\setminus{\mathbin{\backslash}}\let\smallsetminus\setminus
  \def\vdots{\mathord{\vbox{\baselineskip3.2pt\lineskiplimit0pt\kern4pt\hbox{.}\hbox{.}\hbox{.}}}}%
  \def\ddots{\mathinner{\mkern1mu\raise6.5pt\hbox{.}\mkern2mu\raise3.6pt\hbox{.}\mkern2mu\raise0.7pt\hbox{.}\mkern1mu}}%
  \def\triangle{\mathord{\scalebox{0.9}{$\symup{\Delta}$}}}%
  \def\nabla{\mathord{\raisebox{\depth}{\scalebox{1}[-1]{$\symup{\Delta}$}}}}%
  \def\checkmark{\popcheck{popdark}}%
  \def\star{\mathbin{\ast}}\def\bigstar{\mathord{\ast}}%
  \def\diamond{\mathbin{\scalebox{0.6}{$\Diamond$}}}%
  \def\bigcup{\mathop{\mathchoice{\raisebox{-0.3em}{\scalebox{1.7}{$\cup$}}}{\scalebox{1.25}{$\cup$}}{\cup}{\cup}}\displaylimits}%
  \def\bigcap{\mathop{\mathchoice{\raisebox{-0.3em}{\scalebox{1.7}{$\cap$}}}{\scalebox{1.25}{$\cap$}}{\cap}{\cap}}\displaylimits}%
  \def\lesseqqgtr{\mathrel{\lessgtr}}\def\bigoplus{\mathop{\oplus}\displaylimits}%
}
\providecommand{\ssi}{si et seulement si} % abréviation fréquente
\AtBeginDocument{\providecommand{\wideparen}{\overparen}} % arc de cercle

%% FIGFIX-BEGIN v5 — correctifs globaux des figures (docs/onbuch-generation/tools/figfix)
\usepackage{adjustbox}\usepackage{etoolbox}\tcbuselibrary{raster}
% toute figure TikZ/pgfplots plus large que la ligne est réduite à la largeur disponible
\BeforeBeginEnvironment{tikzpicture}{\begin{adjustbox}{max width=\linewidth}}
\AfterEndEnvironment{tikzpicture}{\end{adjustbox}}
% légendes pgfplots lisibles par-dessus les courbes
\pgfplotsset{every axis/.append style={legend style={fill=white,fill opacity=0.92,text opacity=1,draw=black!15,inner sep=3pt}}}
% étiquettes d'arêtes (syntaxe edge["..."]) avec fond blanc
\tikzset{every edge quotes/.append style={fill=white,inner sep=1.2pt}}
%% FIGFIX-END
\begin{document}

\pagegarde{Perturbations du système immunitaire : Le VIH/Sida}{SVTEEHB}{3ème}{SVTEEHB – classe de 3ème}{N07}{3 séances (fiche de progression harmonisée)}{Cette leçon explique le cycle de réplication du VIH au cœur des lymphocytes T CD4+, analyse les phases cliniques de l'infection à travers l'évolution de la charge virale et du taux de CD4, et évalue les stratégies de prévention et l'action ciblée des antirétroviraux (ARV) pour comprendre la prise en charge thérapeutique actuelle au Cameroun.
\vspace{4pt}
\begin{multicols}{2}\small
\begin{itemize}
\item Décrire la structure du VIH et identifier son matériel génétique et ses enzymes clés.
\item Expliquer les étapes du cycle de réplication du VIH (fixation, transcription inverse, intégration, transcription, traduction, assemblage, bourgeonnement) en reliant chaque étape à l'action d'une classe d'ARV.
\item Interpréter les courbes parallèles de la charge virale et du taux de lymphocytes T CD4+ au cours des trois phases de l'infection (primaire, chronique, SIDA).
\item Justifier la latence clinique par le piégeage viral dans les organes lymphoïdes et l'épuisement progressif du système immunitaire.
\item Différencier infection VIH et maladie SIDA à l'aide de critères biologiques (seuil CD4) et cliniques (infections opportunistes).
\item Analyser les résultats d'un test de dépistage (recherche Ag p24 / Ac anti-VIH) pour conclure sur le statut sérologique.
\end{itemize}\end{multicols}}

\begin{sommaire}
\begin{multicols}{2}
\begin{enumerate}
\item 1. Le VIH : identification, structure et cellules cibles
\item 2. Cycle de réplication : de l'ARN viral au nouveau virion infectieux
\item 3. Dynamique de l'infection : Phases cliniques, marqueurs biologiques et physiopathologie
\item 4. Prévention, Dépistage et Traitement Antirétroviral (ARV) au Cameroun
\item 5. Bilan synthèse \& Exercices d'intégration type BEPC
\item Activité d'intégration
\item Méthodes et exercices résolus
\item Exercices
\item Corrigés des exercices
\item Fiche bilan
\end{enumerate}
\end{multicols}
\end{sommaire}

\begin{objectifs}
\begin{itemize}
\item Décrire la structure du VIH et identifier son matériel génétique et ses enzymes clés.
\item Expliquer les étapes du cycle de réplication du VIH (fixation, transcription inverse, intégration, transcription, traduction, assemblage, bourgeonnement) en reliant chaque étape à l'action d'une classe d'ARV.
\item Interpréter les courbes parallèles de la charge virale et du taux de lymphocytes T CD4+ au cours des trois phases de l'infection (primaire, chronique, SIDA).
\item Justifier la latence clinique par le piégeage viral dans les organes lymphoïdes et l'épuisement progressif du système immunitaire.
\item Différencier infection VIH et maladie SIDA à l'aide de critères biologiques (seuil CD4) et cliniques (infections opportunistes).
\item Analyser les résultats d'un test de dépistage (recherche Ag p24 / Ac anti-VIH) pour conclure sur le statut sérologique.
\item Argumenter l'intérêt de la trithérapie (association de 3 ARV) pour prévenir la résistance et atteindre l'indétectabilité (I=I).
\item Appliquer les connaissances aux stratégies de prévention combinée (PrEP, PEP, TASP, PMTCT) dans le contexte camerounais.
\end{itemize}
\end{objectifs}

\begin{prerequis}
\begin{itemize}
\item Structure d'une cellule eucaryote (noyau, cytoplasme, ribosomes, membrane).
\item Rôle des lymphocytes T CD4+ (coordination immunité humorale et cellulaire) et T CD8+ (lyse cellulaire) vus en 4ème.
\item Notion de virus : acapside, matériel génétique (ADN ou ARN), parasitisme intracellulaire obligatoire.
\item Dogme central de la biologie moléculaire : ADN -> ARN -> Protéine (Transcription / Traduction).
\item Principes de base de la sérologie (réaction Antigène-Anticorps).
\end{itemize}
\end{prerequis}

\newpage

\coursec{Le VIH : identification, structure et cellules cibles}

\courssub{Situation d'accroche et problématique}

Au Cameroun, malgré les progrès remarquables du Plan National de Lutte contre le Sida (PNLS), près de \num{500000} personnes vivaient avec le VIH en 2023, dont une proportion significative de jeunes âgés de 15 à 24 ans. Imagine Amina, 15 ans, élève en classe de \textbf{3ème} dans un collège de Yaoundé. Lors d'une campagne de dépistage volontaire organisée au Centre de Santé de son arrondissement, elle découvre sa séropositivité. Elle ne se sent pas malade, n'a pas de fièvre, pas de toux, pas de diarrhée. Pourtant, le médecin lui annonce qu'elle doit commencer un traitement antirétroviral (ARV) à vie, à prendre chaque jour à heure fixe, sans faille.

Comment un virus microscopique, invisible à l'œil nu, peut-il détruire les défenses d'un organisme entier sans provoquer de symptômes immédiats ? Pourquoi la prise du traitement doit-elle être si rigoureuse, sous peine de voir le virus devenir résistant ? Pour répondre à ces questions vitales, nous devons d'abord comprendre l'identité de l'ennemi : sa structure, sa stratégie d'invasion et sa cible privilégiée.

\courssub{Historique et classification : deux virus, une pandémie}

L'histoire de la découverte du VIH est une page majeure de la virologie moderne, marquée par une course scientifique intense au début des années 1980.

\begin{itemize}
    \item \textbf{1983 :} À l'Institut Pasteur de Paris, l'équipe de \textbf{Luc Montagnier} (avec \textbf{Françoise Barré-Sinoussi} et Jean-Claude Chermann) isole un nouveau rétrovirus chez un patient atteint de lymphadénopathie, qu'ils nomment \textbf{LAV} (Lymphadenopathy-Associated Virus).
    \item \textbf{1984 :} Aux États-Unis, \textbf{Robert Gallo} isole un virus qu'il nomme \textbf{HTLV-III} (Human T-Lymphotropic Virus type III) et démontre formellement son rôle étiologique dans le SIDA.
    \item \textbf{1986 :} La communauté scientifique s'accorde sur la dénomination commune : \textbf{VIH} (Virus de l'Immunodéficience Humaine). On distingue deux types :
    \begin{itemize}
        \item \textbf{VIH-1 :} Responsable de la pandémie mondiale (groupe M principal). Origine : transmission zoonotique du chimpanzé (\emph{Pan troglodytes troglodytes}) à l'homme en Afrique centrale (sud-Cameroun) au début du XXe siècle.
        \item \textbf{VIH-2 :} Endémique en Afrique de l'Ouest, moins virulent, transmission moins efficace, progression plus lente vers le SIDA. Origine : singe \emph{Cercocebus atys} (mangabey fuligineux).
    \end{itemize}
    \item \textbf{2008 :} Prix Nobel de Physiologie ou Médecine décerné à \textbf{Françoise Barré-Sinoussi} et \textbf{Luc Montagnier} pour la découverte du VIH.
\end{itemize}

\begin{definition}[Rétrovirus]
Un \textbf{rétrovirus} est un virus à ARN dont le génome est \textbf{rétrotranscrit en ADN} par une enzyme virale spécifique, la \textbf{transcriptase inverse (TI)}, avant d'être intégré dans le génome de la cellule hôte. Cette intégration rend l'infection permanente pour la durée de vie de la cellule.
\end{definition}

\begin{savaistu}[Le VIH au Cameroun : une diversité exceptionnelle]
Le Cameroun est l'un des rares pays au monde où circulent \textbf{les deux types de VIH (VIH-1 et VIH-2)} et où l'on retrouve une grande diversité de groupes (M, N, O, P pour le VIH-1) et de sous-types (CRF02\_AG majoritaire). Cette diversité complique le diagnostic (tests de dépistage) et le suivi de la charge virale, mais elle témoigne de l'origine zoonotique ancienne du virus dans les forêts du bassin du Congo.
\end{savaistu}

\courssub{Morphologie et structure du virion VIH-1}

Le virion (particule virale complète, infectieuse) du VIH-1 est une structure sphérique d'environ \SI{100}{\nano\meter}--\SI{120}{\nano\meter} de diamètre. Bien que minuscule (environ 60 fois plus petit qu'un globule rouge), son architecture est d'une précision moléculaire remarquable, fruit de millions d'années d'évolution.

\begin{popfigure}
\begin{tikzpicture}[scale=1.1, every node/.style={font=\small}]
    % Enveloppe externe
    \draw[popblue, very thick, fill=popblue!10] (0,0) circle (3.5cm);
    % Glycoprotéines de surface (gp120/gp41)
    \foreach \angle in {0,30,60,90,120,150,180,210,240,270,300,330} {
        \draw[poporange, thick] (\angle:3.5) -- (\angle:4.2);
        \fill[poporange] (\angle:4.2) circle (0.08cm);
        \draw[poppink, thick] (\angle:3.5) -- (\angle:4.0);
        \fill[poppink] (\angle:4.0) circle (0.06cm);
    }
    % Matrice (p17) - ligne fine sous l'enveloppe
    \draw[popdark, dashed] (0,0) circle (3.3cm);
    % Capside conique (p24) - représenté par deux triangles accolés (forme conique en coupe)
    \fill[popgreen!60, opacity=0.8] (0,0.5) -- (-1.2,-1.5) -- (1.2,-1.5) -- cycle;
    \fill[popgreen!60, opacity=0.8] (0,-0.5) -- (-1.2,1.5) -- (1.2,1.5) -- cycle;
    % ARN génomique (2 brins)
    \draw[poppurple, very thick, decorate, decoration={snake, amplitude=0.5pt, segment length=4pt}] (-0.8,0) -- (0.8,0);
    \draw[poppurple, very thick, decorate, decoration={snake, amplitude=0.5pt, segment length=4pt}] (-0.8,-0.15) -- (0.8,-0.15);
    % Enzymes
    \node[poppurple, circle, fill=poppurple!20, draw, minimum size=0.5cm] (TI) at (-1.5, 0.5) {TI};
    \node[poppurple, circle, fill=poppurple!20, draw, minimum size=0.5cm] (INT) at (1.5, 0.5) {Int};
    \node[poppurple, circle, fill=poppurple!20, draw, minimum size=0.5cm] (PROT) at (0, -1.8) {Prot};
    
    % Légendes avec flèches
    \draw[->, popblue] (3.8, 2.5) -- (3.5, 1.8) node[right, align=left, text width=3cm] {\textbf{Enveloppe lipidique}\\(d'origine cellulaire)};
    \draw[->, poporange] (4.5, 3.0) -- (4.2, 2.5) node[right, align=left, text width=2.5cm] {\textbf{gp120 (SU)}\\(Fixation au CD4)};
    \draw[->, poppink] (4.5, 1.5) -- (4.0, 1.2) node[right, align=left, text width=2.5cm] {\textbf{gp41 (TM)}\\(Fusion membranaire)};
    \draw[->, popdark] (3.8, -2.5) -- (3.3, -1.8) node[right, align=left, text width=3cm] {\textbf{Matrice p17}\\(Protéine structurale)};
    \draw[->, popgreen] (1.8, -1.5) -- (1.5, -1.0) node[right, align=left, text width=3cm] {\textbf{Capside p24}\\(Protéine majeure,\\forme conique)};
    \draw[->, poppurple] (-2.5, 0.5) -- (-1.7, 0.5) node[left, align=right, text width=3cm] {\textbf{Transcriptase Inverse (TI)}\\(2 copies, RT)};
    \draw[->, poppurple] (2.5, 0.5) -- (1.8, 0.5) node[right, align=left, text width=3cm] {\textbf{Intégrase (IN)}\\(Intégration ADN)};
    \draw[->, poppurple] (0, -2.5) -- (0, -2.0) node[below, align=center, text width=3cm] {\textbf{Protéase (PR)}\\(Maturation virale)};
    \draw[->, poppurple] (-1.5, -0.5) -- (-1.0, -0.2) node[left, align=right, text width=3cm] {\textbf{ARN génomique}\\(2 brins +, identiques,\\$\sim$ 9,7 kb)};
\end{tikzpicture}
\legende{Schéma annoté de la structure du virion VIH-1 en coupe. L'enveloppe lipidique porte les glycoprotéines gp120 (fixation) et gp41 (fusion). La capside conique (p24) protège l'ARN génomique (2 brins) et les 3 enzymes essentielles : Transcriptase Inverse (TI), Intégrase (IN), Protéase (PR).}
\end{popfigure}

\begin{aretenir}[Composition détaillée du virion VIH-1]
\begin{itemize}
    \item \textbf{Enveloppe :} Bicouche lipidique dérivée de la membrane plasmique de la cellule hôte lors du bourgeonnement (budding). Elle incorpore des \textbf{glycoprotéines virales} organisées en trimères (3 gp120 + 3 gp41) : les \textbf{peplomères}.
    \item \textbf{Glycoprotéines de surface :}
    \begin{itemize}
        \item \textbf{gp120 (SU - Surface Unit) :} Extrémité externe. Site de fixation au récepteur \textbf{CD4} et au co-récepteur (CCR5 ou CXCR4). Très variable (boucles variables V1-V5) $\rightarrow$ échappement immunitaire.
        \item \textbf{gp41 (TM - TransMembrane) :} Ancrée dans l'enveloppe. Domaine de fusion hydrophobe (peptide de fusion) masqué dans l'état natif. Médiatise la fusion des membranes virale et cellulaire.
    \end{itemize}
    \item \textbf{Matrice (p17) :} Protéine structurale tapissant la face interne de l'enveloppe, assurant la stabilité et le transport intracellulaire.
    \item \textbf{Capside (p24 / CA) :} Structure protéique \textbf{conique} (caractéristique des lentivirus) formée de \num{1500} à \num{2500} copies de la protéine p24. Elle protège le noyau ribonucléoprotéique.
    \item \textbf{Noyau ribonucléoprotéique :}
    \begin{itemize}
        \item \textbf{2 brins d'ARN simple brin à polarité positive (+)} (chacun $\sim$ \num{9700} bases), identiques, non covalents (diploïde). L'ARN porte les gènes \emph{gag, pol, env, tat, rev, nef, vif, vpr, vpu}.
        \item \textbf{Protéine nucléocapside (p7/NC) :} Lie l'ARN (doigts de zinc), essentielle pour le packaging et la transcription inverse.
        \item \textbf{3 Enzymes virales (codées par \emph{pol}) :}
        \begin{enumerate}
            \item \textbf{Transcriptase Inverse (TI / RT) :} Dimère p66/p51. Activité \textbf{ADN polymérase ARN-dépendante} (synthèse ADN sur matrice ARN) + activité \textbf{RNase H} (dégradation de l'ARN de l'hybride ARN/ADN). \textbf{Pas d'activité de relecture (proofreading)} $\rightarrow$ taux d'erreur élevé.
            \item \textbf{Intégrase (IN) :} Catalyse l'intégration de l'ADN viral (provirus) dans l'ADN chromosomique hôte.
            \item \textbf{Protéase (PR) :} Aspartyl-protéase. Clive les polyprotéines précurseurs \emph{Gag} et \emph{Gag-Pol} lors de la maturation du virion (immature $\rightarrow$ mature infectieux).
        \end{enumerate}
    \end{itemize}
\end{itemize}
\end{aretenir}

\begin{methode}[Identifier les composants du VIH sur une micrographie électronique à transmission (MET)]
\begin{enumerate}
    \item \textbf{Observer la forme globale :} Particule sphérique, diamètre $\approx$ \SI{100}{\nano\meter}--\SI{120}{\nano\meter}.
    \item \textbf{Repérer l'enveloppe :} Double ligne foncée (bicouche lipidique) à la périphérie. Parfois des "piques" (peplomères gp120/gp41) visibles en projection.
    \item \textbf{Identifier la capside :} Structure dense, souvent conique ou arrondie, au centre du virion (protéine p24, très antigénique).
    \item \textbf{Localiser le matériel génétique :} Fils electron-denses (ARN) à l'intérieur de la capside, souvent associés à des points denses (enzymes TI, Intégrase).
    \item \textbf{Conclure :} La présence d'une capside conique centrale + enveloppe à peplomères + ARN central est typique des \textbf{Lentivirus} (genre du VIH).
\end{enumerate}
\end{methode}

\begin{exemplebox}[Micrographie électronique du VIH - Analyse guidée]
\begin{center}
\begin{tikzpicture}[scale=1.5]
    % Simulation d'une MET
    \draw[popdark, very thick, fill=popcream] (0,0) circle (2cm);
    % Enveloppe
    \draw[popblue, thick] (0,0) circle (1.9cm);
    \draw[popblue, thick] (0,0) circle (1.85cm);
    % Piques
    \foreach \a in {10,50,90,130,170,210,250,290,330} {
        \draw[poporange, thick] (\a:1.9) -- (\a:2.15);
    }
    % Capside conique (simplifiée)
    \fill[popgreen!80] (0,0.3) -- (-0.7,-0.9) -- (0.7,-0.9) -- cycle;
    \fill[popgreen!80] (0,-0.3) -- (-0.7,0.9) -- (0.7,0.9) -- cycle;
    % ARN/Enzymes (points)
    \foreach \x/\y in {-0.3/0.1, 0.2/0.0, -0.1/-0.2, 0.4/-0.1, -0.4/-0.3} {
        \fill[poppurple] (\x,\y) circle (0.04cm);
    }
    % Flèches de légende
    \draw[->, popblue] (2.2, 1.5) -- (1.9, 1.0) node[right] {Enveloppe (bicouche)};
    \draw[->, poporange] (2.3, 0.5) -- (1.95, 0.3) node[right] {Peplomères (gp120/gp41)};
    \draw[->, popgreen] (1.2, -1.0) -- (0.9, -0.7) node[right] {Capside p24 (conique)};
    \draw[->, poppurple] (-1.5, -1.2) -- (-0.8, -0.6) node[left] {Nucléoprotéines (ARN + TI/IN)};
\end{tikzpicture}
\end{center}
\legende{Micrographie électronique à transmission (simulée) d'un virion VIH-1. On distingue clairement la bicouche lipidique de l'enveloppe, les projections de surface (peplomères), la capside centrale dense (p24) de forme conique et le matériel ribonucléoprotéique interne. Barre d'échelle : \SI{100}{\nano\meter}.}
\end{exemplebox}

\courssub{Cellules cibles : les porteurs du CD4}

Le VIH ne peut pas infecter n'importe quelle cellule. Son tropisme (spécificité tissulaire) est strictement dicté par la présence de récepteurs spécifiques à la surface de la cellule hôte.

\begin{definition}[Lymphocyte T CD4+ (Cellule T auxiliaire / Helper)]
Cellule coordinatrice centrale de la réponse immune adaptative. Elle porte le récepteur \textbf{CD4} à sa surface. Elle reconnaît les antigènes présentés par les molécules \textbf{CMHC de classe II} sur les cellules présentatrices d'antigène (macrophages, cellules dendritiques, lymphocytes B). Par sécrétion de \textbf{cytokines} (IL-2, IFN-$\gamma$, IL-4, etc.), elle active :
\begin{itemize}
    \item Les \textbf{lymphocytes B} $\rightarrow$ production d'anticorps (immunité humorale).
    \item Les \textbf{lymphocytes T CD8+} cytotoxiques $\rightarrow$ élimination des cellules infectées (immunité cellulaire).
    \item Les \textbf{macrophages} $\rightarrow$ phagocytose et destruction intracellulaire accrue.
\end{itemize}
Sa destruction massive entraîne une \textbf{immunodépression profonde} touchant toutes les branches de l'immunité.
\end{definition}

\begin{propriete}[Récepteurs d'entrée du VIH : Le duo CD4 / Co-récepteur]
L'entrée du VIH dans la cellule cible nécessite \textbf{deux interactions séquentielles} (modèle "clé-serrure" à deux serrures) :
\begin{enumerate}
    \item \textbf{Récepteur principal : CD4.} Protéine de la superfamille des immunoglobulines, exprimée sur \textbf{lymphocytes T CD4+}, \textbf{macrophages}, \textbf{cellules dendritiques} (et quelques cellules cérébrales : microglie, astrocytes). La gp120 se lie au domaine D1 du CD4.
    \item \textbf{Co-récepteur : Chémokine Récepteur (7 domaines transmembranaires, couplé aux protéines G).} Deux principaux :
    \begin{itemize}
        \item \textbf{CCR5 (R5) :} Exprimé sur macrophages, cellules dendritiques, lymphocytes T mémoire (phénotype effecteur). Utilisé par les souches \textbf{R5-tropiques (M-tropiques)}. Majoritaires lors de la \textbf{transmission} et au stade \textbf{asymptomatique}.
        \item \textbf{CXCR4 (X4) :} Exprimé sur lymphocytes T naïfs, cellules souches hématopoïétiques. Utilisé par les souches \textbf{X4-tropiques (T-tropiques)}. Apparaissent souvent au stade \textbf{symptomatique/SIDA}, corrélées à une chute brutale des CD4.
    \end{itemize}
    \item \textbf{Variants à double tropisme (R5X4) :} Utilisent les deux.
\end{enumerate}
\end{propriete}

\begin{savaistu}[La mutation $\Delta$32 du CCR5 : une résistance naturelle]
Une délétion de 32 paires de bases dans le gène \emph{CCR5} (allèle \textbf{CCR5-$\Delta$32}) conduit à l'absence de récepteur CCR5 fonctionnel à la surface des cellules.
\begin{itemize}
    \item \textbf{Homozygotes $\Delta$32/$\Delta$32 (environ 1\% des Européens, rare en Afrique) :} \textbf{Résistants} à l'infection par les souches R5 (majoritaires). Cas célèbre : "Patient de Berlin" (Timothy Ray Brown), guéri du VIH après greffe de moelle d'un donneur $\Delta$32/$\Delta$32.
    \item \textbf{Hétérozygotes $\Delta$32/WT :} Expression réduite de CCR5 $\rightarrow$ \textbf{progression plus lente} vers le SIDA (long-term non-progressors).
\end{itemize}
La découverte du rôle du CCR5 comme co-récepteur du VIH (1996) a été une avancée majeure pour la compréhension de la pathogénie et le développement de nouveaux antiviraux (ex: Maraviroc).
\end{savaistu}

\courssub{Mécanisme d'entrée : de la fixation à la fusion}

L'entrée du VIH est un processus dynamique, déclenché par des changements conformationnels successifs de l'enveloppe virale (gp120/gp41). C'est une cible thérapeutique majeure (inhibiteurs d'entrée).

\begin{popfigure}
\begin{tikzpicture}[scale=0.95, every node/.style={font=\small}]
    % ETAPE 1 : Fixation gp120 - CD4
    \begin{scope}[xshift=0cm]
        \node[align=center] at (0, 3.2) {\textbf{Étape 1 : Fixation gp120 -- CD4}};
        % Cellule
        \draw[popdark, thick, fill=popblue!5] (-2, -1.5) rectangle (2, -0.5);
        \node at (0, -1.0) {Membrane cellulaire};
        % CD4
        \draw[popblue, very thick] (0, -0.5) -- (0, 0.2);
        \fill[popblue] (0, 0.2) circle (0.15cm) node[above=0.1cm] {CD4};
        % Virus
        \draw[poporange, thick, fill=poporange!10] (-0.8, 1.5) -- (0.8, 1.5) -- (0.6, 2.2) -- (-0.6, 2.2) -- cycle; % Virion simplifié
        \draw[poporange, thick] (-0.4, 1.5) -- (-0.4, 0.8); % gp120
        \fill[poporange] (-0.4, 0.8) circle (0.08cm);
        \draw[poppink, thick] (0.4, 1.5) -- (0.4, 0.8); % gp41
        \fill[poppink] (0.4, 0.8) circle (0.06cm);
        \node at (0, 2.5) {Virion VIH};
        \draw[<->, popgreen, thick] (-0.4, 0.8) -- (0, 0.2);
        \node[popgreen, right] at (0.5, 0.5) {Affinitée forte};
    \end{scope}

    % ETAPE 2 : Changement conformationnel & Fixation Co-récepteur
    \begin{scope}[xshift=5.5cm]
        \node[align=center] at (0, 3.2) {\textbf{Étape 2 : Co-récepteur (CCR5/CXCR4)}};
        % Cellule
        \draw[popdark, thick, fill=popblue!5] (-2, -1.5) rectangle (2, -0.5);
        \node at (0, -1.0) {Membrane cellulaire};
        % CD4 + gp120 lié
        \draw[popblue, very thick] (0, -0.5) -- (0, 0.2);
        \fill[popblue] (0, 0.2) circle (0.15cm);
        \draw[poporange, thick] (-0.2, 0.2) -- (-0.2, 0.8);
        \fill[poporange] (-0.2, 0.8) circle (0.08cm) node[right] {gp120};
        % Co-récepteur
        \draw[poppurple, very thick] (1.0, -0.5) -- (1.0, 0.2);
        \fill[poppurple] (1.0, 0.2) circle (0.15cm) node[above=0.1cm] {CCR5 / CXCR4};
        % gp120 se déplace vers co-récepteur
        \draw[->, poporange, dashed] (-0.2, 0.8) .. controls (0.2, 1.0) and (0.8, 1.0) .. (1.0, 0.5);
        \node[poporange, align=center] at (0.5, 1.3) {Changement\\conformationnel\\gp120};
        % gp41 exposé
        \draw[poppink, thick] (0.4, 1.5) -- (0.4, 0.5);
        \fill[poppink] (0.4, 0.5) circle (0.06cm) node[right] {gp41 (exposé)};
    \end{scope}

    % ETAPE 3 : Fusion médiée par gp41
    \begin{scope}[xshift=11cm]
        \node[align=center] at (0, 3.2) {\textbf{Étape 3 : Fusion gp41 \& Entrée}};
        % Membranes fusionnées
        \draw[popdark, very thick, fill=popblue!5] (-2, -1.5) .. controls (-1, -0.2) and (1, -0.2) .. (2, -1.5) -- (2, -1.0) .. controls (1, -0.3) and (-1, -0.3) .. (-2, -1.0) -- cycle;
        \node at (0, -1.0) {Membranes fusionnées};
        % gp41 hairpin
        \draw[poppink, very thick, decorate, decoration={zigzag, amplitude=2pt, segment length=4pt}] (-0.5, 0.0) -- (0.5, 0.0);
        \node[poppink] at (0, 0.5) {gp41 : "Hairpin" (fusion)};
        % Capside libérée
        \fill[popgreen!80] (0, 1.0) circle (0.4cm);
        \node[popgreen] at (0, 1.0) {p24};
        \draw[poppurple, decorate, decoration={snake}] (0, 1.0) -- (0, 1.6);
        \node[poppurple] at (0, 1.8) {ARN + TI/IN};
        \draw[->, popdark] (0, 1.6) -- (0, 2.0) node[above] {Libération dans le cytoplasme};
    \end{scope}
    
    % Flèches entre étapes
    \draw[->, thick, popdark] (3.0, 3.2) -- (3.5, 3.2);
    \draw[->, thick, popdark] (8.5, 3.2) -- (9.0, 3.2);
\end{tikzpicture}
\legende{Mécanisme d'entrée du VIH dans la cellule cible (lymphocyte T CD4+ ou macrophage). 1. Fixation de la gp120 au récepteur CD4. 2. Changement conformationnel de la gp120 exposant le site de liaison au co-récepteur (CCR5 ou CXCR4). 3. Fixation au co-récepteur déclenchant l'insertion du peptide de fusion de la gp41 dans la membrane cellulaire, formation de la structure en "hairpin" (épingle à cheveux) rapprochant les deux membranes $\rightarrow$ Fusion $\rightarrow$ Libération de la capside (p24) contenant l'ARN viral et les enzymes dans le cytoplasme.}
\end{popfigure}

\begin{experience}[Preuve expérimentale : Rôle du CD4 et de la gp120 (Blocage par anticorps)]
\textbf{Question :} Le CD4 et la gp120 sont-ils indispensables et suffisants pour l'infection ?
\textbf{Expérience (in vitro, lignées cellulaires) :}
\begin{enumerate}
    \item \textbf{Cellules cibles :} Lignées de lymphocytes T (ex: SupT1, CEM) exprimant CD4 (+ CCR5/CXCR4). Contrôle : mêmes lignées rendues CD4-négatives (par CRISPR ou sélection).
    \item \textbf{Virus :} VIH-1 souche de laboratoire (ex: NL4-3, X4-tropique).
    \item \textbf{Traitements :}
    \begin{itemize}
        \item Témoin : Virus seul.
        \item Anti-CD4 (anticorps monoclonal anti-CD4, ex: OKT4A) pré-incubé avec cellules cibles.
        \item Anti-gp120 (anticorps monoclonal anti-gp120, ex: b12, VRC01) pré-incubé avec virus.
        \item Soluble CD4 (sCD4, recombinant) pré-incubé avec virus (leurre).
    \end{itemize}
    \item \textbf{Lecture :} Dosage de l'antigène p24 dans le surnageant au jour J+3 à J+7 (marqueur de réplication virale) ou PCR ADN proviral.
\end{enumerate}
\textbf{Résultats attendus :}
\begin{itemize}
    \item Infection forte dans le témoin (cellules CD4+).
    \item \textbf{Blocage quasi-total} de l'infection avec Anti-CD4, Anti-gp120 ou sCD4 sur cellules CD4+.
    \item \textbf{Aucune infection} dans les cellules CD4-négatives (même sans anticorps).
\end{itemize}
\textbf{Interprétation :} L'interaction \textbf{gp120 -- CD4} est \textbf{nécessaire} (pas d'infection sans CD4) et \textbf{spécifique} (blocage par anticorps dirigés contre l'un ou l'autre partenaire). Le sCD4 agit comme un "leurre" (decoy) : il se fixe sur la gp120, déclenche prématurément le changement conformationnel et l'exposition de la gp41, inactivant le virion avant qu'il n'atteigne la cellule.
\textbf{Conclusion :} La fixation gp120/CD4 est l'étape initiale obligatoire et une cible thérapeutique validée (ex: Ibalizumab, anticorps anti-CD4 approuvé pour VIH multi-résistant).
\end{experience}

\courssub{Devenir de l'ARN viral : la nécessité de la transcription inverse}

Une fois la capside libérée dans le cytoplasme, le VIH fait face à un obstacle fondamental : \textbf{la machinerie cellulaire ne sait pas lire l'ARN viral}.

\begin{itemize}
    \item L'ARN viral est un \textbf{ARN simple brin à polarité positive (+)}. Il ressemble à un ARNm (il a une coiffe 5', une queue poly-A 3'), \textbf{MAIS} :
    \begin{itemize}
        \item Il ne possède \textbf{pas de promoteur} reconnu par l'ARN polymérase II cellulaire (le promoteur viral, le LTR 5', est dans l'ADN, pas dans l'ARN entrant).
        \item Il contient des \textbf{structures secondaires complexes} (TAR, poly-A, PBS, DIS, $\Psi$) qui bloquent les ribosomes.
        \item Il code pour des \textbf{polyprotéines} (Gag, Gag-Pol) nécessitant un cadre de lecture décalé (\emph{frameshifting} ribosomal) que la cellule gère mal sans facteurs viraux.
    \end{itemize}
    \item \textbf{Conséquence :} L'ARN viral \textbf{ne peut pas être traduit directement} en protéines virales par les ribosomes de l'hôte. Il ne peut pas non plus s'intégrer dans le génome (l'intégrase agit sur de l'ADN).
    \item \textbf{Solution évolutive :} Le virus apporte sa propre usine : la \textbf{Transcriptase Inverse (TI)}. Elle va synthétiser une copie \textbf{ADN double brin (ADNdb)} à partir de l'ARN matriciel. C'est la \textbf{Transcription Inverse} (ou Rétrotranscription).
\end{itemize}

\begin{definition}[Transcription Inverse (Rétrotranscription)]
Processus enzymatique catalysé par la \textbf{Transcriptase Inverse (TI)} qui synthétise une molécule d'\textbf{ADN double brin} (ADN proviral) à partir du génome \textbf{ARN simple brin (+)} viral. Découverte indépendante par \textbf{Howard Temin} (virus ARN tumoraux) et \textbf{David Baltimore} (virus ARN murins) en \textbf{1970}. Prix Nobel 1975 (partagé avec Renato Dulbecco).
\end{definition}

\begin{exemplebox}[Taux de mutation de la TI : le moteur de l'échappement]
La Transcriptase Inverse du VIH \textbf{ne possède pas d'activité exonucléasique 3'$\rightarrow$5' (proofreading / relecture)}.
\begin{itemize}
    \item Taux d'erreur intrinsèque : $\approx \mathbf{3 \times 10^{-5}}$ \textbf{erreurs par base par cycle de réplication}.
    \item Taille du génome : $\approx \mathbf{9,7 \times 10^3}$ bases (\num{9700} pb).
    \item Nombre de nouveaux virions produits par jour chez un patient non traité : $\approx \mathbf{10^{10}}$ à $\mathbf{10^{11}}$.
\end{itemize}
\textbf{Calcul :} Nombre moyen de mutations par génome et par cycle $\approx 3 \times 10^{-5} \times 9700 \approx \mathbf{0,3}$ mutation/génome/cycle.
Comme $> 1$ virion sur 3 contient au moins une mutation, et que la population virale est immense ($10^{10}$/jour), \textbf{toutes les mutations ponctuelles possibles (y compris celles conférant une résistance aux ARV) sont générées chaque jour}.
\textbf{Conséquence clinique majeure :} La \textbf{monothérapie} (un seul ARV) sélectionne instantanément les mutants résistants préexistants $\rightarrow$ \textbf{Échec thérapeutique rapide}. D'où la nécessité absolue de la \textbf{trithérapie (3 ARV de classes différentes)} : la probabilité qu'un virion ait \emph{simultanément} les 3 mutations de résistance est le produit des probabilités (ex: $10^{-4} \times 10^{-4} \times 10^{-4} = 10^{-12}$), rendant l'émergence de résistance quasi impossible si l'adhérence est parfaite.
\end{exemplebox}

\begin{attention}[Pièges fréquents à éviter absolutely]
\begin{itemize}
    \item \textbf{Confondre VIH et SIDA :} \textbf{VIH} = Virus (l'agent pathogène). \textbf{SIDA} = Syndrome d'Immunodéficience Acquise (le \textbf{stade terminal} de la maladie, défini par l'apparition d'infections opportunistes ou de cancers, ou CD4 $< 200$/mm$^3$). On est \textbf{séropositif} (porteur du virus, anticorps anti-VIH positifs) \textbf{bien avant} d'être \textbf{sidéen} (malade). La période de latence clinique dure en moyenne 8 à 10 ans sans traitement.
    \item \textbf{Dire que le VIH "tue" directement les lymphocytes :} Le virus \textbf{infecte} et \textbf{utilise} la cellule pour se répliquer. La mort cellulaire est indirecte (lysis virale, mort par apoptose "bystander", élimination par lymphocytes T CD8+, activation immune chronique). La destruction massive est une conséquence de la réplication active, pas l'acte d'infection unique.
    \item \textbf{Oublier le co-récepteur :} Dire "Le VIH se fixe au CD4" est incomplet. Sans co-récepteur (CCR5 ou CXCR4), \textbf{pas de fusion, pas d'infection}. C'est la base des inhibiteurs de co-récepteur (Maraviroc, anti-CCR5).
    \item \textbf{Penser que l'ARN viral sert de messager (ARNm) directement :} Faux. Il est le \textbf{matriciel} pour la synthèse d'ADN par la TI. Les ARNm viraux (pour faire les protéines) ne seront produits \textbf{qu'après} intégration de l'ADN proviral et transcription par l'ARN Pol II cellulaire (voir Section 2).
\end{itemize}
\end{attention}

\begin{exoresolu}[Analyse d'une expérience de tropisme viral]
\textbf{Énoncé :} On dispose de trois lignées cellulaires en culture :
\begin{itemize}
    \item Lignée A : Exprime CD4 et CCR5. Pas de CXCR4.
    \item Lignée B : Exprime CD4 et CXCR4. Pas de CCR5.
    \item Lignée C : Exprime CD4 seulement. Ni CCR5, ni CXCR4.
\end{itemize}
On infecte chaque lignée avec deux souches de VIH-1 :
\begin{itemize}
    \item Souche 1 : Isolée d'un patient au stade primaire (asymptomatique). Séquençage de l'enveloppe : boucle V3 typique R5.
    \item Souche 2 : Isolée d'un patient au stade SIDA. Séquençage de l'enveloppe : boucle V3 typique X4.
\end{itemize}
On mesure la production de p24 au bout de 7 jours.

\textbf{Résultats :}
\begin{center}
\begin{tabularx}{\linewidth}{|l|X|X|X|}\hline
 & \textbf{Lignée A (CD4+CCR5)} & \textbf{Lignée B (CD4+CXCR4)} & \textbf{Lignée C (CD4 seul)} \\ \hline
\textbf{Souche 1 (R5)} & \textbf{Forte (+++)} & Nulle (-) & Nulle (-) \\ \hline
\textbf{Souche 2 (X4)} & Nulle (-) & \textbf{Forte (+++)} & Nulle (-) \\ \hline
\end{tabularx}
\end{center}

\textbf{Questions :}
\begin{enumerate}
    \item Interpréter les résultats pour la Souche 1 et la Souche 2.
    \item Expliquer pourquoi la Lignée C n'est pas infectée malgré la présence du CD4.
    \item Relier ces résultats à l'évolution naturelle de l'infection chez le patient (stade primaire vs SIDA).
\end{enumerate}

\tcblower
\textbf{Solution détaillée :}
\begin{enumerate}
    \item \textbf{Souche 1 (R5) :} Infecte uniquement la Lignée A (CD4+CCR5). C'est une souche \textbf{R5-tropique (M-tropique)}. Elle utilise le co-récepteur \textbf{CCR5}. Elle n'infecte pas la B (pas de CCR5) ni la C (pas de co-récepteur).
    \item \textbf{Souche 2 (X4) :} Infecte uniquement la Lignée B (CD4+CXCR4). C'est une souche \textbf{X4-tropique (T-tropique)}. Elle utilise le co-récepteur \textbf{CXCR4}. Elle n'infecte pas la A (pas de CXCR4) ni la C.
    \item \textbf{Lignée C (CD4 seul) :} Le CD4 est \textbf{nécessaire mais non suffisant}. La fixation de la gp120 au CD4 induit un changement conformationnel qui \textbf{expose le site de liaison au co-récepteur}. Sans co-récepteur (CCR5 ou CXCR4), la gp41 ne peut pas s'insérer dans la membrane cellulaire $\rightarrow$ \textbf{pas de fusion, pas d'entrée}. Le virion reste bloqué à l'extérieur ou est internalisé par endocytose non productive.
    \item \textbf{Évolution naturelle :} Au stade \textbf{primaire/asymptomatique}, le virus transmis est quasi exclusivement \textbf{R5-tropique} (sélection à la transmission : les cellules cibles accessibles au niveau muqueux sont riches en CCR5 : macrophages, cellules dendritiques, lymphocytes T mémoire CCR5+). Au stade \textbf{SIDA}, une pression de sélection (disponibilité des cibles, réponse immune) favorise l'émergence de variants \textbf{X4-tropiques} (utilisant CXCR4 sur lymphocytes T naïfs et précurseurs), corrélée à une \textbf{chute brutale des CD4} et une progression rapide. La Lignée A modélise la cible précoce (macrophages/LT mémoire), la B la cible tardive (LT naïfs/précurseurs).
\end{enumerate}
\end{exoresolu}

\courssub{Bilan de la section 1}

\begin{aretenir}[Points clés à retenir : Identification, Structure, Cibles, Entrée]
\begin{enumerate}
    \item \textbf{Identité :} VIH-1 (pandémique) et VIH-2 (Afrique Ouest) sont des \textbf{Rétrovirus} (Famille \emph{Retroviridae}, Genre \emph{Lentivirus}). Découverts en 1983/84 (Montagnier/Barré-Sinoussi, Gallo). Nobel 2008.
    \item \textbf{Structure du virion (~100 nm) :} Enveloppe lipidique (gp120/gp41) $\rightarrow$ Matrice p17 $\rightarrow$ Capside conique p24 $\rightarrow$ 2 ARN (+) + TI + Intégrase + Protéase.
    \item \textbf{Cibles :} Cellules exprimant \textbf{CD4} + \textbf{Co-récepteur} (CCR5 ou CXCR4). Principalement : \textbf{Lymphocytes T CD4+}, \textbf{Macrophages}, \textbf{Cellules dendritiques}.
    \item \textbf{Entrée :} gp120-CD4 $\rightarrow$ Changement conformationnel $\rightarrow$ gp120-Co-récepteur $\rightarrow$ Insertion gp41 (peptide de fusion) $\rightarrow$ Formation "hairpin" $\rightarrow$ \textbf{Fusion membranes} $\rightarrow$ Libération capside dans cytoplasme.
    \item \textbf{Devenir ARN :} L'ARN viral (+) n'est pas un ARNm fonctionnel. Il \textbf{doit} être rétrotranscrit en ADN double brin par la \textbf{Transcriptase Inverse (TI)} (apportée par le virion) pour pouvoir s'intégrer et être exprimé. Absence de proofreading $\rightarrow$ Taux de mutation élevé ($\sim 3 \times 10^{-5}$/base/cycle) $\rightarrow$ Quasi-espèces virales $\rightarrow$ Nécessité de la trithérapie.
\end{enumerate}
\end{aretenir}

\coursec{Cycle de réplication : de l'ARN viral au nouveau virion infectieux}

Dans la section précédente, nous avons identifié le VIH comme un rétrovirus à enveloppe, doté d'un génome \cle{ARN simple brin à polarité positive} et ciblant préférentiellement les lymphocytes \cle{T CD4$^+$} grâce à l'interaction gp120/CD4 et aux co-récepteurs CCR5/CXCR4. Nous allons maintenant suivre pas à pas ce qui se passe \emph{après} la fusion de l'enveloppe virale avec la membrane plasmique de la cellule cible : comment l'ARN viral devient un ADN intégré au génome de l'hôte, puis comment ce \cle{provirus} détourne la machinerie cellulaire pour produire des milliers de nouveaux virions infectieux. Cette compréhension est la clé pour saisir l'action des traitements antirétroviraux (ARV) utilisés au Cameroun et dans le monde.

\courssub{Pénétration et décalottage : le point de départ}

Après la fusion des membranes, le \cle{capside} virale (protéine p24) libère son contenu dans le \cle{cytoplasme} : deux brins d'ARN génomique identiques (+), associés aux enzymes virales \cle{Transcriptase Inverse (TI)}, \cle{Intégrase (IN)} et \cle{Protéase (PR)}, ainsi que des protéines nucléocapsidiques (p7) et de la matrice (p17). C'est le début de la course contre la montre pour le virus : il doit convertir son ARN en ADN avant que les mécanismes de défense innée de la cellule (restriction factors comme APOBEC3G, SAMHD1) ne le dégradent.

\begin{popfigure}
\begin{tikzpicture}[
    node distance=1.2cm and 1.5cm,
    box/.style={rectangle, draw=popdark, thick, fill=popblueL, rounded corners, minimum height=1.2cm, text width=3.5cm, align=center, font=\small},
    arrow/.style={->, thick, popdark},
    labelbox/.style={rectangle, draw=poporange, fill=poporangeL, rounded corners, font=\scriptsize, text width=3cm, align=center, minimum height=0.8cm}
]
% Cytoplasm zone
\node[box, align=center] (entry){Pénétration / Fusion\\Membrane plasmique\\Libération du capside};
\node[box, below=of entry, align=center] (uncoating){Décalottage (Uncoating)\\Libération : 2 ARN (+), TI, IN, PR};
\node[box, below=of uncoating, fill=popgreenL, align=center] (rt){ÉTAPE 1 : Transcription Inverse\\Cytoplasme\\ARN (+) $\to$ ADN double brin};
\node[box, below=of rt, fill=poppurpleL, align=center] (int){ÉTAPE 2 : Intégration\\Noyau (pore nucléaire)\\ADN viral $\to$ ADN hôte};
\node[box, below=of int, fill=poppinkL, align=center] (transc){ÉTAPE 3 : Transcription\\Noyau\\Provirus $\to$ ARN génomique + ARNm};
\node[box, below=of transc, fill=popgoldL, align=center] (transl){ÉTAPE 4 : Traduction\\Cytoplasme (Ribosomes)\\Polyprotéines Gag, Gag-Pol, Env};
\node[box, below=of transl, fill=popblueL, align=center] (assem){ÉTAPE 5 : Assemblage / Bourgeonnement\\Membrane plasmique\\Particule immature};
\node[box, below=of assem, fill=popgreenL, align=center] (mat){ÉTAPE 6 : Maturation\\Extracellulaire / Fin bourgeonnement\\Clivage Protéase $\to$ Virion mature};
% ARV Targets
\node[labelbox, right=of rt, xshift=1cm, align=center] (arv1){Cibles ARV :\textbf{INRTI} (AZT, TDF)\\\textbf{INNRTI} (NVP, EFV)};
\node[labelbox, right=of int, xshift=1cm] (arv2) {Cible ARV :\textbf{II} (DTG, RAL)};
\node[labelbox, right=of transl, xshift=1cm] (arv4) {Cible ARV :\textbf{IP} (LPV/r, DRV/r, ATV/r)};
\node[labelbox, right=of mat, xshift=1cm, align=center] (arv5){Cible ARV :\textbf{IP}\\(Maturation finale)};
% Arrows
\draw[arrow] (entry) -- (uncoating);
\draw[arrow] (uncoating) -- (rt);
\draw[arrow] (rt) -- (int);
\draw[arrow] (int) -- (transc);
\draw[arrow] (transc) -- (transl);
\draw[arrow] (transl) -- (assem);
\draw[arrow] (assem) -- (mat);
% Compartment labels
\node[font=\bfseries\small, text=popdark, left=of rt, xshift=-1.5cm] {Cytoplasme};
\node[font=\bfseries\small, text=popdark, left=of int, xshift=-1.5cm] {Noyau};
\node[font=\bfseries\small, text=popdark, left=of transc, xshift=-1.5cm] {Noyau};
\node[font=\bfseries\small, text=popdark, left=of transl, xshift=-1.5cm] {Cytoplasme};
\node[font=\bfseries\small, text=popdark, left=of assem, xshift=-1.5cm] {Membrane};
\node[font=\bfseries\small, text=popdark, left=of mat, xshift=-1.5cm] {Extracellulaire};
\end{tikzpicture}
\legende{Fluxogramme synthétique du cycle de réplication du VIH-1 avec localisation subcellulaire et cibles des classes d'antirétroviraux (ARV). INRTI = Inhibiteurs Nucléosidiques de la Transcriptase Inverse ; INNRTI = Inhibiteurs Non Nucléosidiques de la TI ; II = Inhibiteurs de l'Intégrase ; IP = Inhibiteurs de la Protéase.}
\end{popfigure}

\courssub{Étape 1 : La Transcription Inverse (Cytoplasme) -- De l'ARN à l'ADN proviral}

C'est l'étape signature des \cle{Rétrovirus} (famille \emph{Retroviridae}). Elle est catalysée par la \cle{Transcriptase Inverse (TI)}, une enzyme à activité double : \cle{ADN polymérase ARN-dépendante} et \cle{Ribonucléase H (RNase H)}. La TI du VIH-1 est un \cle{dimère} hétérodimérique composé d'une sous-unité \cle{p66} (active, sites polymérase et RNase H) et d'une sous-unité \cle{p51} (structurelle, issue du clivage de p66 par la Protéase).

\begin{popfigure}
\begin{tikzpicture}[
    node distance=0.8cm and 1cm,
    mol/.style={rectangle, draw=popdark, thick, fill=popcream, rounded corners, minimum height=0.8cm, text width=3.2cm, align=center, font=\small},
    enz/.style={ellipse, draw=poppurple, thick, fill=poppurpleL, minimum height=0.8cm, text width=2.5cm, align=center, font=\small\bfseries},
    arrow/.style={->, thick, popdark},
    label/.style={font=\scriptsize, text=popdark, midway, above, sloped}
]
% Substrates/Products
\node[mol, align=center] (rna){ARN génomique (+)\\Modèle};
\node[enz, right=of rna, align=center] (ti){Transcriptase Inverse\\(p66/p51)};
\node[mol, right=of ti, align=center] (dna1){ADN (-) simple brin\\Hybride ARN/ADN};
\node[mol, below=of dna1, align=center] (rna_degrad){Fragments d'ARN\\dégradés (RNase H)};
\node[mol, right=of dna1, align=center] (dna2){ADN (+) simple brin\\Synthèse second brin};
\node[mol, below=of dna2, align=center] (proviral){ADN double brin linéaire\\(ADN proviral / pré-intégration)};
% Arrows and labels
\draw[arrow] (rna) -- node[label] {Modèle} (ti);
\draw[arrow] (ti) -- node[label] {Polymérase} (dna1);
\draw[arrow, poporange] (ti) -- node[label, text=poporange] {RNase H} (rna_degrad);
\draw[arrow] (dna1) -- node[label] {Modèle} (ti);
\draw[arrow] (ti) -- node[label] {Polymérase} (dna2);
\draw[arrow] (dna1) -- (dna2);
\draw[arrow] (dna2) -- (proviral);
% Primer tRNA
\node[mol, above=of rna, xshift=-1cm, fill=popblueL, text width=2.5cm] (primer) {Amorce : tRNA$^{Lys3}$\\(primer binding site PBS)};
\draw[arrow, dashed] (primer) -- (rna);
\end{tikzpicture}
\legende{Mécanisme de la Transcriptase Inverse (TI). 1. Fixation sur l'ARN (+) amorcé par un tRNA$^{Lys3}$ cellulaire. 2. Synthèse de l'ADN (-) (activité polymérase). 3. Dégradation de l'ARN modèle (activité RNase H) sauf les séquences PPT (Polypurine Tract) qui servent d'amorce pour le brin (+). 4. Synthèse de l'ADN (+). 5. Formation de l'ADN double brin linéaire (ADN proviral) avec les LTR (Long Terminal Repeats) aux extrémités.}
\end{popfigure}

\begin{definition}[ADN Proviral (Forme linéaire pré-intégration)]
L'\cle{ADN proviral} est la forme double brin linéaire issue de la transcription inverse. Il porte aux deux extrémités des séquences identiques appelées \cle{LTR (Long Terminal Repeats)}, essentielles pour l'intégration (sites de reconnaissance par l'Intégrase) et la transcription future (promoteur, enhanceur, signaux de polyadénylation). Il n'est pas encore intégré au génome hôte.
\end{definition}

\begin{propriete}[Absence de relecture (Proofreading) et variabilité génétique]
La Transcriptase Inverse du VIH \textbf{ne possède pas d'activité exonucléasique 3'$\to$5' (relecture)}. Son taux d'erreur est d'environ \textbf{$3 \times 10^{-5}$ erreurs par base et par cycle de réplication} (contre $10^{-8}$ à $10^{-10}$ pour les ADN polymérases cellulaires). Couplé à une production massive de virions ($\sim 10^{10}$/jour) et à la \cle{recombinaison} fréquente lors de la co-infection d'une cellule par deux souches différentes (changement de modèle par la TI), cela génère une \cle{variabilité génétique extrême} (quasi-espèces).
\par \textbf{Conséquences majeures :} Échappement à la réponse immune (anticorps, CTL), résistance rapide aux ARV en monothérapie, difficulté vaccinale. D'où la \textbf{nécessité absolue de la trithérapie (ARV combinés)} pour créer une barrière génétique élevée.
\end{propriete}

\begin{attention}[Piège fréquent : Confusion TI / ADN Polymérase cellulaire]
Ne confonds pas la Transcriptase Inverse (enzyme \emph{virale}, ARN-dépendante, sans relecture) et l'ADN Polymérase cellulaire (enzyme \emph{hôte}, ADN-dépendante, avec relecture). La TI est une cible thérapeutique \emph{sélective} : les INRTI (analogues de nucléosides) et INNRTI (non nucléosidiques) bloquent la TI virale sans affecter significativement les polymérases cellulaires aux doses thérapeutiques.
\end{attention}

\paragraph{Cibles thérapeutiques : INRTI et INNRTI.}
\begin{itemize}
    \item \textbf{INRTI (ex: Zidovudine/AZT, Ténéfovir/TDF, Lamivudine/3TC, Emtricitabine/FTC) :} Analogues structuraux des désoxyribonucléotides. Ils s'incorporent à la chaîne d'ADN en elongation mais \textbf{manquent du groupement 3'-OH} nécessaire à l'ajout du nucléotide suivant $\to$ \textbf{arrêt de chaîne (chain termination)}. Ex: AZT (thymidine analogue), TDF (adénosine analogue).
    \item \textbf{INNRTI (ex: Névirapine/NVP, Efavirenz/EFV, Doravirine/DOR, Rilpivirine/RPV) :} Se fixent sur un \textbf{site allostérique} (poche hydrophobe) de la sous-unité p66, loin du site actif $\to$ changement conformationnel $\to$ inhibition non compétitive de l'activité polymérase. Ne s'incorporent \emph{pas} dans l'ADN.
\end{itemize}

\begin{exoresolu}[Interprétation d'un résultat de résistance génotypique]
\textbf{Énoncé :} Un patient sous traitement de première ligne (TDF + 3TC + EFV) au Cameroun présente un échec virologique (charge virale $> 1000$ copies/mL). Le test de résistance révèle la mutation \textbf{M184V} sur la Transcriptase Inverse. Expliquer l'impact de cette mutation sur la sensibilité aux médicaments.
\par \textbf{Solution :}
\begin{enumerate}
    \item \textbf{Identification :} M184V = substitution d'une Méthionine (M) par une Valine (V) en position 184 de la TI. Position clé dans le site actif (poche de liaison des dNTP).
    \item \textbf{Effet sur 3TC/FTC (INRTI cytidine analogues) :} La Valine crée un encombrement stérique qui \textbf{réduit drastiquement l'incorporation} de la 3TC-triphosphate (facteur $> 1000$). \cle{Résistance de haut niveau} à la 3TC et FTC.
    \item \textbf{Effet sur TDF (INRTI adénosine analogue) :} Paradoxalement, M184V \textbf{augmente la sensibilité} au TDF (hypersensibilité) en réduisant le taux d'excision (pyrophosphorolyse) de l'analogue incorporé.
    \item \textbf{Effet sur EFV (INNRTI) :} Aucun effet direct (site de liaison différent).
    \item \textbf{Conclusion clinique :} La 3TC/FTC perd son activité antivirale propre mais \textbf{conserve un intérêt} : la mutation M184V réduit la \emph{fitness} (capacité réplicative) du virus ($\sim 30-50\%$) et restaure la sensibilité au TDF. On maintient souvent la 3TC pour cet effet "coût de fitness" et on change l'INNRTI (ex: passage au Dolutégravir, un II).
\end{enumerate}
\end{exoresolu}

\courssub{Étape 2 : L'Intégration (Noyau) -- La formation du provirus définitif}

L'ADN proviral linéaire, complexe avec l'Intégrase (IN) et des protéines cellulaires (LEDGF/p75, Importine), forme le \cle{Complexe de Pré-Intégration (PIC)}. Ce complexe traverse l'\cle{enveloppe nucléaire} via les \cle{pores nucléaires} (transport actif, indépendant de la mitose : le VIH infecte les cellules non divisées comme les macrophages et les lymphocytes au repos G0/G1).

\begin{popfigure}
\begin{tikzpicture}[
    node distance=1cm and 1.5cm,
    stepbox/.style={rectangle, draw=popdark, thick, fill=popblueL, rounded corners, minimum height=1cm, text width=3.5cm, align=center, font=\small},
    dna/.style={rectangle, draw=popgreen, thick, fill=popgreenL, rounded corners, minimum height=1cm, text width=4cm, align=center, font=\small},
    arrow/.style={->, thick, popdark},
]
\node[stepbox, align=center] (pic){Complexe de Pré-Intégration (PIC)\\ADN viral + Intégrase (IN)\\+ LEDGF/p75 + Importines};
\node[stepbox, below=of pic, align=center] (entry_nuc){Traversée pore nucléaire\\Transport actif (Ran-GTP)};
\node[dna, below=of entry_nuc, align=center] (host_dna){ADN chromosomique hôte\\Chromatine active (LEDGF/p75)};
\node[stepbox, below=of host_dna, align=center] (int_steps){Réactions catalysées par l'Intégrase :\\1. \textbf{3'-Processing} : Coupe 2 nucléotides (GT) aux 3' de l'ADN viral (exposant CA-OH 3').\\2. \textbf{Strand Transfer} : Attaque nucléophile des 3'-OH viraux sur l'ADN hôte $\to$ Coupes décalées (5 pb).\\3. Réparation par enzymes cellulaires (ligase, polymérase) $\to$ Duplication cible 5 pb.};
\node[dna, below=of int_steps, fill=poporangeL, draw=poporange, align=center] (provirus){PROVIRUS INTÉGRÉ\\ADN viral inséré dans chromosome hôte\\Flanké par LTR - Duplication 5 pb cible\\FORME LATENTE OU ACTIVE};
\draw[arrow] (pic) -- (entry_nuc);
\draw[arrow] (entry_nuc) -- (host_dna);
\draw[arrow] (host_dna) -- (int_steps);
\draw[arrow] (int_steps) -- (provirus);
\end{tikzpicture}
\legende{Processus d'intégration. L'Intégrase (IN) réalise le "3'-processing" (coupe des 2 derniers nucléotides GT aux extrémités 3' de l'ADN viral) puis le "Strand Transfer" (jonction covalente des extrémités 3' virales sur l'ADN hôte, créant une duplication de 5 paires de bases de la séquence cible). La réparation des brins nickés est assurée par la machinerie cellulaire.}
\end{popfigure}

\begin{definition}[Provirus Intégré]
Le \cle{provirus} est l'ADN viral double brin \textbf{intégré covalemment} dans un chromosome de la cellule hôte. C'est la forme \textbf{latente} (silencieuse) ou \textbf{active} (transcrite). Il est \textbf{réplicatif} : lors de la division cellulaire (mitose), il est dupliqué et transmis aux cellules filles au même titre que les gènes cellulaires. Il constitue le \cle{réservoir viral} principal, invisible au système immunitaire et inaccessible aux ARV actuels.
\end{definition}

\begin{attention}[Irréversibilité de l'intégration et réservoir]
\textbf{Aucun traitement ARV actuel ne permet d'exciser le provirus intégré.} Les Inhibiteurs de l'Intégrase (II : Dolutégravir/DTG, Raltégravir/RAL, Bictegravir/BIC, Cabotégravir/CAB) bloquent l'étape d'intégration \emph{avant} qu'elle n'ait lieu (ils chélatent les ions Mg$^{2+}$ du site actif de l'IN). Une fois intégré, le provirus est là "pour la vie de la cellule". C'est la raison fondamentale pour laquelle \textbf{le traitement ARV est à vie} (sauf rares cas de "contrôleurs post-traitement" ou greffe de moelle osseuse CCR5$\Delta$32). L'arrêt du traitement entraîne une réplication virale à partir des réservoirs (réactivation transcriptionnelle) en quelques semaines.
\end{attention}

\begin{savaistu}[Inhibiteurs d'entrée et de fixation : Maraviroc, Fostemsavir]
Avant même la transcription inverse, on peut bloquer l'entrée. \textbf{Maraviroc (MVC)} : antagoniste du co-récepteur \textbf{CCR5} (utilisé par les souches R5, majoritaires au stade précoce). Nécessite un test de tropisme viral (coûteux). \textbf{Fostemsavir (FTR)} : inhibiteur de fixation sur gp120 (empêche liaison CD4). \textbf{Enfuvirtide (T-20)} : peptide inhibiteur de fusion (gp41), injectable SC 2x/jour. Au Cameroun, ces molécules sont \textbf{peu utilisées en 1ère ligne} (coût élevé, complexité de suivi, tropisme) mais sont des options précieuses en \textbf{3ème ligne (salvage therapy)} pour les multi-résistants, disponibles dans les centres de référence (ex: Centre Pasteur du Cameroun, Hôpital Central de Yaoundé).
\end{savaistu}

\courssub{Étape 3 : Transcription du provirus (Noyau) -- Déclenchement et régulation}

Une fois intégré, le provirus est soumis au contrôle de la \cle{Polymérase II} cellulaire. Son promoteur (dans le LTR 5') contient des sites de fixation pour des facteurs de transcription cellulaires (\cle{NF-$\kappa$B}, Sp1, NFAT) et viraux (\cle{Tat}).

\begin{itemize}
    \item \textbf{Phase latente / Basale :} Faible transcription initiatée par facteurs cellulaires seuls. Production de quelques ARNm \emph{complètement épissés} (codant pour \cle{Rev}, \cle{Nef}, \cle{Tat}). \textbf{Tat} (Trans-activator of transcription) se fixe sur \cle{TAR} (Trans-activation Response element) à l'extrémité 5' des ARNs naissants $\to$ recrute P-TEFb (CDK9/Cycline T1) $\to$ \textbf{élongation processive} $\to$ amplification massive (boucle de rétroaction positive).
    \item \textbf{Activation cellulaire :} Signaux (antigènes, cytokines IL-2, TNF-$\alpha$, co-infections) $\to$ activation NF-$\kappa$B $\to$ translocation nucléaire $\to$ forte transcription.
    \item \textbf{Produits :}
    \begin{itemize}
        \item \textbf{ARN génomiques (non épissés, $\sim$ 9 kb) :} Servent de génome pour nouveaux virions + ARNm pour \cle{Gag} (p55) et \cle{Gag-Pol} (p160, lecture traversante \emph{frameshift} -1).
        \item \textbf{ARNm simple épissés ($\sim$ 4 kb) :} Codent \cle{Env} (gp160), \cle{Vif}, \cle{Vpr}, \cle{Vpu}.
        \item \textbf{ARNm multi-épissés ($\sim$ 2 kb) :} Codent \cle{Tat}, \cle{Rev}, \cle{Nef}.
    \end{itemize}
    \item \textbf{Rôle de Rev :} Se fixe sur \cle{RRE} (Rev Response Element) dans l'intron des ARNs longs $\to$ export nucléaire via CRM1 (exportine) $\to$ permet la traduction des protéines structurales (Gag, Pol, Env) dans le cytoplasme. Sans Rev : rétention nucléaire et dégradation des ARNs longs.
\end{itemize}

\courssub{Étape 4 : Traduction, Maturation et Assemblage (Cytoplasme / Membrane)}

\begin{enumerate}
    \item \textbf{Traduction :} Sur ribosomes libres (Gag, Gag-Pol) ou ribosomes liés au RE (Env $\to$ gp160 clivé en gp120/gp41 par furine dans appareil de Golgi $\to$ transport vésiculaire vers membrane plasmique).
    \item \textbf{Adressage :} Domaine MA (p17) de Gag (myristoylé) + domaine basique (NC, p7) $\to$ ciblage spécifique vers \cle{rafts lipidiques} (domaines riches en cholestérol/sphingolipides) de la membrane plasmique.
    \item \textbf{Assemblage :} Multimérisation de \textbf{Gag (p55)} $\to$ recrutement de \textbf{2 ARN génomiques} (dimérisation via DIS, Dimer Initiation Site, stabilisée par NC/p7) + enzymes (TI, IN, PR) via Gag-Pol $\to$ formation d'une \textbf{particule immature} (sphérique, non infectieuse, capside non formée).
    \item \textbf{Bourgeonnement (Budding) :} Domaine \textbf{L (Late) de p6 (PTAP)} $\to$ recrutement complexe \textbf{ESCRT} (Endosomal Sorting Complex Required for Transport) cellulaire $\to$ fission membranaire $\to$ libération du virion immature dans l'espace extracellulaire.
    \item \textbf{MATURATION (Étape critique) :} Activation de la \cle{Protéase virale (PR)} par auto-clivage (dimérisation) $\to$ clivage séquentiel et ordonné des polyprotéines :
    \[
    \text{Gag (p55)} \xrightarrow{\text{PR}} \text{MA (p17)} + \text{CA (p24)} + \text{NC (p7)} + \text{p6} + \text{peptides spacer (SP1, SP2)}
    \]
    \[
    \text{Gag-Pol (p160)} \xrightarrow{\text{PR}} \text{PR} + \text{RT (p66/p51)} + \text{IN (p32)}
    \]
    \textbf{Résultat :} Réorganisation structurelle $\to$ formation du \cle{cône capsidique} (CA/p24) mature $\to$ \textbf{Virion mature, infectieux}.
\end{enumerate}

\begin{propriete}[Cible thérapeutique : Inhibiteurs de la Protéase (IP)]
Les \textbf{IP (Lopinavir/r, Darunavir/r, Atazanavir/r, Atazanavir/cobicistat)} sont des \textbf{analogues de l'état de transition} (peptidomimétiques) qui se lient \textbf{compétitivement} au site actif de la Protéase (aspartyl-protéase, dimère) $\to$ bloquent le clivage de Gag et Gag-Pol.
\par \textbf{Conséquence :} Production de virions \textbf{immatures, non infectieux} (morphologie "déficiente", capside non condensée).
\par \textbf{Pharmacocinétique :} Faible biodisponibilité, métabolisme hépatique intense (CYP3A4). \textbf{Toujours administrés avec un "booster" pharmacologique} : \textbf{Ritonavir (r)} (faible dose 100 mg) ou \textbf{Cobicistat (c)} (150 mg) $\to$ inhibition CYP3A4 $\to$ augmentation exposition (AUC) de l'IP principal. \textbf{Attention :} Interactions médicamenteuses majeures (rifampicine, contraceptifs oraux, etc.).
\end{propriete}

\courssub{Cinétique de l'infection in vitro : Données expérimentales}

L'étude de cultures de lymphocytes CD4$^+$ (lignées MT-4, CEM ou PBMC primaires) infectées par une multiplicité d'infection (MOI) connue permet de suivre la dynamique temporelle des marqueurs viraux.

\begin{popfigure}
\begin{tikzpicture}
\begin{axis}[
    width=\linewidth, height=9cm,
    xlabel={Temps post-infection (heures)}, ylabel={Quantité relative (log échelle)},
    xmin=0, xmax=72, ymin=0, ymax=7,
    xtick={0,12,24,36,48,60,72}, ytick={0,1,2,3,4,5,6,7},
    legend style={at={(0.02,0.98)}, anchor=north west, fill=popcream, draw=popdark, font=\small},
    grid=major, grid style={popline}
]
% Données simulées représentatives (log10)
% ADN viral total (linéaire + circulaire + intégré) : détection précoce ~6-8h, pic ~24-36h
\addplot[popcurve, domain=0:72, samples=100, thick] {2.5*(1-exp(-(x-4)/6)) * (x>4) + 0.5*exp(-(x-36)/20) * (x>36) + 2.5};
\addlegendentry{ADN viral total (PCR)}
% ADN intégré (Alu-PCR) : détection ~12-18h, plateau
\addplot[poporange, thick, domain=0:72, samples=100] {2.0*(1-exp(-(x-10)/8)) * (x>10) + 2.0};
\addlegendentry{ADN intégré (Alu-PCR)}
% ARN viral cellulaire : suit transcription, pic ~24-48h
\addplot[popgreen, thick, domain=0:72, samples=100] {3.0*(1-exp(-(x-14)/10)) * (x>14) * exp(-(x-48)/30) + 0.5};
\addlegendentry{ARN viral cellulaire (RT-qPCR)}
% p24 extracellulaire (ELISA) : tardif, sécrétion, pic ~48-72h
\addplot[poppurple, thick, domain=0:72, samples=100] {3.5*(1-exp(-(x-24)/12)) * (x>24) * (1 - 0.2*exp(-(x-60)/10))};
\addlegendentry{p24 extracellulaire (ELISA)}
\end{axis}
\end{tikzpicture}
\legende{Cinétique typique de l'infection par le VIH-1 in vitro (lymphocytes CD4$^+$, MOI faible). \textbf{ADN viral total} (courbe bleue) : détection précoce (produits de transcription inverse linéaires et circulaires 1-LTR/2-LTR). \textbf{ADN intégré} (orange) : suit avec un décalage (nécessite transport nucléaire + intégration), plateau stable. \textbf{ARN viral} (vert) : pic transcriptionnel. \textbf{p24 extracellulaire} (violet) : marqueur tardif de production virale (assemblage/bourgeonnement/maturation). Échelle logarithmique.}
\end{popfigure}

\begin{methode}[Reconstituer la chronologie du cycle à partir de données moléculaires]
\textbf{Objectif :} Ordonner les étapes (Transcription inverse, Intégration, Transcription, Traduction/Assemblage, Maturation) à partir de résultats expérimentaux (Western Blot, PCR, ELISA) obtenus à différents temps post-infection.
\par \textbf{Démarche :}
\begin{enumerate}
    \item \textbf{Identifier les marqueurs précoces (Cytoplasme/Noyau, 6-24h) :} Détection d'\textbf{ADN viral total} (PCR sur ADN extrait total) $\to$ signature de la \textbf{Transcription Inverse}. Détection d'\textbf{ADN circulaire 2-LTR} (PCR spécifique) $\to$ marqueur d'échec d'intégration (nucléaire, non intégré).
    \item \textbf{Identifier le marqueur d'intégration définitive (Noyau, 12-24h) :} \textbf{Alu-PCR} (amplification jonction ADN viral / ADN Alu hôte) $\to$ preuve de l'\textbf{Intégration}.
    \item \textbf{Identifier les marqueurs d'expression génique (Noyau/Cytoplasme, 18-48h) :} \textbf{ARN viral cellulaire} (RT-qPCR) $\to$ \textbf{Transcription}. Protéines virales intracellulaires (p24 intracellulaire, Western Blot) $\to$ \textbf{Traduction}.
    \item \textbf{Identifier les marqueurs tardifs (Extracellulaire, 36-72h) :} \textbf{p24 extracellulaire} (ELISA sur surnageant) + \textbf{ARN viral extracellulaire} (Charge virale) $\to$ \textbf{Assemblage, Bourgeonnement, Maturation, Libération}.
    \item \textbf{Valider par inhibition :} Ajout d'ARV spécifiques (AZT $\to$ bloque ADN viral total ; DTG $\to$ bloque ADN intégré mais pas ADN total ; LPV/r $\to$ bloque p24 extracellulaire mature mais pas p24 intracellulaire immature).
\end{enumerate}
\end{methode}

\begin{exemplebox}[Exemple chiffré : Turnover viral massif]
Chez un patient non traité, charge virale plasmatique $\sim 10^5$ copies/mL. Volume plasmatique $\sim 3$ L.
\begin{itemize}
    \item Virions circulants totaux $\approx 3 \times 10^8$ particules.
    \item \textbf{Demi-vie des virions libres in vivo :} $t_{1/2} \approx \textbf{6 heures}$ (études de décroissance sous ARV puissants).
    \item \textbf{Production quotidienne :} Pour maintenir l'état stationnaire, il faut produire $\approx \frac{\ln 2}{t_{1/2}} \times \text{Pool total} \approx \frac{0.693}{6\text{h}} \times 3 \times 10^8 \times 24\text{h} \approx \textbf{$8 \times 10^9$ à $10^{10}$ virions/jour}$.
    \item \textbf{Turnover CD4$^+$ :} Chaque virion produit résulte de l'infection/mort d'une cellule (majoritairement). Turnover massif des lymphocytes CD4$^+$ activés ($\sim 10^9$ à $10^{10}$ cellules/jour détruites/remplacées). Épuisement progressif du stock de cellules naïves et mémoire $\to$ immunodépression.
\end{itemize}
\emph{Source : Modélisation mathématique (Ho et al., Nature 1995 ; Wei et al., Nature 1995) confirmée cliniquement.}
\end{exemplebox}

\begin{aretenir}[Résumé du cycle réplicatif et cibles ARV]
\begin{center}
\begin{tabularx}{\linewidth}{|l|X|X|X|}\hline
\textbf{Étape} & \textbf{Localisation} & \textbf{Événements clés / Enzymes} & \textbf{Cibles ARV (Classes)} \\ \hline
1. Fixation/Entrée & Membrane plasmique & gp120/CD4/Co-récepteur $\to$ Fusion gp41 & IE (Maraviroc/CCR5), IF (Fostemsavir/gp120), IFu (Enfuvirtide/gp41) \\ \hline
2. Décalottage & Cytoplasme & Libération ARN, TI, IN, PR & -- \\ \hline
3. \textbf{Transcription Inverse} & \textbf{Cytoplasme} & \textbf{TI (p66/p51) : ARN $\to$ ADN ds (LTR)} & \textbf{INRTI (AZT, TDF, 3TC), INNRTI (EFV, NVP, DOR)} \\ \hline
4. Transport nucléaire & Pore nucléaire & PIC (IN, LEDGF/p75, Importines) & -- \\ \hline
5. \textbf{Intégration} & \textbf{Noyau (Chromatine)} & \textbf{IN : 3'-Processing + Strand Transfer} $\to$ \textbf{Provirus} & \textbf{II (DTG, RAL, BIC, CAB)} \\ \hline
6. \textbf{Transcription} & \textbf{Noyau} & \textbf{Pol II hôte + Tat/NF-$\kappa$B} $\to$ ARN génomique + ARNm & -- (Recherche : inhibiteurs Tat, BRD4) \\ \hline
7. Export ARN & Pore nucléaire & \textbf{Rev/RRE/CRM1} & -- \\ \hline
8. \textbf{Traduction} & \textbf{Cytoplasme (Ribosomes)} & \textbf{Gag (p55), Gag-Pol (p160), Env (gp160)} & -- \\ \hline
9. Assemblage/Bourgeonnement & Membrane (Rafts) & Gag + ARNg + Enzymes $\to$ Particule immature $\to$ ESCRT & -- (Recherche : inhibiteurs Gag-Gag, ESCRT) \\ \hline
10. \textbf{Maturation} & \textbf{Extracellulaire} & \textbf{Protéase (PR) : Clivage Gag/Gag-Pol} $\to$ Virion mature & \textbf{IP (LPV/r, DRV/r, ATV/r, ATV/c)} \\ \hline
\end{tabularx}
\end{center}
\end{aretenir}

\begin{exercice}[Application : Chronologie et cibles thérapeutiques]
Un chercheur infecte des lymphocytes CD4$^+$ in vitro et ajoute différents inhibiteurs au moment de l'infection (temps 0). Il mesure l'ADN viral total (PCR), l'ADN intégré (Alu-PCR) et le p24 extracellulaire (ELISA) à 48h.
\begin{enumerate}
    \item Prédire le résultat (Augmenté / Diminué / Inchangé / Absent) pour chaque marqueur dans les conditions suivantes :
    \begin{itemize}
        \item Témoin (sans inhibiteur).
        \item + AZT (INRTI) ajouté à T0.
        \item + Dolutégravir (II) ajouté à T0.
        \item + Lopinavir/r (IP) ajouté à T0.
    \end{itemize}
    \item Justifier chaque prédiction en citant l'étape bloquée et l'enzyme ciblée.
    \item Pourquoi l'ajout de Dolutégravir à T0 ne diminue-t-il pas l'ADN viral total mesuré à 48h, alors qu'il abolit l'ADN intégré ?
\end{enumerate}
\end{exercice}

\begin{corrige}[Exercice : Chronologie et cibles thérapeutiques]
\begin{enumerate}
    \item \textbf{Tableau de résultats attendus à 48h :}
    \begin{center}
    \begin{tabular}{|l|c|c|c|}\hline
    \textbf{Condition} & \textbf{ADN viral total} & \textbf{ADN intégré} & \textbf{p24 extracellulaire} \\ \hline
    Témoin & +++ (Positif) & +++ (Positif) & +++ (Positif) \\ \hline
    + AZT (INRTI) & \textbf{Absent / Très faible} & \textbf{Absent} & \textbf{Absent} \\ \hline
    + DTG (II) & +++ (Positif) & \textbf{Absent} & \textbf{Absent / Très faible} \\ \hline
    + LPV/r (IP) & +++ (Positif) & +++ (Positif) & \textbf{Absent / Très faible (immature)} \\ \hline
    \end{tabular}
    \end{center}
    \item \textbf{Justifications :}
    \begin{itemize}
        \item \textbf{AZT :} Bloque la \textbf{Transcriptase Inverse} (Étape 3). Pas de synthèse ADN $\to$ pas de substrat pour l'intégration $\to$ pas de provirus $\to$ pas de transcription $\to$ pas de protéines $\to$ pas de virions.
        \item \textbf{DTG :} Bloque l'\textbf{Intégrase} (Étape 5). La Transcription Inverse a lieu normalement $\to$ ADN viral total (linéaire + circulaires 2-LTR) détecté. Mais pas d'intégration $\to$ pas de provirus stable $\to$ transcription très faible/transitoire (depuis ADN non intégré instable) $\to$ pas de production virale soutenue.
        \item \textbf{LPV/r :} Bloque la \textbf{Protéase} (Étape 10). Tout l'amont se passe normalement (Intégration, Transcription, Traduction, Assemblage, Bourgeonnement). Des particules \textbf{immatures} sont libérées (p24 intracellulaire présent, p24 extracellulaire détectable par ELISA mais \textbf{non infectieuses}, morphologie anormale au ME). Note : certains ELISA détectent le p24 immature, d'autres non ; la charge virale ARN reste élevée mais non infectieuse.
    \end{itemize}
    \item \textbf{Explication question 3 :} L'ADN viral total mesuré par PCR standard amplifie des séquences présentes sur l'\textbf{ADN linéaire (pré-intégration)} et les \textbf{formes circulaires (1-LTR, 2-LTR)}. Ces formes sont des \textbf{produits "cul-de-sac" de la transcription inverse} qui s'accumulent dans le noyau lorsque l'intégration est bloquée (elles ne sont pas dégradées immédiatement). L'Alu-PCR, elle, est spécifique de la jonction virus-hôte, donc négative si l'Intégrase est inhibée.
\end{enumerate}
\end{corrige}

\coursec{Dynamique de l'infection : Phases cliniques, marqueurs biologiques et physiopathologie}

Nous venons d'étudier comment le VIH détourne la machinerie cellulaire du lymphocyte T CD4+ pour se répliquer (Section 2). Mais une infection ne se résume pas à un cycle cellulaire unique : c'est une bataille dynamique à l'échelle de l'organisme entier, qui se joue sur plusieurs années. Comprendre cette dynamique, c'est comprendre pourquoi le malade peut sembler guérir puis rechuter, pourquoi les marqueurs biologiques évoluent en « ciseaux », et comment le clinicien décide d'intervenir. Observons d'abord l'évolution naturelle de l'infection non traitée.

\courssub{Les trois phases cliniques de l'infection VIH non traitée}

L'histoire naturelle de l'infection par le VIH-1 (le type le plus fréquent au Cameroun) se divise en trois phases successives, caractérisées par l'évolution de la charge virale (quantité de virus dans le sang) et du taux de lymphocytes T CD4+.

\begin{experience}[Preuve expérimentale : Suivi de cohortes de patients]
\textbf{Observation :} Dès les années 1980, des groupes de patients (cohortes) ont été suivis régulièrement avec des prises de sang tous les 6 mois pour mesurer la charge virale et le taux de CD4+.

\textbf{Résultats clés :}
\begin{itemize}
    \item La charge virale plasmatique (ARN VIH/mL) suit une courbe en « pic - plateau - remontée terminale ».
    \item Le taux de CD4+ sanguins suit une courbe en « chute brutale - déclin lent - effondrement ».
    \item Les deux courbes se croisent : on parle d'évolution en « \textbf{ciseaux} ».
\end{itemize}

\textbf{Interprétation :} La réplication virale n'est jamais complètement arrêtée pendant la phase « asymptomatique » ; elle est active en continu. Le système immunitaire contrôle partiellement l'infection au début, mais s'épuise progressivement à force de lutter.
\legende{Preuves cliniques fondant la description triphasique de l'infection VIH.}
\end{experience}

\begin{popfigure}
\begin{tikzpicture}
\begin{axis}[
    width=\linewidth, height=0.6\linewidth,
    xlabel={Temps (années)}, 
    ylabel={Charge virale (Log$_{10}$ copies/mL)},
    ymin=0, ymax=7,
    xmin=0, xmax=12,
    xtick={0,1,2,3,4,5,6,7,8,9,10,11,12},
    ytick={0,1,2,3,4,5,6,7},
    grid=major,
    axis y line*=left,
    tick label style={font=\small},
    label style={font=\small},
    legend style={font=\small, at={(0.98,0.98)}, anchor=north east, cells={anchor=west}},
]
% Zones colorées phases
\addplot[popblueL!20, draw=none, fill=popblueL!20] coordinates {(0,0) (0.5,0) (0.5,7) (0,7)} \closedcycle;
\addplot[poporangeL!20, draw=none, fill=poporangeL!20] coordinates {(0.5,0) (0.5,7) (9,7) (9,0)} \closedcycle;
\addplot[popred!15, draw=none, fill=popred!15] coordinates {(9,0) (9,7) (12,7) (12,0)} \closedcycle;
% Courbe Charge Virale (Log10) - Axe gauche
\addplot[popblue, very thick, domain=0:0.5, samples=50] {6.5 - 5.5*x/0.5}; % Pic
\addplot[popblue, very thick, domain=0.5:9, samples=100] {4.0 + 0.3*sin(deg(3*x))}; % Plateau (Set point)
\addplot[popblue, very thick, domain=9:12, samples=50] {4.0 + 2.5*(x-9)/3}; % Remontée terminale
% Courbe CD4+ - Axe droit (second axis)
% On utilise un second axe pour l'échelle réelle des CD4
\end{axis}
% Second axe pour CD4 (superposé)
\begin{axis}[
    width=\linewidth, height=0.6\linewidth,
    xlabel={Temps (années)}, 
    ylabel={Taux CD4+ (cellules/$\mu$L)},
    ymin=0, ymax=1200,
    xmin=0, xmax=12,
    xtick={0,1,2,3,4,5,6,7,8,9,10,11,12},
    ytick={0,200,400,600,800,1000,1200},
    grid=none,
    axis y line*=right,
    axis x line=none,
    tick label style={font=\small},
    label style={font=\small},
    legend style={font=\small, at={(0.98,0.02)}, anchor=south east, cells={anchor=west}},
    hide x axis
]
% Courbe CD4+
\addplot[poporange, very thick, dashed, domain=0:0.25, samples=30] {1000 - 2000*x/0.25}; % Chute brutale
\addplot[poporange, very thick, dashed, domain=0.25:9, samples=100] {500 - 70*(x-0.25)}; % Déclin lent ~70/an
\addplot[poporange, very thick, dashed, domain=9:11, samples=30] {200 - 100*(x-9)}; % Effondrement
\addlegendentry{Charge Virale (Log$_{10}$) -- Axe gauche}
\addlegendentry{Taux CD4+ -- Axe droit}
% Annotations phases (placées sur le premier axe via coordinates relatives ou node absolus)
% On utilise des nodes TikZ absolus par rapport au premier axis environment pour simplicité
\end{axis}
% Annotations manuelles (approximatives en coordonnées relatives page) -> Mieux : les mettre dans le premier axis via \node at (axis cs:...)
% Je les ajoute dans le premier axis ci-dessus via \node mais ici le code est séparé. 
% Solution : Tout mettre dans un seul axis avec deux y axes ? Non, pgfplots gère mal deux courbes échelles différentes dans une légende unique facilement.
% Alternative : Un seul graphique schématique simplifié sans données numériques exactes mais forme qualitative.
\end{tikzpicture}
\legende{Évolution schématique type de la charge virale (bleu, trait plein, échelle logarithmique à gauche) et du taux de CD4+ (orange, trait pointillé, échelle linéaire à droite) au cours de l'infection VIH non traitée. Les zones colorées délimitent les 3 phases cliniques. Notez le « point de consigne » viral (plateau phase 2) et l'inversion des courbes (ciseaux).}
\end{popfigure}

\courssub{Phase 1 : Infection Primaire (Semaines 1 à 12) — Le « raz-de-marée » viral}

C'est la phase d'invasion massive. Après la transmission (sexuelle, sanguine, mère-enfant), le virus franchit la muqueuse, atteint les ganglions lymphatiques et se réplique de façon exponentielle.

\begin{definition}[Séroconversion et Fenêtre sérologique]
La \textbf{séroconversion} est le passage du statut séronégatif (anticorps anti-VIH indétectables) au statut séropositif (anticorps détectables par test rapide ou ELISA). Elle survient en moyenne \textbf{3 à 6 semaines} après l'infection.
La \textbf{fenêtre sérologique} est la période entre l'infection et la séroconversion : l'antigène viral (Ag p24) et l'ARN viral (PCR) sont détectables, mais les anticorps anti-VIH ne le sont pas encore. \textbf{Diagnostic précoce : PCR ARN VIH ou test combiné Ag p24/Ac (4ème génération).}
\end{definition}

\begin{itemize}
    \item \textbf{Virémie :} Pic extrême \textbf{$> 10^6$ copies ARN/mL} (très contagieux).
    \item \textbf{CD4+ :} Chute brutale et transitoire (perte de 30 à 50\% du taux initial).
    \item \textbf{Clinique :} Syndrome pseudo-grippal ou mononucléosique chez \textbf{50 à 90\%} des cas (fièvre, adénopathies, éruption cutanée, maux de gorge, douleurs musculaires). Souvent méconnu ou confondu avec un paludisme ou une grippe au Cameroun.
    \item \textbf{Réponse immunitaire :} Montée des lymphocytes T CD8+ cytotoxiques (CTL) qui détruisent les cellules infectées. Cette réponse fait chuter la virémie de 2 à 3 log (facteur 100 à 1000).
\end{itemize}

\begin{propriete}[Le Point de Consigne Viral (Set Point)]
Vers \textbf{6 mois}, s'établit un équilibre précaire : le \textbf{point de consigne viral}. C'est le niveau de charge virale stabilisé après la réponse immune initiale. C'est un \textbf{marqueur pronostique} majeur :
\begin{itemize}
    \item Set Point bas ($< 10^3$ copies/mL) $\rightarrow$ Progression lente vers le SIDA.
    \item Set Point haut ($> 10^5$ copies/mL) $\rightarrow$ Progression rapide (SIDA en 3-5 ans).
\end{itemize}
\end{propriete}

\courssub{Phase 2 : Phase Chronique / Asymptomatique (Années 1 à 10+) — La guerre d'usure}

Le patient se sent bien, ne présente pas de signes cliniques majeurs, mais la réplication virale est continue (des milliards de nouveaux virions produits \textbf{par jour}).

\begin{attention}[Piège majeur : « Asymptomatique $\neq$ Inactif »]
La phase chronique n'est \textbf{PAS} une latence virale.
\begin{enumerate}
    \item \textbf{Réplication active continue :} Le virus mute, des variants peuvent apparaître.
    \item \textbf{Destruction des organes lymphoïdes :} Les ganglions, d'abord gonflés (adénopathies), voient leur architecture se détruire (fibrose). Le réseau de soutien nécessaire à la maturation des lymphocytes est abîmé.
    \item \textbf{Inflammation chronique :} Le système immunitaire est constamment activé, ce qui l'épuise et favorise d'autres maladies (cardiovasculaires, rénales).
    \item \textbf{Réservoirs viraux :} Le virus se cache dans certaines cellules au repos (mémoire), invisibles aux traitements et au système immunitaire. C'est la barrière à la guérison complète.
\end{enumerate}
\end{attention}

\begin{popfigure}
\begin{tikzpicture}[scale=0.85]
% Ganglion Normal
\begin{scope}[xshift=-4cm]
\node[align=center, font=\bfseries\small] at (0, 3.2) {Ganglion Normal};
\draw[thick, fill=popgreenL!20] (0,0) ellipse (1.6 and 2.2); % Capsule
\draw[thick, fill=popblueL!30] (0, 1.3) ellipse (1.1 and 0.5) node[midway, font=\tiny, align=center]{Zone B\\Follicules};
\draw[thick, fill=popblueL!10] (0, 1.3) ellipse (0.5 and 0.25) node[midway, font=\tiny, align=center]{Centre\\Germinatif};
\draw[thick, fill=poporangeL!30] (-1.1, -0.5) rectangle (1.1, 0.8) node[midway, font=\tiny, align=center]{Zone T (Paracortex)\\CD4+ naïfs};
\draw[->, thick, popdark] (-1.9, 0.2) -- (-1.2, 0.2) node[left, font=\tiny] {Lymphe afférente};
\draw[->, thick, popdark] (1.2, -1.3) -- (1.9, -1.3) node[right, font=\tiny] {Lymphe efférente};
\end{scope}
% Ganglion Infecté VIH (Phase Chronique -> Fibrose)
\begin{scope}[xshift=4cm]
\node[align=center, font=\bfseries\small] at (0, 3.2) {Ganglion Infecté VIH (Phase Chronique)};
\draw[thick, fill=popred!10] (0,0) ellipse (1.6 and 2.2); % Capsule épaissie
\draw[thick, fill=poppurpleL!40, pattern=north east lines, pattern color=poppurple] (0, 1.3) ellipse (1.1 and 0.5) node[midway, font=\tiny, text=popdark, align=center]{Fibrose /\\Destruction};
\draw[thick, fill=poporangeL!10] (-1.1, -0.5) rectangle (1.1, 0.8) node[midway, font=\tiny, text=popdark] {Zone T atrophiée};
% Virions piégés
\foreach \i in {-0.9, -0.4, 0.1, 0.6} {
    \fill[popred] (\i, 1.4) circle (0.04);
    \fill[popred] (\i+0.1, 1.3) circle (0.04);
}
\node[font=\tiny, popred] at (0, 2.1) {Virions piégés};
\draw[->, thick, popdark] (-1.9, 0.2) -- (-1.2, 0.2) node[left, font=\tiny] {Afférente};
\draw[->, thick, popdark] (1.2, -1.3) -- (1.9, -1.3) node[right, font=\tiny] {Efférente};
\end{scope}
% Légende
\node[font=\small, anchor=north, align=center] at (0, -2.8){
\begin{tabular}{ll}
\textcolor{popblueL!80!black}{\rule{3mm}{3mm}} Zone B (Lymphocytes B) & \textcolor{poporangeL!80!black}{\rule{3mm}{3mm}} Zone T (Lymphocytes T CD4+) \\
\textcolor{poppurple}{\rule{3mm}{3mm}} Fibrose (Destruction) & \textcolor{popred}{\rule{2mm}{2mm}} Virions VIH
\end{tabular}
};
\end{tikzpicture}
\legende{Évolution de l'architecture ganglionnaire : à gauche, structure normale ; à droite, destruction des centres germinatifs (fibrose), atrophie de la zone T et piégeage des virions.}
\end{popfigure}

\courssub{Phase 3 : SIDA (Stade Terminal) — L'effondrement}

Sans traitement, la phase chronique dure en médiane 8 à 10 ans. Le seuil critique est franchi quand le capital immunitaire ne permet plus de se défendre contre des microbes habituellement inoffensifs.

\begin{definition}[Critères diagnostiques du SIDA (OMS Stade 4)]
Le SIDA est défini par \textbf{UN} des critères suivants chez un séropositif VIH :
\begin{enumerate}
    \item \textbf{Immunologique :} Taux de CD4+ \textbf{$< 200$ cellules/$\mu$L} (ou $< 14\%$ des lymphocytes totaux).
    \item \textbf{Clinique :} Survenue d'une \textbf{infection opportuniste (IO) définissante} ou d'un \textbf{cancer définissant}, indépendamment du taux de CD4+.
\end{enumerate}
Sans traitement, le décès survient en \textbf{1 à 2 ans}.
\end{definition}

\begin{center}
\begin{tabularx}{\linewidth}{|l|X|}\hline
\rowcolor{popblueL} \textbf{Stade OMS} & \textbf{Manifestations cliniques majeures (Contexte Cameroun)} \\ \hline
\textbf{Stade 1} & Asymptomatique ou Adénopathies généralisées persistantes (PGL). \\ \hline
\textbf{Stade 2} & Perte de poids $< 10\%$, zona, infections cutanées, angines répétées. \\ \hline
\textbf{Stade 3} & Perte de poids $> 10\%$, fièvre/diarrhée chroniques, tuberculose pulmonaire, candidose buccale, pneumonies bactériennes sévères. \\ \hline
\textbf{Stade 4 (SIDA)} & \textbf{Infections Opportunistes :} Pneumocystose, Toxoplasmose cérébrale, Cryptococcose méningée, Candidose \textbf{œsophagienne}, Tuberculose disséminée, Cytomégalovirus. \\
& \textbf{Cancers :} Sarcome de Kaposi, Lymphomes, Cancer du col utérin invasif. \\ \hline
\end{tabularx}
\end{center}
\legende{Classification OMS des 4 stades. Le passage au Stade 4 définit le SIDA. Au Cameroun, la Tuberculose et la Pneumocystose sont fréquentes.}

\courssub{Marqueurs biologiques de suivi : La boussole du clinicien}

Deux examens de référence (disponibles dans les Centres de Prise en Charge - CPEC - au Cameroun) guident la prise en charge.

\begin{propriete}[Marqueurs de suivi standard]
\begin{enumerate}
    \item \textbf{Charge Virale (CV) :} Quantification de l'ARN VIH plasmatique par \textbf{PCR temps réel}. Unité : \textbf{copies/mL} (ou Log$_{10}$ copies/mL).
        \begin{itemize}
            \item Reflet direct de l'activité de réplication du virus.
            \item Objectif sous ARV : \textbf{Indétectable ($< 50$ copies/mL)}.
            \item Échec virologique : CV $> 1000$ copies/mL (confirmé) sous traitement bien pris.
        \end{itemize}
    \item \textbf{Numération des lymphocytes T CD4+ :} \textbf{Cytométrie de flux}. Unité : \textbf{cellules/$\mu$L} (et pourcentage).
        \begin{itemize}
            \item Reflet du \textbf{capital immunitaire} (risque d'infections opportunistes).
            \item Normal : $500$--$1500$ cellules/$\mu$L.
            \item Ratio CD4/CD8 : Normal $> 1$. Inversé ($< 1$) dès l'infection.
        \end{itemize}
\end{enumerate}
\end{propriete}

\begin{methode}[Interpréter un bilan pour classer le stade et décider l'urgence]
\textbf{Scénario :} Patient 32 ans, nouveau diagnostiqué VIH+. Résultats : CV = 250 000 copies/mL ; CD4+ = 180 cellules/$\mu$L.

\textbf{Étapes :}
\begin{enumerate}
    \item \textbf{Classer le stade immunologique :} CD4+ $< 200$ $\rightarrow$ \textbf{Stade 4 (SIDA) immunologique}.
    \item \textbf{Évaluer le risque immédiat :} CD4+ $< 200$ $\rightarrow$ Risque élevé de \textbf{Pneumocystose} et \textbf{Toxoplasmose cérébrale}.
    \item \textbf{Prescrire les prophylaxies primaires d'URGENCE :} \textbf{Cotrimoxazole (Bactrim\textsuperscript{\textregistered}) 960 mg/j} (prévention Pneumocystose + Toxoplasmose + bactéries).
    \item \textbf{Rechercher une infection opportuniste en cours :} Radio thorax (TB), examen neurologique, etc.
    \item \textbf{Initiation ARV :} \textbf{Dès que possible (idéalement J0 à J7)}.
    \item \textbf{Suivi :} CV à 3 mois (cible $< 50$), CD4+ à 6 mois (cible remontée $> 50-100$ cellules/$\mu$L/an).
\end{enumerate}
\end{methode}

\begin{exemplebox}[Seuils CD4+ et Prophylaxies Primaires (Guide National Cameroun)]
\begin{center}
\begin{tabular}{|c|c|l|}\hline
\rowcolor{popblueL} \textbf{Seuil CD4+} & \textbf{Médicament} & \textbf{Infection prévenue} \\ \hline
$< 200$ / $\mu$L & \textbf{Cotrimoxazole} 960 mg/j & \emph{Pneumocystis}, \emph{Toxoplasma}, bactéries \\ \hline
$< 100$ / $\mu$L & \textbf{Azithromycine} 1200 mg/sem & \emph{Mycobacterium avium} complexe (MAC) \\ \hline
$< 50$ / $\mu$L & \textbf{Fluconazole} 200 mg/j & \emph{Cryptococcus} (Méningite) \\ \hline
\end{tabular}
\end{center}
\legende{Prophylaxies primaires selon le niveau d'immunodépression. Le Cotrimoxazole est la pierre angulaire.}
\end{exemplebox}

\courssub{Physiopathologie intégrée : Pourquoi le système s'effondre-t-il ?}

La déplétion des CD4+ ne vient pas seulement de l'infection directe (seulement $\sim 0,01\%$ à $0,1\%$ des CD4+ circulants sont infectés à un instant $t$). Les mécanismes \textbf{indirects} (effet « spectateur » ou \emph{bystander effect}) sont majoritaires :

\begin{enumerate}
    \item \textbf{Apoptose des cellules non infectées :} Protéines virales libres (gp120, Tat, Nef) se fixent sur des cellules saines $\rightarrow$ déclenchent leur mort programmée (apoptose) ou fusion cellulaire (syncytia). L'activation chronique épuise les cellules (mort par activation).
    \item \textbf{Destruction du nichée lymphoïde :} Fibrose ganglionnaire $\rightarrow$ perte du microenvironnement (cellules de soutien, interleukine-7) nécessaire à la survie et la maturation des nouveaux lymphocytes T. Le thymus s'atrophie prématurément.
    \item \textbf{Perte du répertoire immunitaire :} Disparition des clones spécifiques d'antigènes et incapacité à en générer de nouveaux $\rightarrow$ « trous » dans la défense $\rightarrow$ infections opportunistes.
    \item \textbf{Dysfonction des cellules résiduelles :} Les CD4+ et CD8+ restants sont « épuisés » (perdent leur capacité à se multiplier et sécréter plusieurs cytokines à la fois).
\end{enumerate}

\begin{savaistu}[Les « Contrôleurs d'élite » : L'exception qui guide la recherche]
Moins de \textbf{0,5\%} des infectés (dits « contrôleurs d'élite ») maintiennent spontanément une CV indétectable et des CD4+ normaux \textbf{sans traitement} pendant des décennies.
\textbf{Leçons tirées :}
\begin{itemize}
    \item Ils ont une réponse \textbf{CD8+ exceptionnellement efficace} (tuent bien les cellules infectées).
    \item Ils possèdent souvent des \textbf{allèles HLA protecteurs} (ex: HLA-B*57, HLA-B*27) qui présentent des morceaux de virus très conservés, forçant le virus à muter en s'affaiblissant.
\end{itemize}
Étudier ces patients guide la recherche vaccinale (cibler les parties conservées du virus) et les stratégies de rémission sans traitement à vie.
\legende{Les contrôleurs d'élite prouvent qu'un contrôle immunitaire durable est possible.}
\end{savaistu}

\courssub{Prévention, Dépistage et Traitement Antirétroviral (ARV) au Cameroun}

Le programme 3ème insiste sur la prévention et le principe de l'action des ARV.

\begin{propriete}[Stratégies de prévention du VIH (Approche combinée)]
\begin{itemize}
    \item \textbf{Comportementales :} Abstinence, Fidélité mutuelle, \textbf{Préservatif} (masculin/féminin) à chaque rapport à risque.
    \item \textbf{Biomédicales :} Dépistage volontaire (connaître son statut), \textbf{PTME} (Prévention Transmission Mère-Enfant : ARV mère + AZT nourrisson + allaitement sécurisé), \textbf{PEP} (Prophylaxie Post-Exposition : ARV d'urgence dans les 72h après accident), \textbf{PrEP} (Prophylaxie Pré-Exposition : ARV préventifs pour personnes à haut risque), Circoncision médicale volontaire (réduit risque acquisition homme $\sim 60\%$).
    \item \textbf{Structurelles :} Lutte contre la stigmatisation, éducation sexuelle complète, réduction des inégalités de genre.
\end{itemize}
\end{propriete}

\begin{propriete}[Traitement Antirétroviral (ARV) : Principe et Objectifs]
\begin{itemize}
    \item \textbf{Principe :} Association de \textbf{3 molécules} (Trithérapie) ciblant différentes étapes du cycle viral (Entrée, Transcription inverse, Intégration, Protéase). Empêche la réplication $\rightarrow$ charge virale chute $\rightarrow$ système immunitaire se reconstitue.
    \item \textbf{Indication :} \textbf{Tout séropositif, quel que soit le stade ou le taux de CD4+} (Recommandation OMS « Treat All » depuis 2015, confirmée par l'essai START).
    \item \textbf{Objectifs :}
        \begin{enumerate}
            \item \textbf{Virologique :} Charge virale **indétectable ($< 50$ copies/mL)** $\rightarrow$ Arrêt transmission (TasP : \emph{Treatment as Prevention} / U=U : \emph{Undetectable = Untransmittable}).
            \item \textbf{Immunologique :} Remontée durable des CD4+ ($> 500$ cellules/$\mu$L).
            \item \textbf{Clinique :} Suppression des IO, vie normale, espérance de vie quasi normale.
        \end{enumerate}
    \item \textbf{Observance :} \textbf{Condition sine qua non} de l'efficacité. Oubli de doses $\rightarrow$ sélection mutants résistants $\rightarrow$ échec thérapeutique.
\end{itemize}
\end{propriete}

\courssub{Exercices d'intégration type BEPC}

\begin{exoresolu}[Analyse de l'évolution d'un patient non traité]
\textbf{Énoncé :} Le graphique ci-dessous représente l'évolution de la charge virale (CV) et du taux de CD4+ chez un patient infecté par le VIH, non traité, suivi pendant 10 ans.
\begin{center}
\begin{tikzpicture}[scale=0.7]
\begin{axis}[width=0.9\linewidth, height=0.5\linewidth, xlabel={Années}, ylabel={Valeurs}, xmin=0, xmax=10, ymin=0, ymax=7, grid=major, legend pos=north east]
\addplot[popblue, thick] coordinates {(0,6.5) (0.2,6.8) (0.5,4.2) (2,4.0) (5,4.3) (8,4.5) (9,5.5) (10,6.5)};
\addplot[poporange, thick, dashed] coordinates {(0,6.0) (0.2,4.0) (2,4.5) (5,3.5) (8,2.0) (9,1.0) (10,0.5)};
\addlegendentry{CV (Log$_{10}$)};
\addlegendentry{CD4+ ($\times 100$ cellules/$\mu$L)};
\end{axis}
\end{tikzpicture}
\end{center}
\textbf{Questions :}
\begin{enumerate}
    \item Identifier les phases A, B, C sur le graphique (donner les années approximatives).
    \item Expliquer la forme de la courbe de charge virale en phase A (pic puis chute).
    \item Pourquoi la courbe de CD4+ chute-t-elle brutalement en phase A alors que peu de cellules sont infectées ?
    \item À partir de quelle année le patient entre-t-il au stade SIDA (critère immunologique) ? Justifier.
    \item Citer deux infections opportunistes à risque élevé pour ce patient à l'année 9.
\end{enumerate}

\textbf{Solution détaillée :}
\begin{enumerate}
    \item \textbf{Phase A (Infection Primaire) :} Années 0 à 0,5. Pic CV $\to$ Chute. Séroconversion vers 0,1-0,2 an.
    \textbf{Phase B (Phase Chronique) :} Années 0,5 à $\sim$ 8,5. Plateau CV (Set point $\sim 10^4$), déclin lent CD4+.
    \textbf{Phase C (SIDA) :} Années $\sim$ 8,5 à 10. Remontée CV, effondrement CD4+.
    \item \textbf{Phase A :} Infection massive $\rightarrow$ réplication exponentielle (Pic). Montée de la réponse immune (CD8+, Ac) $\rightarrow$ lyse cellules infectées + blocage nouvelle infection $\rightarrow$ Chute virémie. Établissement du Set Point.
    \item \textbf{Mécanismes indirects (bystander) :} Protéines virales (gp120, Tat, Nef) induisent l'apoptose des CD4+ non infectés. Activation immune chronique $\rightarrow$ mort par activation. Destruction précoce de l'architecture ganglionnaire (perte niches de survie).
    \item \textbf{Année 8,5 - 9.} Le seuil CD4+ $< 200$ cellules/$\mu$L (soit $< 2$ sur l'échelle $\times 100$ du graphique) est franchi vers l'année 8,5-9. Critère immunologique du SIDA rempli.
    \item \textbf{Pneumocystose (\emph{Pneumocystis jirovecii})} et \textbf{Toxoplasmose cérébrale} (risque majeur CD4 $< 200$). Aussi : Candidose œsophagienne, Tuberculose disséminée, Cryptococcose méningée.
\end{enumerate}
\end{exoresolu}

\begin{exercice}[Cas clinique : Prise en charge au CPEC]
\textbf{Énoncé :} M. K., 35 ans, se présente au CPEC de Ngaoundéré. Il a été dépisté VIH+ il y a 2 semaines. Il tousse depuis 3 semaines, a de la fièvre le soir et a perdu 5 kg en 2 mois. Bilan initial :
\begin{itemize}
    \item Charge Virale : 450 000 copies/mL.
    \item CD4+ : 150 cellules/$\mu$L.
    \item Radiographie thorax : Infiltrats interstitiels bilatéraux diffuses.
    \item Hémoglobine : 11 g/dL. Fonction rénale et hépatique normales.
\end{itemize}
\textbf{Questions :}
\begin{enumerate}
    \item Classer le stade clinique et immunologique de M. K. (OMS).
    \item Quelle infection opportuniste suspecter fortement devant ce tableau (toux, fièvre, infiltrats interstitiels, CD4 $< 200$) ?
    \item Quel traitement d'urgence prescrire \textbf{avant même les ARV} pour cette infection suspectée ?
    \item Citer le schéma ARV de 1ère ligne recommandé au Cameroun (3 molécules) et son objectif principal.
    \item Expliquer pourquoi l'observance stricte (prise quotidienne sans oubli) est vitale pour M. K.
\end{enumerate}
\end{exercice}

\begin{corrige}[Exercice n°1 : Cas clinique M. K.]
\begin{enumerate}
    \item \textbf{Stade Clinique :} Stade 4 (SIDA) -- Perte de poids $> 10\%$ + fièvre chronique + suspicion pneumocystose (IO définissante). \textbf{Stade Immunologique :} CD4+ = 150 ($< 200$) $\rightarrow$ Stade 4 (SIDA) immunologique.
    \item \textbf{Pneumocystose (\emph{Pneumocystis jirovecii})}. C'est la pneumonie opportuniste la plus fréquente au stade SIDA (CD4 $< 200$) avec ce tableau radiologique (infiltrats interstitiels diffuses « en verre dépoli »).
    \item \textbf{Cotrimoxazole (Bactrim\textsuperscript{\textregistered})} : Dose curative (souvent 960 mg $\times$ 3/j ou 1920 mg $\times$ 2/j pendant 21 jours) par voie orale ou IV selon gravité. C'est le traitement de 1ère ligne de la pneumocystose. \textit{Note : Le Cotrimoxazole en prophylaxie primaire (960 mg/j) n'aurait pas dû être omis vu le CD4 $< 200$}.
    \item \textbf{Schéma 1ère ligne Cameroun :} \textbf{TDF + 3TC (ou FTC) + DTG} (Ténofovir + Lamivudine/Emtricitabine + Dolutégravir). \textbf{Objectif principal :} Obtenir une **charge virale indétectable ($< 50$ copies/mL)** durablement pour permettre la remontée des CD4+ et la guérison clinique.
    \item L'observance stricte empêche la réplication virale résiduelle. Tout oubli laisse le virus se multiplier $\rightarrow$ mutations $\rightarrow$ \textbf{résistance aux ARV} $\rightarrow$ échec du traitement (charge virale remonte, CD4+ rechutent, IO réapparaissent). Les molécules de 2ème ligne sont plus chères, plus toxiques, moins disponibles.
\end{enumerate}
\end{corrige}

\begin{aretenir}[Bilan synthèse : Dynamique de l'infection VIH]
\begin{itemize}
    \item \textbf{3 Phases :} Primaire (Pic CV, chute CD4, séroconversion) $\rightarrow$ Chronique (Set Point CV, déclin lent CD4 $\sim$ 70/an, destruction ganglionnaire) $\rightarrow$ SIDA (CD4 $< 200$, IO/Cancers, décès 1-2 ans sans ARV).
    \item \textbf{Marqueurs :} CV (PCR ARN) = Réplication virale ; CD4+ (Cytométrie) = Capital immunitaire ; Ratio CD4/CD8 inversé.
    \item \textbf{Physiopathologie :} Déplétion CD4+ mostly indirecte (apoptose bystander, destruction architecture lymphoïde, épuisement, perte thymique).
    \item \textbf{Prise en charge :} ARV précoces pour tous $\rightarrow$ CV indétectable, remontée CD4, vie quasi-normale, PTME efficace, Prévention (TasP/U=U).
    \item \textbf{Prophylaxies :} CD4 $< 200$ (Cotrimoxazole curatif/préventif), $< 100$ (Azithro), $< 50$ (Fluconazole).
    \item \textbf{Prévention :} ABC (Abstinence, Fidélité, Préservatif), Dépistage, PTME, PEP, PrEP, Circoncision.
\end{itemize}
\end{aretenir}

\coursec{Prévention, Dépistage et Traitement Antirétroviral (ARV) au Cameroun}

Après avoir analysé la dynamique de l'infection à VIH, depuis l'entrée du virus jusqu'à l'effondrement immunitaire caractérisant le Sida (section 3), nous comprenons qu'une intervention est possible à plusieurs niveaux : empêcher le virus d'entrer (prévention), le détecter le plus tôt possible (dépistage) et bloquer sa réplication s'il est déjà présent (traitement). Au Cameroun, la riposte nationale s'appuie sur une stratégie \textbf{combinée} intégrant ces trois piliers, gratuite dans le secteur public grâce à l'appui de l'État et des partenaires (Fonds mondial, PEPFAR, OMS).

\courssub{Dépistage du VIH : Principe des Tests Rapides et Algorithme National}

Le dépistage est la \textbf{porte d'entrée} de la prise en charge. Connaître son statut sérologique permet d'accéder précocement au traitement ou de renforcer la prévention. Au Cameroun, on utilise des \textbf{Tests Rapides de Diagnostic (TRD) de 4ème génération} sur sang total (prélèvement au doigt), détectant simultanément les \textbf{anticorps anti-VIH-1/2} (réponse immunitaire) et l'\textbf{antigène p24} (protéine virale précoce).

\begin{experience}[Le Test Rapide VIH (TRD 4ème génération) : Principe et Lecture]
\textbf{Matériel :} Cassette de test (détection Ag/Ab), lancette stérile, capillaire, tampon de chase, chronomètre, gants.
\textbf{Protocole :}
\begin{enumerate}
    \item Désinfecter le doigt, piquer avec la lancette.
    \item Recueillir une goutte de sang total, déposer dans le puits \og Sample \fg{} de la cassette.
    \item Ajouter 2 à 3 gouttes de tampon (chase buffer).
    \item Lire le résultat \textbf{entre 15 et 20 minutes} (lecture avant 15 min = faux négatif possible ; après 20 min = bandes fantômes).
\end{enumerate}
\textbf{Observations (Bandettes) :} La cassette possède une zone de lecture avec des zones de capture : \textbf{C} (Contrôle), \textbf{T1} (VIH-1), \textbf{T2} (VIH-2), \textbf{Ag} (Antigène p24).
\begin{itemize}
    \item \textbf{Bande C seule visible} $\rightarrow$ \textbf{Négatif} (Test valide, pas d'Ag/Ac détectés).
    \item \textbf{Bande C + Bande T1 et/ou T2 et/ou Ag} $\rightarrow$ \textbf{Réactif (Positif)}.
    \item \textbf{Pas de bande C} $\rightarrow$ \textbf{Invalide} (Test raté, à refaire).
\end{itemize}
\textbf{Interprétation :} La détection de l'Ag p24 réduit la \textbf{fenêtre sérologique} (période entre infection et positivité du test) à \textbf{2 à 3 semaines}. Un test négatif récent ne garantit pas l'absence d'infection si l'exposition date de moins de 3 semaines.
\legende{Schéma de lecture d'un TRD 4ème génération : Bandes C (contrôle), T1 (VIH-1), T2 (VIH-2), Ag (Antigène p24).}
\end{experience}

\begin{popfigure}
\begin{tikzpicture}[
    node distance=1.5cm and 2cm,
    box/.style={draw, rounded corners, minimum width=3.5cm, minimum height=1.2cm, text centered, font=\small},
    decision/.style={diamond, draw, aspect=2, minimum width=3cm, text centered, font=\small},
    arrow/.style={-Stealth, thick},
    poslabel/.style={font=\scriptsize, text width=2.5cm, align=center}
]
% Nodes
\node[box, fill=popblueL, align=center] (start){Dépistage : Test de Dépistage\\(Haute Sensibilité)};
\node[decision, below of=start, fill=poporangeL] (testA) {Test Réactif ?};
\node[box, right of=testA, xshift=4cm, fill=popgreenL, align=center] (neg){NÉGATIF\\(Non infecté\\ou Fenêtre sérologique)\\$\rightarrow$ Conseil prévention\\$\rightarrow$ Re-test à 6 sem};
\node[box, below of=testA, fill=poppinkL, align=center] (testB){Test de Confirmation\\(Haute Spécificité)};
\node[decision, below of=testB, fill=poporangeL] (testBdec) {Test Confirmation Positif ?};
\node[box, right of=testBdec, xshift=4cm, fill=popgold!18, align=center] (discord){DISCORDANT\\$\rightarrow$ Nouveau prélèvement\\à J+14 ou Test PCR};
\node[box, below of=testBdec, fill=poppurpleL, align=center] (pos){POSITIF CONFIRMÉ\\(VIH-1 / VIH-2 / Dual)\\$\rightarrow$ Orientation prise en charge\\$\rightarrow$ Bilan CD4 / CV / Clinique};
% Arrows
\draw[arrow] (start) -- (testA);
\draw[arrow] (testA) -- node[poslabel, left] {NON} (neg);
\draw[arrow] (testA) -- node[poslabel, left] {OUI} (testB);
\draw[arrow] (testB) -- (testBdec);
\draw[arrow] (testBdec) -- node[poslabel, right] {NON} (discord);
\draw[arrow] (testBdec) -- node[poslabel, left] {OUI} (pos);
\end{tikzpicture}
\legende{Algorithme national de dépistage VIH au Cameroun (Stratégie 2 tests séquentiels). Test 1 = Sensibilité max (peu de faux négatifs). Test 2 = Spécificité max (confirmation). En cas de discordance : nouveau test ou PCR ARN VIH.}
\end{popfigure}

\begin{methode}[Lire et interpréter un résultat de TRD et conduite à tenir]
\begin{enumerate}
    \item \textbf{Vérifier la validité :} La bande \textbf{C (Contrôle)} \textbf{doit} apparaître. Si absente $\rightarrow$ \textbf{Invalide} $\rightarrow$ Refaire le test (nouvelle cassette).
    \item \textbf{Lire les bandes de test :} Observer les bandes \textbf{T1}, \textbf{T2}, \textbf{Ag}. Toute couleur (même faible) = Réactif.
    \item \textbf{Interpréter selon l'algorithme :}
    \begin{itemize}
        \item \textbf{C seul} = Négatif. Rassurer, conseiller prévention, proposer re-test à 6 semaines si exposition récente (< 3 sem).
        \item \textbf{C + (T1 ou T2 ou Ag)} = Réactif. \textbf{Ne pas annoncer "Sida" ni "Positif définitif"}. Dire : "Le test est réactif, nous devons confirmer par un deuxième test".
        \item Enchaîner immédiatement \textbf{Test de Confirmation}. Si Confirmation négatif $\rightarrow$ Discordance $\rightarrow$ Nouveau prélèvement à 2 semaines ou PCR ARN VIH.
        \item Si Confirmation positif $\rightarrow$ \textbf{Positif confirmé}. Orientation vers le service de prise en charge pour bilan pré-thérapeutique (CD4, Charge Virale, Hépatites, Fonction rénale, Grossesse).
    \end{itemize}
    \item \textbf{Conseil post-test :} Annonce du résultat confirmé, soutien psychologique, orientation vers la prise en charge.
\end{enumerate}
\end{methode}

\begin{attention}[Pièges fréquents au dépistage]
\begin{itemize}
    \item \textbf{Confondre "Réactif" et "Positif confirmé" :} Un seul test positif \textbf{ne suffit pas} pour diagnostiquer l'infection. L'annonce prématurée cause un traumatisme psychologique majeur.
    \item \textbf{Négliger la fenêtre sérologique :} Un test négatif chez un exposé récent (< 3 semaines) ne rassure pas. Reprogrammer un TRD à 6 semaines ou prescrire une \textbf{PCR ARN VIH} (détection directe du virus, fenêtre ~10 jours).
    \item \textbf{Oublier le typage VIH-1 / VIH-2 :} Au Cameroun, le VIH-2 circule. Il est \textbf{résistant nature} aux INNRTI (Névirapine, Éfavirenz). La confirmation doit discriminer VIH-1 et VIH-2 pour adapter le traitement.
\end{itemize}
\end{attention}

\courssub{Prévention Combinée : Une boîte à outils adaptée au risque}

La prévention ne repose pas sur une seule méthode. Le Cameroun promeut la \textbf{prévention combinée} : association de mesures biomédicales, comportementales et structurelles.

\begin{definition}[Prévention Combinée]
Ensemble cohérent d'interventions de prévention du VIH (préservatifs, dépistage, traitement comme prévention, PrEP, PEP, PMTCT, circoncision) ciblées selon le niveau de risque, visant à réduire l'incidence (nouveaux cas) à zéro.
\end{definition}

\begin{center}
\begin{tabularx}{\linewidth}{|l|X|X|}\hline
\rowcolor{popblueL} \textbf{Stratégie} & \textbf{Principe / Cible} & \textbf{Efficacité / Contexte Cameroun} \\\hline
\textbf{Préservatifs (M/F)} & Barrière mécanique (latex/nitrile). Homme (externe) / Femme (interne). & \textbf{> 95\%} si usage correct/constant. Gratuits en formations sanitaires. Seule méthode protégeant aussi des IST et grossesses non désirées. \\\hline
\textbf{TASP (Traitement comme Prévention)} & \cle{I = I} : Indétectable = Intransmissible. PVVIH sous ARV efficace (Charge Virale < 200 copies/mL stable > 6 mois). & \textbf{0\% transmission sexuelle} (preuves scientifiques solides). Pilier de l'élimination. \\\hline
\textbf{PrEP (Prophylaxie Pré-Exposition)} & Médicament ARV (TDF/FTC) pris \textbf{quotidiennement} par personne séronégative à risque élevé. & \textbf{> 99\%} si observance parfaite. Disponible au Cameroun pour populations clés (HSH, TS, partenaires de PVVIH, travailleurs du sexe). Test VIH négatif requis avant initiation. \\\hline
\textbf{PEP (Prophylaxie Post-Exposition)} & Trithérapie d'urgence pendant \textbf{28 jours}. \textbf{< 72h} post-exposition (viol, accident sanguin, rupture préservatif). & Efficacité \textbf{> 80\%} si début précoce (< 24h idéale). Kit PEP disponible en urgences. Tests VIH à J0, J28, 3 mois. \\\hline
\textbf{PMTCT (Prévention Mère-Enfant)} & \textbf{Option B+ :} ARV à vie pour \textbf{toutes} femmes enceintes/allaitantes VIH+. Sirop AZT nourrisson 6 sem. Allaitement maternel exclusif 6 mois + ARV mère. & Taux transmission Mère-Enfant \textbf{< 5\%} (objectif < 2\%). Suivi "Mère-Enfant" intégré. \\\hline
\textbf{Circoncision Médicale (Hommes)} & Ablation prépuce (réduit cellules cibles, micro-ulcérations). & Réduction acquisition VIH hétérosexuel \textbf{~ 60\%} (Homme). Programme actif régions Adamaoua, Nord, Extrême-Nord. \\\hline
\end{tabularx}
\end{center}

\begin{propriete}[I = I : Indétectable = Intransmissible (Preuve scientifique)]
Une personne vivant avec le VIH (PVVIH) sous traitement antirétroviral efficace, maintenant une \textbf{charge virale (CV) plasmatique < 200 copies/mL de façon durable (> 6 mois)}, \textbf{ne transmet pas le VIH par voie sexuelle}, même sans préservatif.
\textbf{Preuves majeures :} Études HPTN 052 (2011) et PARTNER 1 \& 2 (2016-2019) sur des milliers de couples sérodifférents et > 100 000 actes sexuels sans préservatif : \textbf{Zéro transmission liée au partenaire traité} quand CV indétectable.
\emph{Condition sine qua non :} Observance stricte (CV durablement indétectable). "I=I" est un message majeur pour réduire la stigmatisation et favoriser l'adhésion au traitement.
\end{propriete}

\begin{savaistu}[PrEP au Cameroun : De la recherche à l'accès]
Le Cameroun a participé à des essais cliniques majeurs sur la PrEP (IPERGAY, ANRS 12174). Depuis 2021, la PrEP quotidienne (TDF/FTC) est intégrée dans le paquet de prévention pour les \textbf{populations clés}. Elle est délivrée gratuitement dans les centres agréés (CEPEC, ONG partenaires) après conseil, test VIH négatif, dépistage Hépatite B (vaccination si négatif) et vérification fonction rénale. L'observance est soutenue par des \textbf{pairs éducateurs} et des rappels (SMS).
\end{savaistu}

\courssub{Traitement Antirétroviral (ARV) : La Trithérapie Standard}

Le traitement du VIH ne guérit pas (éradication impossible à cause des \textbf{réservoirs viraux} : lymphocytes T CD4+ mémoire au repos, macrophages, SNC), mais il \textbf{contrôle la réplication} de façon durable.

\begin{definition}[Trithérapie (ARV - Antirétroviraux)]
Association de \textbf{3 molécules} actives appartenant à \textbf{au moins 2 classes pharmacologiques différentes}, ciblant \textbf{étapes distinctes} du cycle de réplication (ex: Entrée, Transcription inverse, Intégration, Assemblage/Maturation).
\textbf{Objectif :} Synergie antivirale maximale + \textbf{barrière génétique élevée} (nécessité de multiples mutations simultanées pour échapper) $\rightarrow$ Prévention de l'émergence de résistances.
\end{definition}

\begin{popfigure}
\begin{tikzpicture}[
    node distance=0.8cm and 0.5cm,
    class/.style={draw, rounded corners, fill=popblueL, minimum width=3cm, minimum height=1cm, text centered, font=\small\bfseries},
    mol/.style={draw, rounded corners, fill=popcream, minimum width=3cm, minimum height=0.9cm, text centered, font=\small},
    target/.style={draw, rounded corners, fill=popgreenL, minimum width=3.5cm, minimum height=0.8cm, text centered, font=\small},
    arrow/.style={-Stealth, thick, popdark}
]
% Classes
\node[class, align=center] (c1){INRTI\\(Inhibiteurs Nucléosidiques\\de la Transcriptase Inverse)};
\node[class, right of=c1, xshift=1.5cm, align=center] (c2){INNRTI\\(Inhibiteurs Non-Nucléosidiques\\de la Transcriptase Inverse)};
\node[class, right of=c2, xshift=1.5cm, align=center] (c3){II / INI\\(Inhibiteurs de l'Intégrase)};
\node[class, right of=c3, xshift=1.5cm, align=center] (c4){IP/r\\(Inhibiteurs de Protéase\\boostés)};
% Molecules 1ere ligne
\node[mol, below of=c1, align=center] (m1){\textbf{Colonne vertébrale 1ère ligne}\\TDF + 3TC (ou FTC)};
\node[mol, below of=c2, align=center] (m2){\textbf{Ancienne 1ère ligne}\\EFV (Éfavirenz)\\(Retrait progressif)};
\node[mol, below of=c3, align=center] (m3){\textbf{1ère ligne Préférée (OMS/CM)}\\DTG (Dolutégravir)\\Barrière haute, 1 cp/jour};
\node[mol, below of=c4, align=center] (m4){\textbf{2ème ligne / Alternative}\\ATV/r ou LPV/r\\(Charge pilulaire +)};
% Targets
\node[target, below of=m1, align=center] (t1){Cible : Transcriptase Inverse\\(Arrêt synthèse ADN viral)};
\node[target, below of=m3, align=center] (t3){Cible : Intégrase\\(Blocage intégration ADN viral)};
\node[target, below of=m4, align=center] (t4){Cible : Protéase\\(Blocage maturation virion)};
% Arrows
\draw[arrow] (c1) -- (m1); \draw[arrow] (m1) -- (t1);
\draw[arrow] (c2) -- (m2);
\draw[arrow] (c3) -- (m3); \draw[arrow] (m3) -- (t3);
\draw[arrow] (c4) -- (m4); \draw[arrow] (m4) -- (t4);
\end{tikzpicture}
\legende{Classes ARV, molécules de référence 1ère/2ème ligne Cameroun et cibles dans le cycle viral. INRTI = "fausses briques" (arrêt chaîne ADN). INNRTI = "clé cassée dans serrure" (blocage enzyme). II = blocage "ciseaux" (intégrase). IP = blocage "ciseaux" (protéase).}
\end{popfigure}

\begin{exemplebox}[Pourquoi Dolutégravir (DTG) est le standard 1ère ligne actuel ?]
\begin{itemize}
    \item \textbf{Barrière génétique haute :} Nécessite plusieurs mutations rares pour résister. Très peu d'échecs en 1ère ligne.
    \item \textbf{Puissance / Vitesse :} Charge virale indétectable souvent en 1-2 mois (vs 3-4 mois pour EFV). Crucial pour grossesse et TASP.
    \item \textbf{Tolérance :} Moins d'effets neuropsychiques (pas de troubles du sommeil, dépression comme EFV). Prise unique quotidienne, sans contrainte alimentaire.
    \item \textbf{Coût :} Génériques accessibles, \textbf{gratuits pour le patient} dans le public.
    \item \textbf{Vigilance :} \textbf{Prise de poids} possible (surtout femme + TDF/FTC). Surveillance poids/glycémie. \textbf{Grossesse :} Supplémentation \textbf{Acide Folique 5 mg/j} systématique chez femme en âge de procréer (réduction risque défaut tube neural).
\end{itemize}
\end{exemplebox}

\begin{popfigure}
\begin{tikzpicture}
\begin{axis}[
    width=\linewidth, height=7cm,
    grid=major, grid style={popline},
    xlabel={Temps sous ARV (Mois)}, ylabel={Log$_{10}$ Charge Virale (copies/mL)},
    xmin=0, xmax=12, ymin=1, ymax=6,
    legend style={at={(0.98,0.98)}, anchor=north east, fill=white, draw=popline, font=\small},
    tick label style={font=\small}, label style={font=\small},
]
% CV DTG
\addplot[thick, popblue, domain=0:12, samples=100] 
    ({x}, {5.5 - 1.5*ln(x+1) - 0.4*x}); 
\addlegendentry{DTG (1ère ligne actuelle)}
% CV EFV
\addplot[thick, poporange, dashed, domain=0:12, samples=100] 
    ({x}, {5.5 - 1.0*ln(x+1) - 0.25*x}); 
\addlegendentry{EFV (Ancienne 1ère ligne)}
% Seuil indétectable
\draw[popred, dashed, thick] (0, 1.7) -- (12, 1.7) node[pos=0.5, above, sloped, font=\tiny, popred] {Seuil détection (50 copies/mL)};
% Annotations
\node[popblue, font=\small, align=center] at (1.5, 4.5) {Chute rapide\\et indétectable};
\node[poporange, font=\small, align=center] at (4, 3.5) {Chute plus lente};
\end{axis}
\end{tikzpicture}
\legende{Cinétique de décroissance de la Charge Virale sous Trithérapie efficace. DTG (bleu) atteint l'indétectabilité (< 50 cp/mL) plus vite qu'EFV (orange pointillé).}
\end{popfigure}

\courssub{Objectifs Thérapeutiques, Observance et Résistances}

\begin{aretenir}[Objectifs de la Trithérapie (Cible "95-95-95" ONUSIDA / Cameroun)]
\begin{enumerate}
    \item \textbf{Virologique :} \textbf{Charge Virale (CV) < 50 copies/mL (Indétectable)} à \textbf{6 mois (M6)}. Critère principal de succès.
    \item \textbf{Immunologique :} \textbf{Restauration CD4+} (Gain $\approx$ 50-150 cellules/µL/an au début). Objectif : CD4 > 500/µL. Risque infections opportunistes négligeable si CD4 > 200/µL.
    \item \textbf{Clinique :} Prévention morbidité/mortalité (Sida, infections opportunistes, cancers). Qualité de vie normale.
    \item \textbf{Prévention (Santé Publique) :} \textbf{I = I} (CV < 200 cp/mL stable) $\rightarrow$ Zéro transmission sexuelle. PMTCT efficace.
\end{enumerate}
\textbf{Condition absolue :} \textbf{Observance > 95\%} (ne pas rater $> 1$ prise / mois pour 1 prise/jour). L'observance est le \textbf{premier prédicteur} de succès virologique.
\end{aretenir}

\begin{attention}[L'arrêt du traitement = Danger majeur]
\begin{itemize}
    \item \textbf{Rebond viral rapide :} En 2-4 semaines, CV remonte à niveau pré-traitement (réservoirs latents se réactivent).
    \item \textbf{Sélection de résistances :} Pendant la descente des ARV (demi-vies inégales), il y a une période de \textbf{monothérapie fonctionnelle} $\rightarrow$ Sélection mutants résistants.
    \item \textbf{Aggravation clinique :} Progression vers Sida, syndrome inflammatoire à la reprise tardive.
    \item \textbf{Message clé :} \textbf{ARV = À VIE.} (Sauf essais cliniques encadrés de rémission exceptionnelle).
\end{itemize}
\end{attention}

\begin{definition}[Résistance aux ARV : Mécanisme et Surveillance]
Sous pression sélective des ARV, le VIH (mutation spontanée fréquente) sélectionne des mutants porteurs de mutations conférant une \textbf{résistance} (le médicament ne bloque plus le virus).
\begin{itemize}
    \item \textbf{Résistance aux INRTI (ex: 3TC/FTC) :} Apparition rapide. Peut diminuer la "fitness" (capacité de réplication) du virus et resensibiliser au TDF/AZT.
    \item \textbf{Résistance aux INNRTI (ex: EFV/NVP) :} Mutation unique $\rightarrow$ Résistance croisée à toute la classe. Barrière génétique basse.
    \item \textbf{Résistance aux II (DTG) :} Voies multiples, barrière haute = émergence lente/rare en 1ère ligne.
\end{itemize}
\textbf{Test de résistance (Génotypage) :} Indiqué en \textbf{échec virologique confirmé} (CV > 1000 copies/mL à deux reprises, \textbf{sous observance vérifiée > 95\%}). Guide le choix de la 2ème/3ème ligne.
\end{definition}

\begin{exemplebox}[Coût et Enjeu Économique des Lignes de Traitement au Cameroun (Données approx. 2024)]
\begin{itemize}
    \item \textbf{1ère Ligne (TDF/3TC/DTG) :} Coût matière ~ \textbf{30 000 - 50 000 FCFA/an/patient}. \textbf{Gratuit pour le patient} (Budget État + Fonds Mondial). 1 comprimé/jour.
    \item \textbf{2ème Ligne (ex: AZT/3TC + ATV/r) :} Coût \textbf{5 à 10 fois plus élevé}. Charge pilulaire élevée (2 à 4 comprimés 2x/jour). Effets secondaires métaboliques majeurs.
    \item \textbf{3ème Ligne (Nouvelles classes) :} Coût très élevé, accès très limité (programmes spéciaux).
    \item \textbf{Enjeu :} La \textbf{préservation de la 1ère ligne (DTG)} par une observance parfaite est la meilleure assurance vie (médicale et économique) pour le patient et le système de santé.
\end{itemize}
\end{exemplebox}

\courssub{Exercices d'Intégration Type BEPC}

\begin{exoresolu}[Interprétation d'algorithme de dépistage et conduite à tenir]
\textbf{Énoncé :} Un jeune homme de 24 ans se présente au CDV après un rapport non protégé il y a 10 jours. Le Test de Dépistage (4ème génération) montre les bandes \textbf{C} et \textbf{Ag} visibles (faibles), Bandes T1 et T2 absentes.
\begin{enumerate}
    \item Interpréter ce résultat.
    \item Quelle est la conduite à tenir immédiate ? Justifier.
    \item Quel test demander pour lever le doute ? Pourquoi ?
\end{enumerate}
\tcblower
\textbf{Solution détaillée :}
\begin{enumerate}
    \item \textbf{Interprétation :} Test \textbf{Réactif pour Ag p24 seul}. La bande C valide le test. La présence de l'antigène p24 sans anticorps (T1/T2) suggère une \textbf{infection très récente (primo-infection)} : l'Ag p24 apparaît vers J10-15, les Ac apparaissent plus tard (J21-28). L'exposition date de 10 jours (compatibilité).
    \item \textbf{Conduite immédiate :} \textbf{Ne pas annoncer "Séropositif définitif".} Annoncer : "Le test montre la présence de l'antigène du virus, suggérant une infection récente. Nous devons confirmer par un test de confirmation." $\rightarrow$ Réaliser immédiatement le \textbf{Test de Confirmation}.
    \item \textbf{Test pour lever le doute :}
        \begin{itemize}
            \item \textbf{Si Confirmation Positif :} $\rightarrow$ Infection confirmée $\rightarrow$ Prise en charge urgente (bénéfice immunologique majeur en primo-infection).
            \item \textbf{Si Confirmation Négatif (Discordance) :} Hypothèses : Faux positif Test 1 (rare sur Ag) OU Fenêtre sérologique Test Confirmation (Ac pas encore détectables). \textbf{Gold Standard : PCR ARN VIH (Charge Virale) quantitative.}
        \end{itemize}
        \textbf{Justification PCR ARN :} Détecte l'ARN viral directement (fenêtre ~10 jours). Quantifie la réplication (CV souvent > $10^6$ cp/mL en primo-infection). Permet diagnostic certain avant séroconversion. Si PCR + $\rightarrow$ Traitement d'urgence. Si PCR - $\rightarrow$ Faux positif $\rightarrow$ Rassurer, re-test TRD à 6 semaines.
\end{enumerate}
\end{exoresolu}

\begin{exoresolu}[Gestion d'un échec virologique sous 1ère ligne DTG]
\textbf{Énoncé :} Une patiente de 32 ans, sous TDF/3TC/DTG depuis 18 mois (Observance déclarée "bonne"), présente une CV à \textbf{1 200 copies/mL} à M18 (M6 : < 50 ; M12 : < 50). CD4 = 450/µL.
\begin{enumerate}
    \item Qualifier cette situation selon les critères OMS.
    \item Quelle est la conduite à tenir \textbf{avant} de changer de traitement ? (Étapes clés).
    \item Si l'échec est confirmé (CV > 1000 à J+30 sous observance prouvée), quel test demander ? Quelle 2ème ligne probable ?
\end{enumerate}
\tcblower
\textbf{Solution détaillée :}
\begin{enumerate}
    \item \textbf{Échec virologique suspecté/confirmé.} Critère OMS : CV > 1000 copies/mL sous traitement (après 6 mois). Ici CV = 1200 cp/mL $\rightarrow$ Seuil dépassé.
    \item \textbf{Conduite avant changement (Impératif) :}
        \begin{enumerate}
            \item \textbf{Vérifier l'observance réelle :} Entretien non-jugeant, comptage pilules. Cause n°1 d'échec = sous-observance (oubli, rupture, effets secondaires, stigmatisation).
            \item \textbf{Corriger les causes :} Conseil adherence renforcé (pair éducateur, boîte à pilules, alarme, gestion effets indésirables).
            \item \textbf{Contrôle CV à 1 mois (M19) :} Si CV redevient < 50 cp/mL $\rightarrow$ "Blip" ou non-observance transitoire $\rightarrow$ \textbf{Maintenir 1ère ligne}.
        \end{enumerate}
    \item \textbf{Si échec confirmé (CV M19 > 1000 cp/mL sous observance prouvée) :}
        \begin{itemize}
            \item \textbf{Test de résistance (Génotypage) :} Obligatoire pour guider le choix.
            \item \textbf{2ème Ligne Standard Cameroun :} \textbf{AZT + 3TC + ATV/r (ou LPV/r).}
            \item \textbf{Raisonnement :} Changement de classe (IP/r) + Colonne vertébrale INRTI modifiée (AZT remplace TDF si résistance K65R suspectée, ou maintien TDF si M184V seule). AZT/3TC + IP/r a barrière génétique haute. Surveillance : Hb (AZT), Lipides/Glycémie/Foie (IP/r).
        \end{itemize}
\end{enumerate}
\end{exoresolu}

\begin{exercice}[Prévention Combinée : Cas Pratiques BEPC]
\textbf{Cas 1 :} Un couple sérodifférent (Homme VIH+ sous DTG/3TC/TDF depuis 8 mois, CV = 35 copies/mL ; Femme VIH-). Ils souhaitent un enfant.
\begin{enumerate}
    \item La femme doit-elle prendre la PrEP pendant la période de conception ? Justifier.
    \item Quels conseils donner pour l'accouchement et l'allaitement ?
\end{enumerate}

\textbf{Cas 2 :} Un infirmier se pique au doigt avec une aiguille ayant servi sur un patient VIH+ (CV = 50 000 cp/mL). L'accident survient à 14h.
\begin{enumerate}
    \item Quels sont les premiers gestes immédiats sur le site de piqûre ?
    \item Quel traitement PEP initier ? Pour quelle durée ? Quel délai maximal pour le débuter ?
    \item Quels examens de surveillance biologique prévoir ?
\end{enumerate}

\textbf{Cas 3 :} Une adolescente de 17 ans, travailleuse du sexe, VIH-, demande la PrEP. Elle a une hépatite B chronique (Ag HBs +, Foie sain).
\begin{enumerate}
    \item La PrEP (TDF/FTC) est-elle indiquée ? Justifier le double bénéfice.
    \item Quelle surveillance biologique spécifique mettre en place ?
\end{enumerate}
\end{exercice}

\begin{corrige}[Exercice Prévention Combinée]
\textbf{Cas 1 :}
1. \textbf{Non.} L'homme est sous ARV efficace, CV indétectable (< 50 cp/mL) depuis > 6 mois $\rightarrow$ \textbf{I = I s'applique}. Risque transmission sexuelle = 0. PrEP inutile pour la conception. (Dépistage IST systématique avant arrêt préservatif recommandé).
2. \textbf{Accouchement :} Voie basse possible si CV < 50 cp/mL à terme. Prophylaxie nourrisson : AZT sirop 4-6 semaines. \textbf{Allaitement :} Recommandé (OMS/Cameroun) \textbf{sous couverture ARV mère (CV indétectable)} + AZT nourrisson 6 sem. Allaitement exclusif 6 mois, puis diversifié jusqu'à 12-24 mois. Arrêt progressif.

\textbf{Cas 2 :}
1. \textbf{Premiers gestes :} Ne pas presser le site. Lavage abondant \textbf{eau + savon} (ou sérum physiologique) 5 min. Désinfection (Dakin / Bétadine / Alcool 70\%). \textbf{Ne pas saigner le doigt.} Déclaration immédiate (Registre AES, Médecin du travail).
2. \textbf{PEP :} Initiation \textbf{< 2h idéale, < 72h max}. Trithérapie standard : \textbf{TDF 300 mg + 3TC 300 mg (ou FTC 200 mg) + DTG 50 mg} (1 cp/jour) pendant \textbf{28 jours}. \textbf{Délai max : 72 heures.} Au-delà, efficacité non prouvée.
3. \textbf{Surveillance :} VIH J0, J28, 3 mois (6 mois si source en échec thérapeutique). Créatinine J0, J28 (TDF). Transaminases J0, J28. Grossesse (β-HCG) si femme.

\textbf{Cas 3 :}
1. \textbf{Oui, fortement indiquée.} Population clé + risque élevé. \textbf{Bénéfice double :} TDF/FTC traite \textbf{efficacement l'Hépatite B chronique} (suppression réplication VHB) \textbf{ET} prévient le VIH (PrEP). \textbf{Attention :} Arrêt PrEP $\rightarrow$ Risque \textbf{flare hépatique sévère} (réactivation VHB). $\rightarrow$ PrEP = Traitement VHB + Prévention VIH. Ne pas arrêter sans avis hépatologue.
2. \textbf{Surveillance :} VIH (TRD) tous les 3 mois. \textbf{Fonction rénale (Créatinine, DFG) tous les 6 mois} (TDF). \textbf{Marqueurs VHB :} ADN VHB quantitatif, Transaminases (ALAT/ASAT) tous les 6 mois. HBeAg/Ac anti-HBe annuel. Vaccination Hépatite A si non immunisée.
\end{corrige}

\coursec{Bilan synthèse \& Exercices d'intégration type BEPC}

Après avoir détaillé la prise en charge du VIH/Sida au Cameroun (dépistage, algorithme sérologique, schémas ARV, prophylaxies), nous rassemblons maintenant toutes les notions de la leçon dans une vision synthétique. Cette section te prépare à l'épreuve du BEPC : elle t'entraîne à lire des documents scientifiques (courbes, schémas, tableaux, textes), à croiser tes connaissances et à rédiger des réponses argumentées, structurées et rigoureuses.

\courssub{Synthèse schématique : Du virus à la prise en charge globale}

Pour réviser efficacement, rien ne vaut une carte mentale qui relie la structure du virus, son cycle de réplication, les cibles thérapeutiques, la dynamique immunitaire, les phases cliniques et la stratégie de santé publique.

\begin{popfigure}
\begin{center}\tcbox[colback=popink,colframe=popink,coltext=white,arc=9pt,boxsep=4pt,left=6pt,right=6pt]{\bfseries\small VIH / Sida (Leçon 7)}\end{center}
\vspace{-2pt}
\begin{tcbraster}[raster columns=3,raster equal height=rows,raster column skip=6pt,raster row skip=6pt,enhanced,blankest]
\begin{tcolorbox}[enhanced,colback=white,colframe=poporange,boxrule=0.9pt,arc=3pt,title={Structure \& Cibles},fonttitle=\bfseries\scriptsize,coltitle=white,colbacktitle=poporange,left=4pt,right=4pt,top=2pt,bottom=2pt,boxsep=1pt,toptitle=1.5pt,bottomtitle=1.5pt,halign=left]
\begin{itemize}[nosep,leftmargin=*,itemsep=1pt,font=\scriptsize]
\item ARN (+), RT, Intégrase, Protéase
\item Env : gp120/gp41
\item Cible : CD4+ (LT4, Macrophages, DC)
\end{itemize}
\end{tcolorbox}
\begin{tcolorbox}[enhanced,colback=white,colframe=popgreen,boxrule=0.9pt,arc=3pt,title={Cycle Réplication},fonttitle=\bfseries\scriptsize,coltitle=white,colbacktitle=popgreen,left=4pt,right=4pt,top=2pt,bottom=2pt,boxsep=1pt,toptitle=1.5pt,bottomtitle=1.5pt,halign=left]
\begin{itemize}[nosep,leftmargin=*,itemsep=1pt,font=\scriptsize]
\item Fixation CD4/CCR5/CXCR4
\item RT : ARN $\to$ ADN (Provirus)
\item Intégration (Intégrase)
\item Transcription / Traduction
\item Assemblage / Bourgeonnement (Protéase)
\end{itemize}
\end{tcolorbox}
\begin{tcolorbox}[enhanced,colback=white,colframe=poppurple,boxrule=0.9pt,arc=3pt,title={Cibles ARV (Classes)},fonttitle=\bfseries\scriptsize,coltitle=white,colbacktitle=poppurple,left=4pt,right=4pt,top=2pt,bottom=2pt,boxsep=1pt,toptitle=1.5pt,bottomtitle=1.5pt,halign=left]
\begin{itemize}[nosep,leftmargin=*,itemsep=1pt,font=\scriptsize]
\item IEP : TDF, 3TC, ABC
\item INNRTI : EFV, NVP, DOR
\item IP : LPV/r, ATV/r, DRV/r
\item II : DTG, RAL, BIC
\item IE : MVC (entrée)
\end{itemize}
\end{tcolorbox}
\begin{tcolorbox}[enhanced,colback=white,colframe=poppink,boxrule=0.9pt,arc=3pt,title={Dynamique Infection},fonttitle=\bfseries\scriptsize,coltitle=white,colbacktitle=poppink,left=4pt,right=4pt,top=2pt,bottom=2pt,boxsep=1pt,toptitle=1.5pt,bottomtitle=1.5pt,halign=left]
\begin{itemize}[nosep,leftmargin=*,itemsep=1pt,font=\scriptsize]
\item CV : Pic aigu $\to$ Set-point
\item CD4 : Chute progressive
\item Séroconversion (sem. 2-12)
\item Phases : Aiguë, Chronique, Sida
\end{itemize}
\end{tcolorbox}
\begin{tcolorbox}[enhanced,colback=white,colframe=popgold,boxrule=0.9pt,arc=3pt,title={Prise en Charge \& Prévention},fonttitle=\bfseries\scriptsize,coltitle=white,colbacktitle=popgold,left=4pt,right=4pt,top=2pt,bottom=2pt,boxsep=1pt,toptitle=1.5pt,bottomtitle=1.5pt,halign=left]
\begin{itemize}[nosep,leftmargin=*,itemsep=1pt,font=\scriptsize]
\item Dépistage : Algo 3 tests
\item ARV : TAR à vie (2 IEP + 1 II/INNRTI)
\item Objectif : CV $<$ 50 cp/mL (Indétectable)
\item PrEP / PPE / PMTCT
\item 95-95-95 ONUSIDA
\end{itemize}
\end{tcolorbox}
\end{tcbraster}

\legende{Carte mentale récapitulative de la leçon 7 : structure du VIH, cycle de réplication avec cibles ARV, dynamique charge virale / CD4, phases cliniques et continuum de prise en charge.}
\end{popfigure}

\begin{methode}[Réviser avec une carte mentale]
\begin{enumerate}
\item Place le thème central au milieu.
\item Identifie les \textbf{4-5 grands piliers} (Structure, Cycle, Traitement, Dynamique, Santé publique).
\item Pour chaque pilier, note les \textbf{mots-clés} (enzymes, molécules, valeurs seuils, noms de phases).
\item Trace des \textbf{flèches de lien} : ex. "RT $\xrightarrow{\text{cible}}$ IEP/INNRTI", "CD4 $< 200 \to$ Phase Sida".
\item Récite le cours en suivant les branches : c'est la meilleure preuve que tu as compris la logique globale.
\end{enumerate}
\end{methode}

\courssub{Analyse de documents types BEPC : Entraînement ciblé}

Au BEPC, l'épreuve de SVTEEHB propose souvent un dossier de 3 à 4 documents (courbes, schémas, tableaux, extraits de texte) suivi de questions de synthèse. Voici les 4 grands types de documents à maîtriser.

\courssub{Type 1 : Courbes Charge Virale (CV) et Taux de CD4 à annoter}

\begin{center}
\textit{Évolution typique non traitée de la charge virale (orange, échelle log à droite) et des lymphocytes T CD4+ (bleu, échelle linéaire à gauche). Repère : pic aigu, set-point, phase chronique, phase Sida (CD4 $< 200$), initiation ARV (an 10) et suppression virale.}
\end{center}

\begin{exoresolu}[Analyse de courbes CV/CD4]
\textbf{Document :} La figure ci-dessus montre l'évolution de la charge virale (CV) et du taux de lymphocytes T CD4+ chez un patient non traité pendant 10 ans, puis sous ARV.

\textbf{Questions :}
\begin{enumerate}
\item Identifier sur la courbe la période correspondant à l'infection primaire. Justifier par deux arguments tirés du document.
\item Quel est le seuil de CD4 qui définit le passage à la phase Sida ? Indiquer la valeur et l'année approximative sur le graphique.
\item Le patient commence un traitement ARV à l'année 10. Décrire l'évolution attendue de la CV et des CD4 sous traitement efficace.
\item Expliquer pourquoi la charge virale devient "indétectable" ($< 50$ copies/mL) mais que le patient n'est pas guéri.
\end{enumerate}

\tcblower
\textbf{Solution détaillée :}
\begin{enumerate}
\item \textbf{Infection primaire (an 0 à $\sim$0,3 an) :} 
    \begin{itemize}
    \item Pic de charge virale très élevé ($> 10^6$ copies/mL) dû à la réplication massive non contrôlée.
    \item Chute brutale et transitoire des CD4 (descendent sous 500/$\mu$L) avant remontée partielle.
    \end{itemize}
    \textit{Preuve doc :} Pic orange à 6 log ($10^6$) et flèche bleue descendante au début.
\item \textbf{Seuil Sida : CD4 $< 200$ cellules/$\mu$L.} Sur le graphique, la courbe bleue franchit 200 vers l'année 8-9. Apparition d'infections opportunistes (zone rouge).
\item \textbf{Sous ARV efficace (dès an 10) :}
    \begin{itemize}
    \item CV : Chute exponentielle (1 à 2 log par mois) jusqu'à $< 50$ copies/mL (indétectable) en 3-6 mois. Courbe verte descendante.
    \item CD4 : Remontée progressive ($\sim$ 50 à 150 cellules/an) mais souvent incomplète si traitement tardif. Courbe bleue qui repart à la hausse.
    \end{itemize}
\item \textbf{Indétectable $\neq$ Guéri :} Les ARV bloquent la réplication (RT, Intégrase, Protéase) mais \textbf{ne éliminent pas le provirus intégré} dans l'ADN des cellules réservoirs (LT4+ mémoire au repos). À l'arrêt du traitement, le rebond viral est quasi certain (semaines). \textit{Notion de réservoir viral (voir Définition).}
\end{enumerate}
\end{exoresolu}

\begin{attention}[Piège BEPC : "Charge virale indétectable = Guérison" ?]
\textbf{FAUX.} "Indétectable" ($< 50$ copies/mL) signifie \textbf{succès thérapeutique} : le virus ne se réplique plus activement, le système immunitaire se reconstruit, le risque de transmission sexuelle est nul (I = I : Indétectable = Intransmissible). Mais le \textbf{provirus persiste} dans les réservoirs (LT4+ mémoire, macrophages, SNC, ganglions). L'arrêt du traitement entraîne un rebond viral dans 99\% des cas. La guérison (éradication) reste un objectif de recherche (voir Ouverture Bac).
\end{attention}

\courssub{Type 2 : Schéma du cycle de réplication à compléter (enzymes + ARV)}

\begin{popfigure}
\begin{tikzpicture}[
  node distance=1.2cm and 1.5cm,
  box/.style={rectangle, draw=popdark, thick, fill=popcream, rounded corners, minimum width=2.8cm, minimum height=1.1cm, align=center, font=\small},
  arrow/.style={->, thick, popdark},
  enzyme/.style={rectangle, draw=poporange, fill=poporangeL, rounded corners, minimum width=2.2cm, minimum height=0.7cm, align=center, font=\scriptsize},
  arv/.style={rectangle, draw=popgreen, fill=popgreenL, rounded corners, minimum width=2.2cm, minimum height=0.7cm, align=center, font=\scriptsize},
]
% Virus entry
\node[box, align=center] (virus){Virion VIH-1\\ (ARN +, RT, Int, Protéase)};
\node[box, right=of virus, align=center] (fix){Fixation : gp120 $\to$ CD4 \\ + Co-récepteur (CCR5/CXCR4)};
\node[box, right=of fix, align=center] (fusion){Fusion membranaire \\ $\to$ Libération capside dans cytoplasme};
% Reverse Transcription
\node[enzyme, below=of fix, xshift=1.5cm, align=center] (rt){Transcriptase Inverse (RT)\\ \textbf{ARN $\to$ ADN double brin}};
\node[arv, below=of rt, yshift=-0.3cm, align=center] (arv_rt){Cibles ARV :\\ IEP (TDF, 3TC, ABC)\\ INNRTI (EFV, NVP, DOR)};
% Integration
\node[box, below=of fusion, yshift=-1.5cm, align=center] (cyto){ADN viral double brin\\ (Pré-complexe d'intégration)};
\node[enzyme, right=of cyto, xshift=1.5cm, align=center] (int){Intégrase\\ \textbf{Intégration dans ADN hôte}};
\node[arv, below=of int, yshift=-0.3cm, align=center] (arv_int){Cible ARV :\\ II (DTG, RAL, BIC, CAB)};
% Provirus
\node[box, below=of cyto, align=center] (prov){Provirus intégré\\ (ADN viral dans chromosome hôte)};
% Transcription/Translation
\node[box, left=of prov, xshift=-1.5cm, align=center] (trans){Transcription (ARNm viraux)\\ $\to$ Traduction (Polyprotéines)};
\node[box, left=of trans, xshift=-1.5cm, align=center] (assemb){Assemblage : Polyprotéines\\ + ARN génomique $\to$ Particule immature};
% Protease / Budding
\node[enzyme, below=of assemb, yshift=-0.5cm, align=center] (prot){Protéase virale\\ \textbf{Clivage polyprotéines $\to$ Virion mature}};
\node[arv, below=of prot, yshift=-0.3cm, align=center] (arv_prot){Cibles ARV :\\ IP (LPV/r, DRV/r, ATV/r)};
\node[box, below=of prot, yshift=-0.5cm, align=center] (mature){Virion mature infectieux\\ (Bourgeonnement)};
% Arrows
\draw[arrow] (virus) -- (fix) node[midway, above, font=\tiny] {1};
\draw[arrow] (fix) -- (fusion) node[midway, above, font=\tiny] {2};
\draw[arrow] (fusion) -- (cyto) node[midway, right, font=\tiny] {3};
\draw[arrow] (cyto) -- (prov) node[midway, right, font=\tiny] {4};
\draw[arrow] (prov) -- (trans) node[midway, above, font=\tiny] {5};
\draw[arrow] (trans) -- (assemb) node[midway, above, font=\tiny] {6};
\draw[arrow] (assemb) -- (mature) node[midway, right, font=\tiny] {7};
% Enzyme links
\draw[poporange, dashed, thick] (fix.south) |- (rt.north);
\draw[poporange, dashed, thick] (cyto.east) |- (int.west);
\draw[poporange, dashed, thick] (assemb.south) |- (prot.north);
% ARV links
\draw[popgreen, dashed, thick] (rt.south) -- (arv_rt.north);
\draw[popgreen, dashed, thick] (int.south) -- (arv_int.north);
\draw[popgreen, dashed, thick] (prot.south) -- (arv_prot.north);
% Entry inhibitor
\node[arv, above=of fix, yshift=0.5cm, align=center] (arv_entry){IE : Maraviroc (CCR5)\\ Enfuvirtide (Fusion)};
\draw[popgreen, dashed, thick] (arv_entry.south) -- (fix.north);
\end{tikzpicture}
\legende{Schéma du cycle de réplication du VIH-1 avec les 3 enzymes cibles (RT, Intégrase, Protéase) et les classes d'ARV correspondantes (IEP/INNRTI, II, IP). En vert : Inhibiteurs d'Entrée (IE).}
\end{popfigure}

\begin{exercice}[Compléter le cycle - Type BEPC]
\textbf{Document :} Schéma incomplet du cycle de réplication (figure ci-dessus avec cases vides).
\textbf{Consignes :}
\begin{enumerate}
\item Sur le schéma, écrire le nom de l'enzyme virale agissant à chaque étape numérotée 3, 4, 7.
\item Indiquer, pour chaque enzyme, la classe(s) d'ARV qui l'inhibe(nt) et donner un exemple de molécule (DCI).
\item L'étape 2 (fixation/fusion) est la cible d'une classe d'ARV peu utilisée en 1ère ligne. Laquelle ? Citer un exemple.
\item Une mutation sur le gène codant la transcriptase inverse (ex: K103N) confère une résistance à l'EFV. Cette mutation rend-elle le virus résistant au DTG ? Justifier.
\end{enumerate}
\end{exercice}

\begin{corrige}[Exercice Compléter le cycle]
\begin{enumerate}
\item \textbf{Étape 3 (RT) :} Transcriptase Inverse. \textbf{Étape 4 (Intégration) :} Intégrase. \textbf{Étape 7 (Maturation) :} Protéase.
\item \begin{itemize}
    \item RT : IEP (TDF, 3TC, ABC, FTC) + INNRTI (EFV, NVP, DOR, ETV, RPV).
    \item Intégrase : II (DTG, RAL, BIC, CAB). \textit{Classe de choix en 1ère ligne OMS 2024.}
    \item Protéase : IP (LPV/r, ATV/r, DRV/r). Utilisés en 2ème ligne ou cas particuliers.
  \end{itemize}
\item \textbf{Inhibiteurs d'Entrée (IE) :} Maraviroc (antagoniste CCR5), Enfuvirtide (inhibiteur fusion, injectable). Peu utilisés en 1ère ligne (coût, voie admin, tropisme viral à tester pour MVC).
\item \textbf{Non.} K103N est une mutation conférant une résistance aux \textbf{INNRTI} (EFV, NVP). Le DTG (Dolutégravir) est un \textbf{Inhibiteur d'Intégrase (II)}. Cibles différentes $\to$ pas de résistance croisée (sauf cas rares multi-résistances). \textit{Voir Propriété : Résistance croisée.}
\end{enumerate}
\end{corrige}

\courssub{Type 3 : Cas clinique complet (Sérologie, Stadification, Prise en charge)}

\begin{exoresolu}[Cas clinique type BEPC - Intégration]
\textbf{Situation :} Marie, 28 ans, enceinte de 14 SA (semaines d'aménorrhée), consulte au Centre de Santé de District de Bafoussam pour sa 1ère CPN (Consultation Prénatale). Le test rapide Determine\textsuperscript{TM} HIV-1/2 est \textbf{réactif}. Le test de confirmation Uni-Gold\textsuperscript{TM} est \textbf{réactif}. Le test de discrimination Stat-Pak\textsuperscript{TM} est \textbf{réactif pour VIH-1, non réactif pour VIH-2}. Le taux d'hémoglobine est 10,8 g/dL. Elle n'a pas de signes cliniques d'infection opportuniste. Son mari refuse de se faire tester.

\textbf{Documents fournis :}
\begin{itemize}
\item \textbf{Doc 1 :} Algorithme national de dépistage VIH (3 tests : Dépistage $\to$ Confirmation $\to$ Discrimination).
\item \textbf{Doc 2 :} Classification clinique OMS (Stade 1 à 4) et Immunologique (CD4).
\item \textbf{Doc 3 :} Schéma thérapeutique 1ère ligne Cameroun 2024 : TDF + 3TC + DTG (TLD).
\item \textbf{Doc 4 :} Protocole PMTCT : ARV mère (TLD à vie) + AZT+NVP nourrisson 6 sem + Allaitement maternel exclusif 6 mois + ARV mère.
\end{itemize}

\textbf{Questions :}
\begin{enumerate}
\item Interpréter le résultat sérologique de Marie selon l'algorithme national. Quel est son statut VIH ? Quel type ?
\item Stadifier cliniquement Marie selon l'OMS. Justifier.
\item Proposer la prise en charge ARV de Marie (schéma, moment de début, surveillance).
\item Détailler la prophylaxie pour le nouveau-né et les recommandations d'allaitement.
\item Quel conseil donner à Marie concernant son mari ? (Notion de "Partenaire non testé").
\end{enumerate}

\tcblower
\textbf{Solution détaillée :}
\begin{enumerate}
\item \textbf{Statut : VIH-1 POSITIF.}
    \begin{itemize}
    \item Test 1 (Dépistage) : Réactif $\to$ Poursuivre.
    \item Test 2 (Confirmation) : Réactif $\to$ Infection VIH confirmée.
    \item Test 3 (Discrimination) : VIH-1 Réactif / VIH-2 Non réactif $\to$ \textbf{Infection à VIH-1}.
    \end{itemize}
    \textit{Connaissance :} Algo 3 tests obligatoire au Cameroun. Un seul test réactif ne suffit pas.
\item \textbf{Stade clinique OMS : Stade 1 (Asymptomatique).}
    \begin{itemize}
    \item Aucune infection opportuniste, pas de fièvre prolongée, pas de perte de poids $> 10\%$, pas de diarrhée chronique.
    \item Grossesse $\neq$ stade 3 ou 4. Anémie modérée (Hb 10,8) non spécifique au VIH.
    \end{itemize}
    \textit{Note :} Le stade immunologique nécessite le taux de CD4 (non fourni ici). Si CD4 $< 350$ $\to$ Stade immunologique avancé.
\item \textbf{Prise en charge ARV :}
    \begin{itemize}
    \item \textbf{Schéma :} TLD = \textbf{TDF 300 mg + 3TC 300 mg + DTG 50 mg} (1 cp/jour). \textit{1ère ligne recommandée OMS/Cameroun.}
    \item \textbf{Moment :} \textbf{Immédiat} (Test \& Treat), dès l'annonce, même jour si possible. Grossesse = urgence PMTCT.
    \item \textbf{Surveillance :} CV à 3 mois (cible $< 50$ cp/mL), puis 6 mois. CD4 à M6, M12. Hémoglobine, créatinine (TDF), glycémie (DTG), grossesse (échographies).
    \end{itemize}
\item \textbf{PMTCT Nouveau-né :}
    \begin{itemize}
    \item \textbf{Prophylaxie :} AZT (zidovudine) sirop 2 fois/jour + NVP (névirapine) sirop 1 fois/jour pendant \textbf{6 semaines} (schéma combiné recommandé si mère sous ARV $< 4$ sem avant accouchement ou CV non supprimée). Si mère CV indétectable $> 4$ sem : NVP seul 6 sem.
    \item \textbf{Allaitement :} \textbf{Allaitement maternel exclusif (AME) 6 mois}, puis introduction alimentation complémentaire + allaitement jusqu'à 12-24 mois. \textbf{Mère SOUS ARV avec CV indétectable} = risque transmission $< 1\%$.
    \item \textbf{Dépistage enfant :} PCR-ADN à 6 sem (naissance si risque haut), 6 mois, fin allaitement, 18 mois (sérologie).
    \end{itemize}
\item \textbf{Conseil partenaire (Mari) :}
    \begin{itemize}
    \item \textbf{Tester d'urgence} : Proposer autotest (VIH Auto-Test) ou venue au CSD. "Test \& Treat" s'applique à lui aussi.
    \item \textbf{PrEP d'urgence} : Si mari séronégatif, proposer PrEP (TDF+3TC) pendant grossesse/allaitement (période à haut risque).
    \item \textbf{Communication :} Approche non jugeante, confidentialité, bénéfice pour le couple et l'enfant. Pair-éducateur peut aider.
    \end{itemize}
\end{enumerate}
\end{exoresolu}

\courssub{Type 4 : Document sur la résistance aux ARV (Séquence, Mutation, Adaptation)}

\begin{exemplebox}[Résistance au VIH : Du nucléotide au changement de ligne]
\textbf{Document :} Extrait de séquence nucléotidique du gène \textit{pol} (codant RT) chez un patient sous TLD (TDF+3TC+DTG) en échec virologique (CV $> 1000$ cp/mL à M12).
\begin{center}
\texttt{... AAG AAG \textbf{AAC} CTT GGT ...} \quad (Codons 101-105 RT) \\
\texttt{... Lys Lys \textbf{Asn} Leu Gly ...} \quad (Acides aminés : K K \textbf{N} L G)
\end{center}
Séquence sauvage (référence) au codon 103 : \texttt{AAA} (Lysine, K). Séquence patiente : \texttt{AAC} (Asparagine, N). Mutation \textbf{K103N}.

\textbf{Questions :}
\begin{enumerate}
\item Identifier le type de mutation (transition/transversion) et son effet sur l'acide aminé.
\item Quelle classe d'ARV est concernée par la mutation K103N ? Citer les molécules touchées.
\item Le patient est sous DTG (Inhibiteur d'Intégrase). K103N affecte-t-elle l'efficacité du DTG ? Justifier.
\item Proposer un schéma de 2ème ligne adapté pour ce patient (OMS 2024).
\end{enumerate}

\textbf{Réponses :}
\begin{enumerate}
\item \textbf{Transition} (A $\to$ C en 3ème position du codon AAA $\to$ AAC). \textbf{Non-synonyme} : changement d'acide aminé Lysine (K) $\to$ Asparagine (N) en position 103.
\item \textbf{INNRTI (Inhibiteurs Non Nucléosidiques de la Transcriptase Inverse).} Confère une résistance de haut niveau à l'\textbf{EFV} (Éfavirenz) et à la \textbf{NVP} (Névirapine). Résistance croisée au sein des INNRTI de 1ère génération. L'\textbf{ETV} (Étravirine) et la \textbf{DOR} (Doravirine) conservent souvent une activité partielle.
\item \textbf{Non.} K103N est sur la \textbf{Transcriptase Inverse}. Le DTG cible l'\textbf{Intégrase}. Pas de résistance croisée entre classes. Le DTG reste pleinement actif.
\item \textbf{2ème ligne OMS 2024 (échec sous TLD) :} \textbf{TDF + 3TC + DTG} \textbf{EST} la 1ère ligne. En échec sous DTG (rare mais possible si mutations intégrase), on passe à un schéma à base de \textbf{IP boosté (DRV/r ou ATV/r) + 2 IEP (souvent AZT + 3TC ou TDF + 3TC)}. \textit{Exemple : AZT + 3TC + DRV/r 2 fois/jour.} Si échec sous ancien schéma EFV-based : passer à TLD (DTG-based).
\end{enumerate}
\end{exemplebox}

\begin{propriete}[Résistance croisée]
La résistance croisée existe \textbf{au sein d'une même classe} d'ARV (même cible enzymatique).
\begin{itemize}
\item \textbf{Exemple :} Mutation K103N (RT) $\to$ Résistance à tous les INNRTI de 1ère génération (NVP, EFV). Résistance partielle à ETV, RPV. DOR souvent active.
\item \textbf{Exception :} Les Inhibiteurs d'Intégrase de 2ème génération (DTG, BIC, CAB) ont un \textbf{barrière génétique élevée} : plusieurs mutations nécessaires pour résistance. Résistance croisée faible entre RAL/EVG (1ère gén) et DTG/BIC (2ème gén).
\end{itemize}
\textbf{Pas de résistance croisée entre classes différentes} (ex: mutation RT ne protège pas contre II ou IP).
\end{propriete}

\courssub{Argumentation écrite : Dépistage précoce et Gratuité des ARV, un investissement rentable}

\begin{methode}[Rédiger une argumentation santé publique (Type BEPC/Bac)]
\begin{enumerate}
\item \textbf{Thèse claire :} "Le dépistage précoce et la gratuité des ARV sont des investissements hautement rentables pour la santé publique."
\item \textbf{Argument 1 (Individuel) :} Traitement précoce $\to$ CV indétectable $\to$ Immunité préservée (CD4 hauts) $\to$ Espérance de vie quasi normale, qualité de vie, zéro transmission sexuelle (I=I).
\item \textbf{Argument 2 (Collectif - Prévention) :} "Treatment as Prevention" (TasP). Chaque personne sous ARV efficace ne transmet plus. Modélisation : 95-95-95 $\to$ Fin de l'épidémie (incidence $\to$ 0).
\item \textbf{Argument 3 (Économique - Coût-efficacité) :} Coût ARV/an/patient $\approx$ 60 000 - 100 000 FCFA (génériques). Coût hospitalisation Sida (infections opportunistes, séjours réanimation) $\gg$ 500 000 FCFA/épisode. Coût orphelins, perte productivité. \textbf{Investir 1 FCFA en prévention/traitement $\to$ économise 10-20 FCFA en soins curatifs.}
\item \textbf{Argument 4 (Équité/Droits) :} Gratuité = suppression barrière financière. Droit à la santé (Constitution, OMS). Lutte contre stigmatisation (accès aux soins sans discrimination).
\item \textbf{Conclusion :} Dépistage + ARV gratuits = \textbf{meilleur retour sur investissement sanitaire} (OMS, ONUSIDA, Banque Mondiale). Au Cameroun, progression vers 95-95-95 prouve la faisabilité.
\end{enumerate}
\end{methode}

\courssub{Éducation à la santé : Démystification, Lutte contre la stigmatisation, Rôle des pairs}

\begin{aretenir}[Messages clés à retenir et à transmettre (Éducation par les pairs)]
\begin{itemize}
\item \textbf{Transmission :} \textbf{SEULEMENT} 3 voies : Sexuelle (non protégée), Sanguine (aiguilles, transfusion non sécurisée), Mère-Enfant (grossesse, accouchement, allaitement \textbf{SANS} ARV).
\item \textbf{PAS de transmission :} Salut (main, baiser social), Moustiques/Insectes, Repas/Ustensiles partagés, Toilettes/Piscine, Transpiration/Larmes/Salive (sauf volume massif + lésion buccale), Soins quotidiens (hygiène, change).
\item \textbf{Stigmatisation = Obstacle majeur :} Peur du test, cachetonnage, abandon traitement, dépression, propagation silencieuse.
\item \textbf{Droits du patient (Loi 2023 Cameroun) :} Confidentialité absolue, Non-discrimination (travail, école, santé), Accès aux soins gratuits, Consentement éclairé (sauf urgence vitale).
\item \textbf{Rôle des Pairs-Éducateurs (PE) :} Jeunes formés $\to$ sensibilisation pairs (langage adapté), accompagnement au dépistage, soutien observance (rappels, groupes de parole), lutte rumeurs, orientation vers services (CSD, Hôpital de Jour, ONG).
\end{itemize}
\end{aretenir}

\begin{savaistu}[Cameroun : Chiffres clés \& Progrès (EDS-MICS 2018, Rapports PNLS 2023)]
\begin{itemize}
\item \textbf{Prévalence adultes (15-49 ans) :} \textbf{2,7\%} (Femmes \textbf{4,1\%} > Hommes \textbf{2,1\%}). Facteurs : biologiques (surface muqueuse), sociaux (violences, dépendance économique, mariages précoces).
\item \textbf{PMTCT (Prévention Mère-Enfant) :} Taux transmission $\mathbf{< 5\%}$ (objectif $< 2\%$). Couverture ARV femmes enceintes $> 95\%$. \textbf{Gain majeur :} De $> 20\%$ (sans ARV) à $< 5\%$ aujourd'hui. Objectif \textbf{Élimination (validation OMS)} en cours.
\item \textbf{Cascade 95-95-95 (2023) :} $\sim$ \textbf{90\% - 85\% - 90\%}. Progression notable depuis 2015 (60-40-70). Gap principal : \textbf{2ème 95 (Lien aux soins / Rétention)}. Hommes, adolescents, populations clés (HSH, TS, UD) moins bien couverts.
\item \textbf{ARV Gratuits depuis 2007} (Décret Premier Ministre). Approvisionnement : Fonds Mondial, PEPFAR, Budget État. Génériques préqualifiés OMS.
\end{itemize}
\end{savaistu}

\begin{exemplebox}[Objectifs 95-95-95 ONUSIDA (2025) \& Continuum de soins]
\begin{center}
\begin{tikzpicture}[node distance=1.5cm, >=Stealth]
\node[draw=popblue, thick, fill=popblueL, rounded corners, minimum width=3.5cm, minimum height=1.2cm, align=center] (s1) {\textbf{1. DÉPISTER}\\\textbf{95\%} PVVIH connaissent\\leur statut};
\node[draw=poporange, thick, fill=poporangeL, rounded corners, minimum width=3.5cm, minimum height=1.2cm, align=center, right=of s1] (s2) {\textbf{2. LIER AUX SOINS}\\\textbf{95\%} des diagnostiqués\\sous ARV};
\node[draw=popgreen, thick, fill=popgreenL, rounded corners, minimum width=3.5cm, minimum height=1.2cm, align=center, right=of s2] (s3) {\textbf{3. TRAITER}\\\textbf{95\%} des traités ont\\CV supprimée ($< 50$ cp/mL)};
\node[draw=poppurple, thick, fill=poppurpleL, rounded corners, minimum width=3.5cm, minimum height=1.2cm, align=center, below=of s2, yshift=-1cm] (impact) {\textbf{IMPACT :}\\\textbf{Fin de l'épidémie}\\\textit{Incidence $\to$ 0}\\\textit{Mortalité $\downarrow\downarrow$}\\\textit{Stigmatisation $\downarrow$}};
\draw[->, thick, popdark] (s1) -- (s2);
\draw[->, thick, popdark] (s2) -- (s3);
\draw[->, thick, popdark] (s3) -- (impact);
\draw[->, thick, popdark, dashed] (impact) -- (s1) node[midway, left, font=\small, align=center]{Rétention\\Vie longue};
\end{tikzpicture}
\end{center}
\textbf{Cameroun 2023 :} \textbf{90\% - 85\% - 90\%}. \textit{Progression :} 1er 95 presque atteint (dépistage index, autotests). 2ème 95 : défi rétention (hommes, jeunes, stigma). 3ème 95 : bon (DTG puissant, observance améliorée). \textbf{Investir sur le 2ème 95 (accueil, pairs, horaires élargis, multi-mois) = levier maximal.}
\end{exemplebox}

\courssub{Ouverture Bac (Complément) : Stratégies de guérison en recherche}

\begin{definition}[Réservoir viral]
Ensemble des cellules (principalement \textbf{lymphocytes T CD4+ mémoire au repos}, mais aussi macrophages, cellules dendritiques, cellules du SNC, tissu lymphoïde) contenant un \textbf{provirus intact, réplication-compétent}, intégré dans l'ADN hôte, \textbf{silencieux} (pas de transcription virale) sous ARV. Ces cellules ont une durée de vie très longue (années-décennées) et ne sont pas éliminées par le système immunitaire ni par les ARV actuels. Elles sont la cause du rebond viral systématique à l'arrêt du traitement.
\end{definition}

\begin{savaistu}[Recherche guérison : 3 pistes principales (Niveau Bac / Supérieur)]
\begin{enumerate}
\item \textbf{"Shock and Kill" (Réactiver \& Éliminer) :}
    \begin{itemize}
    \item \textit{Shock :} Agents de réversibilité de la latence (LRAs : inhibiteurs HDAC, agonistes TLR, inhibiteurs BET) pour forcer l'expression du provirus $\to$ présentation antigènes viraux à la surface.
    \item \textit{Kill :} Renforcer l'immunité (Vaccins thérapeutiques, anticorps neutralisants larges bnAbs, cellules CAR-T anti-VIH, NK activées) pour éliminer les cellules réactivées.
    \item \textit{Statut :} Essais cliniques (phase I/II). "Shock" partiel (réactivation faible), "Kill" insuffisant. Toxicité des LRAs. Pas de guérison durable pour l'instant.
    \end{itemize}
\item \textbf{"Block and Lock" (Bloquer \& Enfermer) :}
    \begin{itemize}
    \item Inverse : Maintenir le provirus \textbf{éteint définitivement} (silencement épigénétique profond) sans le réactiver. Molécules : Tat inhibiteurs, LEDGINs (modulateurs intégrase), CRISPR-épi (éditeur épigénétique).
    \item Objectif : "Remission fonctionnelle" sans ARV à vie. Virus présent mais inoffensif.
    \end{itemize}
\item \textbf{Édition Génique (CRISPR/Cas9) :}
    \begin{itemize}
    \item \textbf{Cible 1 : CCR5 (Co-récepteur).} Rendre les cellules résistantes à l'entrée (mutation $\Delta$32 naturelle). Ex : Essai \textit{EB-001} (ex vivo sur cellules souches hématopoïétiques autologues). Succès partiel (chimericité).
    \item \textbf{Cible 2 : Provirus intégré.} Couper l'ADN viral hors du génome hôte. \textit{Défi :} Spécificité (effets hors cible), livraison \textit{in vivo} à tous les réservoirs, diversité virale (échappement).
    \item \textit{Cas "Berlin Patient" \& "London Patient" :} Greffe de moelle donneur CCR5 $\Delta$32/$\Delta$32 pour leucémie $\to$ Guérison VIH. Preuve de concept mais \textbf{non généralisable} (risque mortel 15-20\%, donneurs rares).
    \end{itemize}
\item \textbf{Vaccins (Préventifs \& Thérapeutiques) :}
    \begin{itemize}
    \item \textbf{Échecs récents :} \textit{Mosaico} (Janssen, vecteur Ad26 + protéine, HSH/TS Amériques/Europe) \textbf{arrêté 2023} (inefficacité). \textit{Imbokodo} (Afrique sub-saharienne, femmes) \textbf{arrêté 2021} (inefficacité).
    \item \textbf{Nouveaux espoirs :} Vaccins \textbf{ARNm} (Moderna, IAVI, Scripps) : immunogènes "germline-targeting" pour induire bnAbs (anticorps neutralisants larges). Essais Phase I (ex: IAVI G001, eOD-GT8 60mer mRNA) montrent induction précurseurs bnAbs. \textbf{Vecteurs CMV} (rhCMV) : réponse T CD8+ inhabituelle (non restreinte MHC-Ia), contrôle viral chez singes (SIV). Long chemin vers l'homme.
    \end{itemize}
\end{enumerate}
\textbf{Message pour toi :} La guérison n'est pas pour demain, mais la science progresse. \textbf{Aujourd'hui : ARV à vie = vie longue, saine, non transmissible.} C'est déjà une victoire immense.
\end{savaistu}

\courssub{Sujet type BEPC complet (Énoncé + Documents + Corrigé détaillé)}

\textit{Ce sujet est conçu pour t'entraîner en conditions réelles. Chronomètre-toi : 1h30.}

\begin{exercice}[Sujet BEPC 2025 - SVTEEHB : Le VIH/Sida]
\textbf{Durée : 1h30} \quad \textbf{Coefficient : 1} \quad \textbf{Note sur 20}

\textbf{Consignes générales :} Lisez attentivement tous les documents avant de répondre. Les réponses doivent être justifiées par des éléments tirés des documents \textbf{ET} par vos connaissances. La qualité de la rédaction, la précision du vocabulaire scientifique et la logique du raisonnement sont évaluées.

\vspace{0.5cm}
\textbf{DOCUMENT 1 : Évolution de la charge virale et des CD4 chez un patient non traité}
\begin{center}
\begin{tikzpicture}
\begin{axis}[
  width=0.9\linewidth, height=7cm,
  xlabel={Temps (mois après infection)}, ylabel={CD4 (/$\mu$L)},
  xmin=0, xmax=60, ymin=0, ymax=1200,
  axis y line*=left,
  ytick={0,200,350,500,800,1000},
  grid=major, grid style={popline!20},
  legend style={at={(0.02,0.98)}, anchor=north west, font=\footnotesize, fill=white},
]
\addplot[popblue, very thick, domain=0:60, samples=300] 
  {1000 - 800/(1+exp(-0.8*(x-3))) - 12*x};
\addlegendentry{CD4}
\end{axis}
\begin{axis}[
  width=0.9\linewidth, height=7cm,
  xmin=0, xmax=60, ymin=0, ymax=7,
  axis y line*=right,
  ytick={0,2,4,6},
  yticklabels={$10^0$, $10^2$, $10^4$, $10^6$},
  ylabel={CV (copies/mL) $\log_{10}$},
  hide x axis,
  legend style={at={(0.98,0.98)}, anchor=north east, font=\footnotesize, fill=white},
]
\addplot[poporange, very thick, domain=0:1, samples=50] {6 - 5.5*x};
\addplot[poporange, very thick, domain=1:12, samples=50] {4.2 + 0.3*sin(deg(x))};
\addplot[poporange, very thick, domain=12:40, samples=100] {4.5 + 0.03*(x-12)};
\addplot[poporange, very thick, domain=40:60, samples=100] {5.1 + 0.08*(x-40)^1.5};
\addlegendentry{Charge Virale}
\end{axis}
\node[font=\small, align=center, popblue] at (3.5, 10.5) {Phase 1};
\node[font=\small, align=center, poporange] at (15, 10.5) {Phase 2};
\node[font=\small, align=center, popred] at (50, 10.5) {Phase 3};
\end{tikzpicture}
\end{center}

\vspace{0.5cm}
\textbf{DOCUMENT 2 : Schéma simplifié du cycle de réplication du VIH-1}
\begin{center}
\begin{tikzpicture}[node distance=1cm, box/.style={draw, thick, fill=popcream, rounded corners, minimum width=2.5cm, minimum height=0.9cm, align=center, font=\small}, arrow/.style={->, thick}]
\node[box] (a) {Virion extracellulaire};
\node[box, right=of a] (b) {Fixation gp120-CD4-CCR5};
\node[box, right=of b] (c) {Pénétration / Décapsidage};
\node[box, below=of c, yshift=-0.5cm, align=center] (d){\textbf{Étape X}\\ ARN viral $\to$ ADN viral};
\node[box, right=of d, align=center] (e){\textbf{Étape Y}\\ Intégration ADN viral / ADN hôte};
\node[box, right=of e] (f) {Transcription / Traduction};
\node[box, below=of f, yshift=-0.5cm, align=center] (g){\textbf{Étape Z}\\ Assemblage / Maturation};
\node[box, left=of g] (h) {Virion immature};
\draw[arrow] (a) -- (b); \draw[arrow] (b) -- (c); \draw[arrow] (c) -- (d);
\draw[arrow] (d) -- (e); \draw[arrow] (e) -- (f); \draw[arrow] (f) -- (g);
\draw[arrow] (g) -- (h); \draw[arrow] (h) -- ++(-2,0) |- (a);
\node[font=\small, poporange] at (3, -3.5) {Enzymes virales : \textbf{RT}, \textbf{Intégrase}, \textbf{Protéase}};
\end{tikzpicture}
\end{center}

\vspace{0.5cm}
\textbf{DOCUMENT 3 : Extrait du Guide National de Prise en Charge VIH Cameroun 2024}
\begin{itemize}
\item \textbf{1ère ligne (Adultes/Ados $> 30$ kg) :} \textbf{TLD} = Ténofovir (TDF) + Lamivudine (3TC) + Dolutégravir (DTG). 1 cp/jour.
\item \textbf{2ème ligne (Échec 1ère ligne) :} AZT + 3TC + DRV/r (Darunavir/ritonavir) 2 cp/jour (matin/soir).
\item \textbf{Échec virologique :} CV $> 1000$ copies/mL à deux reprises (3 mois d'intervalle) sous bonne observance.
\item \textbf{Surveillance :} CV à M3, M6, puis tous les 6 mois. Cible : \textbf{Indétectable ($< 50$ cp/mL)}.
\item \textbf{PMTCT :} Femme enceinte $\to$ TLD à vie (Test \& Treat). Nourrisson : Prophylaxie NVP 6 sem (si mère CV indétectable $> 4$ sem) ou AZT+NVP 6 sem (autres cas). Allaitement maternel exclusif 6 mois recommandé.
\end{itemize}

\vspace{0.5cm}
\textbf{DOCUMENT 4 : Résistance aux ARV - Mutations majeures}
\begin{center}
\begin{tabularx}{\linewidth}{|l|X|X|}\hline
\textbf{Classe ARV} & \textbf{Mutation exemple (RT / Intégrase / Protéase)} & \textbf{Effet} \\ \hline
IEP (TDF, 3TC) & M184V (RT) & Résistance haute 3TC/FTC. \textbf{Hypersensibilise} au TDF/AZT. \\ \hline
INNRTI (EFV, NVP) & \textbf{K103N} (RT) & Résistance croisée haute EFV/NVP. \\ \hline
II (DTG, RAL) & R263K (Intégrase) & Résistance DTG (rare, barrière haute). \\ \hline
IP (LPV/r, DRV/r) & M46I/L (Protéase) & Résistance IP (accumulation mutations). \\ \hline
\end{tabularx}
\end{center}

\vspace{0.5cm}
\textbf{QUESTIONS}

\textbf{Partie A : Dynamique de l'infection naturelle (Document 1)} \hfill (6 pts)
\begin{enumerate}
\item Nommer les trois phases (1, 2, 3) illustrées sur la courbe. Pour chaque phase, indiquer l'évolution de la charge virale (CV) et du taux de CD4. \hfill (3 pts)
\item Définir la \textbf{séroconversion}. À quelle période approximative (en mois) se situe-t-elle sur le graphique ? \hfill (1,5 pts)
\item Expliquer pourquoi le taux de CD4 finit par s'effondrer dans la phase 3 malgré la réponse immunitaire initiale. \hfill (1,5 pts)
\end{enumerate}

\textbf{Partie B : Cycle de réplication et Cibles thérapeutiques (Document 2)} \hfill (5 pts)
\begin{enumerate}
\item Identifier les enzymes virales agissant aux \textbf{Étapes X, Y, Z}. \hfill (1,5 pts)
\item Compléter le tableau ci-dessous en associant chaque enzyme à sa classe d'ARV correspondante et à un exemple de molécule (DCI) utilisée en 1ère ligne au Cameroun. \hfill (2,5 pts)
\begin{center}
\begin{tabular}{|c|c|c|}\hline
\textbf{Enzyme} & \textbf{Classe ARV (1ère ligne Cameroun)} & \textbf{Exemple DCI} \\ \hline
Transcriptase Inverse & & \\ \hline
Intégrase & & \\ \hline
Protéase & & \\ \hline
\end{tabular}
\end{center}
\item Pourquoi les Inhibiteurs de Protéase (IP) ne sont-ils pas en 1ère ligne recommandée au Cameroun actuellement ? \hfill (1 pt)
\end{enumerate}

\textbf{Partie C : Prise en charge thérapeutique et Résistance (Documents 3 \& 4)} \hfill (6 pts)
\begin{enumerate}
\item Un patient sous TLD (TDF+3TC+DTG) présente un échec virologique confirmé (CV $= 12\,000$ cp/mL à M12 et M15). Le génotypage révèle la mutation \textbf{M184V} sur la RT et \textbf{R263K} sur l'Intégrase.
\begin{itemize}
\item[a)] Interpréter l'impact de ces deux mutations sur l'efficacité du schéma TLD. \hfill (2 pts)
\item[b)] Proposer le schéma de 2ème ligne adapté, en justifiant le choix de chaque molécule par rapport aux mutations trouvées. \hfill (2 pts)
\end{itemize}
\item Une femme enceinte de 20 SA découvre sa séropositivité VIH-1 (CV $= 85\,000$ cp/mL, CD4 $= 420$/$\mu$L).
\begin{itemize}
\item[a)] Quel schéma ARV initier ? Quand ? \hfill (1 pt)
\item[b)] Détailler la prophylaxie du nouveau-né et le mode d'alimentation recommandés. \hfill (1 pt)
\end{itemize}
\end{enumerate}

\textbf{Partie D : Argumentation - Santé Publique (Connaissances + Documents)} \hfill (3 pts)
\begin{enumerate}
\item En vous appuyant sur les documents et vos connaissances, justifier en 10-15 lignes l'affirmation suivante : \textbf{"Atteindre les objectifs 95-95-95 de l'ONUSIDA au Cameroun nécessite de lever les freins au dépistage et à la rétention sous traitement, plus que d'attendre de nouveaux médicaments."}
\end{enumerate}
\end{exercice}

\begin{corrige}[Corrigé détaillé Sujet BEPC 2025]
\textbf{Partie A : Dynamique de l'infection naturelle}
\begin{enumerate}
\item \textbf{Phase 1 (Infection primaire / Aiguë) :} Mois 0-3. CV : Pic très élevé ($> 10^6$ cp/mL), puis chute brutale. CD4 : Chute brutale transitoire ($< 500$), puis remontée partielle.
\textbf{Phase 2 (Phase chronique / Latence clinique) :} Mois 3 - $\sim$40. CV : Plateau "set-point" ($\sim 10^4 - 10^5$ cp/mL), relativement stable. CD4 : Décline lente et progressive ($\sim 50-80$/an).
\textbf{Phase 3 (Phase Sida / Symptomatique) :} Mois $> 40$. CV : Remontée exponentielle ($> 10^5$ cp/mL). CD4 : Chute accélérée, franchissement seuil $< 200$/$\mu$L. Infections opportunistes.
\item \textbf{Séroconversion :} Apparition des anticorps anti-VIH détectables dans le sang. Période : \textbf{2 à 12 semaines} (mois 0,5 à 3) après infection. Sur le graphique : pendant la chute de la CV du pic vers le set-point (fin Phase 1).
\item \textbf{Effondrement CD4 Phase 3 :} 1) Destruction directe des LT4 infectés par le virus (lyse, syncytia). 2) Activation immunitaire chronique $\to$ apoptose des cellules "bystanders" (non infectées). 3) Épuisement thymique (production nouvelles cellules insuffisante). 4) Fibrose tissu lymphoïde $\to$ perte niches de survie. 5) Variants viraux plus cytopathiques (tropisme CXCR4).
\end{enumerate}

\textbf{Partie B : Cycle et Cibles}
\begin{enumerate}
\item \textbf{Étape X :} Transcriptase Inverse (RT). \textbf{Étape Y :} Intégrase. \textbf{Étape Z :} Protéase.
\item \begin{center}
\begin{tabular}{|c|c|c|}\hline
\textbf{Enzyme} & \textbf{Classe ARV (1ère ligne Cameroun)} & \textbf{Exemple DCI} \\ \hline
Transcriptase Inverse & IEP (Inhibiteurs Nucléosidiques) & \textbf{TDF} (Ténofovir) ou \textbf{3TC} (Lamivudine) \\ \hline
Intégrase & II (Inhibiteurs d'Intégrase) & \textbf{DTG} (Dolutégravir) \\ \hline
Protéase & \textit{Pas en 1ère ligne} & (LPV/r, DRV/r en 2ème ligne) \\ \hline
\end{tabular}
\end{center}
\item IP : Barrière génétique plus faible que DTG (besoin de booster ritonavir/cobicistat $\to$ interactions, effets métaboliques dyslipidémie, insulinorésistance), posologie 2 fois/jour (observance), coût plus élevé. DTG : Barrière haute, 1 cp/jour, bien toléré, peu d'interactions, coût bas (générique).
\end{enumerate}

\textbf{Partie C : Prise en charge et Résistance}
\begin{enumerate}
\item[a)] \textbf{M184V (RT) :} Résistance haute à la \textbf{3TC} (et FTC). \textbf{Mais :} Hypersensibilise au TDF (et AZT) $\to$ TDF reste partiellement actif. \textbf{R263K (Intégrase) :} Mutation majeure de résistance au \textbf{DTG} (et BIC, CAB). \textbf{Conclusion :} Le schéma TLD est \textbf{compromis} : 3TC inactif, DTG résistant (R263K), seul TDF conserve une activité (renforcée par M184V). Échec virologique confirmé.
\item[b)] \textbf{Schéma 2ème ligne recommandé :} \textbf{AZT + 3TC + DRV/r (boosté)}.
\begin{itemize}
\item \textbf{AZT (Zidovudine) :} IEP. Actif malgré M184V (pas de résistance croisée M184V $\to$ AZT). Synergie M184V+AZT.
\item \textbf{3TC (Lamivudine) :} Maintenue malgré M184V pour \textbf{conserver la mutation M184V} (coût fitness viral, hypersensibilisation AZT/TDF). "Effet M184V".
\item \textbf{DRV/r (Darunavir/ritonavir) :} IP de 2ème génération. Barrière génétique très élevée. Aucune mutation Protéase détectée $\to$ pleinement actif. Pris 2 fois/jour.
\end{itemize}
\item[a)] \textbf{TLD (TDF+3TC+DTG) immédiatement} (Test \& Treat). Grossesse = urgence PMTCT. DTG sûr au 1er trimestre (données Botswana, surveillance : pas d'excès défauts tube neural).
\item[b)] \textbf{Prophylaxie nourrisson :} \textbf{AZT + NVP pendant 6 semaines} (schéma combiné car mère CV $> 1000$ cp/mL à l'accouchement probable si début ARV tardif 20 SA, ou par prudence). \textbf{Allaitement :} \textbf{Allaitement maternel exclusif 6 mois} + poursuite allaitement jusqu'à 12-24 mois. \textbf{Condition sine qua non :} Mère sous ARV avec \textbf{CV indétectable ($< 50$ cp/mL)} à l'accouchement et pendant l'allaitement.
\end{enumerate}

\textbf{Partie D : Argumentation}
\textbf{Proposition de rédaction :}
"Atteindre les 95-95-95 au Cameroun (90-85-90 en 2023) bute moins sur l'arsenal thérapeutique (DTG, TLD, génériques gratuits, efficaces, barrière haute) que sur le \textbf{continuum de soins}. Le 1er 95 (dépistage) stagne à 90\% : freins = stigmatisation (peur du positif), faible perception du risque (hommes, jeunes), accès géographique (zones rurales), absence d'autotest généralisé. Le 2ème 95 (liaison aux soins/rétention) est le maillon faible (85\%) : pertes de vue (déménagements, effets secondaires mal gérés, coûts indirects transport, stigmatisation au travail/famille, rupture stock ponctuelle). Or, \textbf{un patient non dépisté transmet (CV haute) ; un patient dépisté non traité transmet ; un patient traité mal observant développe des résistances (coût 2ème ligne $\times$10)}. Investir dans \textbf{pairs-éducateurs, accueil bienveillant, dispensation multi-mois (MMD 3-6 mois), horaires élargis, intégration VIH/SR/PN, lutte anti-stigma} a un \textbf{coût-efficacité supérieur} à l'attente d'hypothétiques nouveaux médicaments (guérison, vaccin incertains). Les outils actuels (TLD, PrEP, PMTCT, CV) \textbf{permettent déjà la fin de l'épidémie SI ils atteignent tout le monde}."
\end{corrige}

\courssub{Bilan des compétences visées (Auto-évaluation)}

\begin{center}
\begin{tabularx}{\linewidth}{|l|X|c|}\hline
\textbf{Compétence (Programme 3ème)} & \textbf{Indicateur de réussite (Je sais...)} & \textbf{Validation} \\ \hline
Connaître structure VIH, cycle, cibles & Dessiner/annoter cycle ; Citer 3 enzymes + 3 classes ARV + 1 molécule/classe & $\square$ \\ \hline
Expliquer dynamique CV/CD4 & Décrire 3 phases ; Situer séroconversion ; Expliquer effondrement CD4 & $\square$ \\ \hline
Interpréter sérologie (Algo 3 tests) & Conclure statut (Pos/Nég, VIH-1/2) à partir de 3 résultats & $\square$ \\ \hline
Stadifier (OMS) & Associer signes cliniques $\to$ Stade 1-4 ; Seuil CD4 Sida ($< 200$) & $\square$ \\ \hline
Proposer prise en charge ARV & Choisir 1ère ligne (TLD), 2ème ligne (AZT+3TC+DRV/r) selon mutations & $\square$ \\ \hline
Détailler PMTCT & Schéma mère, prophylaxie nourrisson, allaitement, calendrier dépistage enfant & $\square$ \\ \hline
Analyser résistance (génotypage) & Identifier mutation $\to$ classe touchée $\to$ prédire résistance croisée $\to$ adapter schéma & $\square$ \\ \hline
Argumenter santé publique & Structurer thèse/arguments (individuel, collectif, éco, droits) $\to$ conclusion & $\square$ \\ \hline
Communiquer prévention (Pairs) & Citer 3 voies transmission, 5 non-voies, droits patient, rôle pair-éducateur & $\square$ \\ \hline
\end{tabularx}
\end{center}

\begin{attention}[Derniers conseils pour l'épreuve du BEPC]
\begin{itemize}
\item \textbf{Lis TOUS les documents} avant d'écrire. Les réponses sont souvent \textbf{dans les documents} (valeurs, noms, schémas).
\item \textbf{Croise Doc + Connaissance :} "D'après le document 1, la CV est de... Mes connaissances m'indiquent que cela correspond à la phase..."
\item \textbf{Vocabulaire précis :} "Charge virale", "Copies/mL", "Lymphocytes T CD4+", "Provirus", "Intégrase", "Indétectable", "Échec virologique", "Observance", "Résistance croisée". Pas de "virus dans le sang", "défenses", "médicament fort".
\item \textbf{Unités obligatoires :} CD4 en cellules/$\mu$L (ou /mm$^3$), CV en copies/mL, Temps en mois/années.
\item \textbf{Soigne la rédaction :} Phrases complètes, connecteurs logiques (Donc, Or, En effet, Par conséquent), paragraphes aérés.
\end{itemize}
\end{attention}

\begin{savaistu}[Ton atout pour le BEPC : La rigueur scientifique]
Le jury BEPC valorise les copies qui montrent une \textbf{démarche scientifique} : Observation (Doc) $\to$ Hypothèse/Interprétation $\to$ Connaissance (Cours) $\to$ Conclusion argumentée. Ne récite pas le cours : \textbf{utilise-le pour expliquer le document}. C'est la différence entre une note moyenne et une excellente note. Courage, tu es prêt(e) ! \textbf{Le VIH se contrôle, le BEPC se réussit.}
\end{savaistu}

\newpage

\coursec{Activité d'intégration}

\courssub{Objectifs de l'activité}
Cette activité d'intégration vise à mobiliser l'ensemble des savoirs et savoir-faire de la leçon sur le VIH/Sida à travers une situation clinique réaliste, contextualisée au Cameroun. Elle permet à l'élève de :
\begin{itemize}
    \item Interpréter un algorithme de sérologie VIH et comprendre la notion de fenêtre sérologique.
    \item Stadifier l'infection selon la classification OMS (immunologique et clinique) à partir de numérations CD4.
    \item Analyser la dynamique rétrospective de la charge virale (CV) et des lymphocytes T CD4 pour estimer la date de contamination et identifier le \emph{set point} viral.
    \item Justifier une prescription antirétrovirale (ARV) de première ligne recommandée par l'OMS et le Programme National de Lutte contre le Sida (PNLS) au Cameroun, en tenant compte de la grossesse de la conjointe (sécurité du Dolutégravir).
    \item Proposer une prise en charge globale du couple : Prophylaxie Pré-Exposition (PrEP), Prophylaxie Post-Exposition (PEP), Prévention de la Transmission Mère-Enfant (PTME), Traitement comme Prévention (TasP), observance et éducation thérapeutique.
    \item Rédiger une note de synthèse professionnelle, claire et argumentée, adaptée à un personnel soignant (infirmier(ère) de CSI).
\end{itemize}

\courssub{Situation : Le dossier clinique de Monsieur K.}
\textbf{Contexte :} Nous sommes au Centre de Santé Intégré (CSI) de \textbf{Nkolbisson}, arrondissement de Yaoundé 6. Monsieur K., 32 ans, maçon, se présente pour une consultation de suivi après un dépistage VIH réalisé lors d'une campagne « \emph{Holidays Free of HIV} » organisée par le Ministère de la Santé Publique (MINSANTE).

\textbf{1. Résultats de la sérologie VIH (Algorithme national 3 tests rapides) :}
\begin{center}
\begin{tabularx}{\linewidth}{|l|X|c|}\hline
\textbf{Test} & \textbf{Type / Marque (Exemple)} & \textbf{Résultat} \\ \hline
Test 1 (Dépistage) & Determine\texttrademark\ HIV-1/2 (Abbott) & \textbf{Réactif (+)} \\ \hline
Test 2 (Confirmation) & Uni-Gold\texttrademark\ HIV (Trinity Biotech) & \textbf{Réactif (+)} \\ \hline
Test 3 (Arbitrage) & SD Bioline HIV-1/2 3.0 (Standard Diagnostics) & \textbf{Réactif (+)} \\ \hline
\end{tabularx}
\end{center}
\emph{Interprétation du laboratoire :} « Séropositivité VIH-1 confirmée selon l'algorithme national. »

\textbf{2. Bilan biologique initial (J0 - Avant mise sous ARV) :}
\begin{center}
\begin{tabular}{|l|c|c|}\hline
\textbf{Paramètre} & \textbf{Résultat} & \textbf{Valeurs de référence / Seuil décisionnel} \\ \hline
Charge Virale (CV) VIH-1 (ARN) & \SI{450 000}{copies/mL} & $<$ \SI{50}{copies/mL} (Indétectable) \\ \hline
Numération CD4 absolus & \SI{180}{cellules/\micro\liter} & \SIrange{500}{1500}{cellules/\micro\liter} \\ \hline
Pourcentage CD4 & \textbf{12\%} & $>$ 29\% \\ \hline
Hémoglobine & \SI{11,0}{g/dL} & \SIrange{13}{17}{g/dL} (Anémie modérée) \\ \hline
Créatininémie & \SI{9,5}{mg/L} (\SI{84}{\micro mol/L}) & \SIrange{6}{12}{mg/L} (Normal) \\ \hline
Débit de filtration glomérulaire (DFG - MDRD) & \SI{95}{mL/min/1,73m^2} & $>$ 60 (Normal) \\ \hline
Grossesse conjointe & \textbf{Oui, 12 SA (3 mois)} & --- \\ \hline
\end{tabular}
\end{center}

\textbf{3. Courbes rétrospectives estimées (Modélisation clinique sur 3 ans) :}
Le clinicien, à partir de l'anamnèse (relations non protégées il y a ~3 ans, épisode fébrile aigu il y a 34 mois), reconstitue l'évolution probable des marqueurs virologiques et immunologiques.

\begin{center}
\textit{Figure 1 : Modélisation rétrospective de la dynamique VIH chez M. K. sur 36 mois. Haut : Charge virale (échelle logarithmique). Bas : Numération CD4 absolus (échelle linéaire). Les zones colorées délimitent les grandes phases physiopathologiques.}
\end{center}

\textbf{4. Proposition thérapeutique (Ordonnance médicale du Médecin traitant du CSI) :}
\begin{itemize}
    \item \textbf{TDF 300 mg} (Ténofovir Disoproxil Fumarate) : 1 cp/jour.
    \item \textbf{3TC 300 mg} (Lamivudine) : 1 cp/jour.
    \item \textbf{DTG 50 mg} (Dolutégravir) : 1 cp/jour.
    \item \textbf{Association fixe TDF/3TC/DTG (TLD)} : 1 cp/jour le matin (recommandation PNLS 2023+).
    \item \textbf{Cotrimoxazole 960 mg} (Sulfaméthoxazole/Triméthoprime) : 1 cp/jour (Prophylaxie primaire des Infections Opportunistes).
    \item \textbf{Acide folique 5 mg} : 1 cp/jour (contexte grossesse conjointe + Cotrimoxazole).
\end{itemize}

\textbf{5. Document d'information remis au couple :} Fiche « \emph{Vivre avec le VIH au Cameroun} » : PrEP, PEP, PTME, Allaitement, TasP, Effets indésirables DTG (insomnie, prise de poids, hyperglycémie), Importance de l'observance ($>$ 95\%).

\courssub{Consignes}
À partir du dossier ci-dessus et de vos connaissances, répondez aux questions suivantes en justifiant chaque réponse par des éléments du dossier ou du cours.

\textbf{Tâche A : Interprétation de la sérologie}
\begin{enumerate}
    \item Quel est le statut sérologique définitif de M. K. ? Justifiez en citant l'algorithme utilisé.
    \item Expliquez pourquoi trois tests rapides de marques différentes sont nécessaires.
    \item M. K. rapporte une prise de risque unique il y a 3 semaines. Le résultat serait-il fiable ? Expliquez la notion de \emph{fenêtre sérologique} (période d'éclipse sérologique) pour les tests rapides d'anticorps.
\end{enumerate}

\textbf{Tâche B : Stadification de la maladie (Classifications OMS et CDC)}
\begin{enumerate}
    \item En vous basant \textbf{uniquement} sur le taux de CD4 (180 cellules/µL, 12\%) et l'absence de signes cliniques déclarés (asymptomatique), classez M. K. selon la stadification \textbf{immunologique OMS} (Stade 1 à 4) et la classification \textbf{CDC (Catégorie A, B, C)}.
    \item Justifiez pourquoi le seuil de \SI{200}{cellules/\micro\liter} est un tournant majeur dans la prise en charge (pronostic, prophylaxie IO).
\end{enumerate}

\textbf{Tâche C : Analyse de la dynamique rétrospective (Figure 1)}
\begin{enumerate}
    \item Estimez la date probable de l'infection primaire de M. K. (repère temporel sur l'axe des abscisses).
    \item Identifiez sur la courbe de charge virale la phase de \emph{set point} viral. Donnez sa valeur approximative en copies/mL et en log$_{10}$. Quelle est son importance pronostique ?
    \item Décrivez l'évolution des CD4 entre le mois 24 et le mois 36. Calculez la pente moyenne de perte (cellules/µL/an) sur cette période et comparez-la à la perte physiologique moyenne en phase de latence clinique ($\sim$ 50-80 cellules/µL/an). Quelle conclusion tirez-vous sur l'urgence thérapeutique ?
\end{enumerate}

\textbf{Tâche D : Justification de l'ordonnance ARV et de la prophylaxie}
\begin{enumerate}
    \item Pour chaque molécule de l'association fixe TLD (TDF, 3TC, DTG), citez sa \textbf{classe thérapeutique}, sa \textbf{cible précise dans le cycle de réplication du VIH} (enzyme virale ou étape) et précisez s'il s'agit d'une \textbf{première ligne (L1)} ou deuxième ligne (L2) selon l'OMS/PNLS Cameroun.
    \item Justifiez spécifiquement le choix du \textbf{Dolutégravir (DTG)} au vu de la grossesse de la conjointe (12 SA) et du désir de grossesse du couple. Citez les données de sécurité (études Tsepamo, ODYSSEY).
    \item Justifiez la prescription de \textbf{Cotrimoxazole} (CTX) à M. K. en vous référant au seuil de CD4 et aux infections opportunistes ciblées (Pneumocystose, Toxoplasmose, Infections bactériennes, Paludisme).
\end{enumerate}

\textbf{Tâche E : Conseil au couple : Prévention combinée et PTME}
\begin{enumerate}
    \item \textbf{PTME (Conjointe) :} La conjointe est séronégative (testé négatif à la première consultation prénatale - CPN1). Quelle stratégie de prévention de la transmission mère-enfant (PTME) proposer pour \textbf{elle} ? (PrEP ?). Quelle stratégie pour \textbf{l'enfant à naître} ?
    \item \textbf{TasP (M. K.) :} Expliquez le concept \emph{Undetectable = Untransmittable} (I = I). Quel est le délai estimé pour atteindre l'indétectabilité sous TLD ? Quelle conduite tenir pour les rapports sexuels pendant ce délai ?
    \item \textbf{Allaitement :} La conjointe souhaite allaiter. Quelle est la recommandation officielle OMS/PNLS Cameroun 2024 pour une mère séronégative dont le partenaire est sous ARV efficace ? Et si la mère était séropositive ?
\end{enumerate}

\textbf{Tâche F : Rédaction d'une note de synthèse pour l'infirmier(ère) du CSI}
Rédigez une note professionnelle (style « Note de service / Transfert de dossier ») adressée à l'infirmier(ère) référent(e) VIH du CSI, résumant :
\begin{itemize}
    \item Le diagnostic, le stade, l'urgence.
    \item Le traitement prescrit (TLD + CTX) et sa posologie.
    \item Les points clés d'\textbf{éducation thérapeutique} à renforcer lors de la prochaine visite (J7, J30, M3) : Observance (objectif $>$ 95\%), Prise le matin pour DTG (insomnie), Surveillance prise de poids / glycémie (DTG), Tolérance digestive (TDF/3TC), Importance de ne pas arrêter sans avis médical.
    \item Le plan de suivi biologique (CV à M3, M6 puis annuel ; CD4 à M6 puis si CV indétectable ; Créatinine/Hb à M3, M6).
    \item La prise en charge de la conjointe (PrEP, dépistage répété, PTME).
\end{itemize}

\courssub{Ressources à mobiliser (Rappels de cours)}
\begin{definition}[Algorithme de dépistage VIH Cameroun]
Stratégie 3 tests rapides séquentiels : Test 1 (Haute Sensibilité) $\rightarrow$ Si (+) Test 2 (Haute Spécificité) $\rightarrow$ Si (+) Test 3 (Arbitrage). 3 tests (+) = Séropositif. 1er (+), 2ème (-), 3ème (-) = Séronégatif. Discordance = Envoi au labo de référence (ELISA/Western Blot/Charge Virale).
\end{definition}

\begin{definition}[Stadification Immunologique OMS (Adulte/Adolescent $>$ 5 ans)]
\begin{itemize}
    \item Stade 1 : CD4 $>$ 500 / $\mu$L ($>$ 29\%)
    \item Stade 2 : CD4 350 - 499 / $\mu$L (19-28\%)
    \item Stade 3 : CD4 200 - 349 / $\mu$L (14-18\%)
    \item Stade 4 (Sida) : CD4 $<$ 200 / $\mu$L ($<$ 14\%) \textbf{OU} Infection Opportuniste définissant le Sida.
\end{itemize}
\end{definition}

\begin{propriete}[Cycle de réplication du VIH \& Cibles ARV]
\begin{itemize}
    \item \textbf{Entrée / Fusion :} Enfuvirtide (T20) - Inhibiteur de fusion (peu utilisé en 1ère ligne).
    \item \textbf{Transcription Inverse (ARN $\rightarrow$ ADN) :} \textbf{INTI} (TDF, 3TC, AZT, ABC) - Inhibiteurs Nucléos(t)idiques ; \textbf{INNTI} (EFV, NVP, DOR, RPV) - Inhibiteurs Non Nucléosidiques.
    \item \textbf{Intégration (ADN viral $\rightarrow$ ADN hôte) :} \textbf{II} (DTG, RAL, BIC, CAB) - Inhibiteurs de l'Intégrase. \textbf{Cible de 1ère ligne privilégiée (Barrière génétique haute)}.
    \item \textbf{Maturation (Clivage polyprotéines) :} \textbf{IP} (ATV/r, DRV/r, LPV/r) - Inhibiteurs de Protéase (2ème ligne).
\end{itemize}
\end{propriete}

\begin{aretenir}[DTG et Grossesse (Données Tsepamo / Botswana / OMS 2019-2024)]
Le Dolutégravir (DTG) est \textbf{recommandé en 1ère ligne chez la femme en âge de procréer et pendant la grossesse} (y compris 1er trimestre). Le risque de défauts du tube neural (DTN) initialement signalé (0,9\% vs 0,1\%) n'a pas été confirmé dans les vastes cohortes ultérieures (Tsepamo mise à jour 2021, DolPHIN-2, IMPAACT 2010). DTG offre une suppression virale plus rapide, meilleure tolérance, barrière génétique élevée. \textbf{Prise le matin conseillée si insomnie.}
\end{aretenir}

\begin{aretenir}[Indications Cotrimoxazole (CTX) Prophylaxie Primaire IO - OMS/PNLS]
\begin{itemize}
    \item \textbf{CD4 $<$ 350 cellules/$\mu$L} (ou Stade OMS 3/4 clinique).
    \item \textbf{Arrêt possible si CD4 $>$ 350 cellules/$\mu$L} durablement (ex: 6 mois) sous ARV efficace.
    \item Cible : \emph{Pneumocystis jirovecii} (Pneumocystose), \emph{Toxoplasma gondii}, Bactéries encapsulées, \emph{Plasmodium} (Paludisme), Isosporose.
\end{itemize}
\end{aretenir}

\begin{aretenir}[PTME / PrEP / TasP - Contexte Cameroun]
\begin{itemize}
    \item \textbf{PrEP (Prophylaxie Pré-Exposition) :} TDF/3TC (ou TDF/FTC) 1 cp/jour. Indiquée pour partenaire séronégatif(ve) à risque élevé (conjoint séropositif non indétectable, travailleur(se) du sexe, HSH). \textbf{Efficacité $>$ 99\% si observance parfaite}.
    \item \textbf{TasP (Treatment as Prevention) :} « I = I » (Indétectable = Intransmissible). CV $<$ 50-200 copies/mL maintenue 6 mois = Risque transmission sexuelle \textbf{nul}.
    \item \textbf{Allaitement (Mère séronégative, Père séropositif sous ARV) :} Allaitement maternel exclusif 6 mois \textbf{recommandé} (Risque transmission par lait négligeable si père indétectable + PrEP mère possible). Si Mère séropositive : Allaitement maternel exclusif 6 mois \textbf{sous ARV mère + ARV nourrisson (NVP 6-12 sem)} recommandé au Cameroun (Option B+).
\end{itemize}
\end{aretenir}

\courssub{Démarche et corrigé commenté}
\begin{exoresolu}[Corrigé détaillé de l'activité d'intégration - M. K.]
\noindent\textbf{Tâche A : Interprétation de la sérologie}
\begin{enumerate}
    \item \textbf{Statut :} M. K. est \textbf{séropositif pour le VIH-1} (infection confirmée).
    \item \textbf{Justification algorithme :} L'algorithme national camerounais (conforme OMS) utilise 3 tests rapides de marques différentes (Différents antigènes/épitopes ciblés pour éviter faux positifs croisés).
    \begin{itemize}
        \item Test 1 (Determine\texttrademark) : Réactif (+) $\rightarrow$ Dépistage.
        \item Test 2 (Uni-Gold\texttrademark) : Réactif (+) $\rightarrow$ Confirmation.
        \item Test 3 (SD Bioline\texttrademark) : Réactif (+) $\rightarrow$ Arbitrage (Confirmation définitive).
    \end{itemize}
    Trois tests réactifs (+ + +) valident la séropositivité. Pas besoin de Western Blot ou CV pour le diagnostic (sauf nourrisson $<$ 18 mois).
    \item \textbf{Fenêtre sérologique :} Si la prise de risque date de 3 semaines (21 jours), le résultat \textbf{N'EST PAS FIABLE} pour exclure l'infection.
    \begin{itemize}
        \item Les tests rapides détectent les \textbf{anticorps anti-VIH} (Ac anti-p24, gp41, gp120/gp160).
        \item La \emph{fenêtre sérologique} (période entre infection et détection des Ac) est en moyenne de \textbf{3 à 6 semaines} (jusqu'à 3 mois rares) pour les tests de 3ème/4ème génération.
        \item À 3 semaines (phase aiguë), M. K. serait probablement \textbf{Antigène p24 positif / ARN VIH positif (CV très haute) mais Anticorps NÉGATIFS ou INDÉTERMINÉS}.
        \item \textbf{Conduite :} Refaire une sérologie à 6 semaines (42 jours) et 3 mois, ou mieux : doser la \textbf{Charge Virale (ARN VIH)} ou l'\textbf{Antigène p24} (Test 4ème génération Ag/Ac) immédiatement.
    \end{itemize}
\end{enumerate}

\noindent\textbf{Tâche B : Stadification}
\begin{enumerate}
    \item \textbf{Stade Immunologique OMS : Stade 4 (Sida).}
    
    Justification : CD4 = \SI{180}{cellules/\micro\liter} ($<$ \SI{200}{cellules/\micro\liter}) \textbf{ET} Pourcentage CD4 = 12\% ($<$ 14\%). Le seuil des 200 cellules/$\mu$L (ou 14\%) définit le Stade 4 (Sida) immunologique, \textbf{même en l'absence de signes cliniques} (Stade clinique 1 = Asymptomatique). Le patient est donc \textbf{Stade Clinique 1 / Stade Immunologique 4}.
    
    \item \textbf{Classification CDC : Catégorie A3 (ou B3/C3 si IO passée/actuelle).}
    
    CDC utilise : Lettre (A=Asymptomatique/Adénopathies, B=Symptomatique modéré, C=Sida définitif) + Chiffre (1=$>$500, 2=200-499, 3=$<$200).
    Ici : Asymptomatique (A) + CD4 $<$ 200 (3) $\rightarrow$ \textbf{A3}.
    
    \item \textbf{Importance du seuil 200 cellules/$\mu$L :}
    \begin{itemize}
        \item \textbf{Pronostic :} Risque exponentiel d'Infections Opportunistes (IO) graves : Pneumocystose (\emph{P. jirovecii}), Toxoplasmose cérébrale, Cryptococcose, Cytomégalovirus, Mycobactérioses atypiques.
        \item \textbf{Prophylaxie :} Déclenche \textbf{OBLIGATOIRE} de la prophylaxie primaire par Cotrimoxazole (CTX) (arrêt seulement si CD4 $>$ 350 durable sous ARV).
        \item \textbf{Urgence ARV :} Mise sous ARV \textbf{immédiate (J0 ou J7 max)} recommandée (Guidelines OMS « Treat All » + urgence clinique Stade 4). Risque de mortalité élevé sans ARV rapide (syndrome de reconstitution immunitaire - SRI - à surveiller).
    \end{itemize}
\end{enumerate}

\noindent\textbf{Tâche C : Analyse dynamique (Figure 1)}
\begin{enumerate}
    \item \textbf{Date infection : Mois 0 (il y a 36 mois, soit 3 ans).} Pic de CV à \SI{5e6}{copies/mL} (5 millions) au mois 1 et chute CD4 brutale (850 $\rightarrow$ 450) au mois 3 = \textbf{Syndrome rétroviral aigu (Infection primaire)}.
    \item \textbf{Set Point Viral :} Période \textbf{Mois 6 à Mois 24 (18 mois)}. CV stable autour de \textbf{\num{3e4} à \SI{6e4}{copies/mL}} ($\approx$ \textbf{4,5 - 4,8 log$_{10}$} copies/mL). CD4 stable autour de \textbf{500-550 cellules/$\mu$L} (perte lente $\sim$ 50/an).
    
    \textbf{Importance :} Le \emph{set point} reflète l'équilibre entre réplication virale et réponse immune (CTL CD8+). Un set point bas ($<$ 10 000 copies/mL = 4 log) = « Contrôleur élite » ou bon pronostic (lente progression). Un set point haut ($>$ 100 000 = 5 log) = progression rapide vers Sida. Ici set point $\sim$ 4,7 log = progression \textbf{intermédiaire/modérée}.
    
    \item \textbf{Évolution Mois 24-36 (12 derniers mois) :}
    \begin{align*}
    \text{CD4 Mois 24} &= 500 \\
    \text{CD4 Mois 36} &= 180 \\
    \text{Perte totale} &= 320 \text{ cellules/$\mu$L en 12 mois} \\
    \text{Pente moyenne} &= \mathbf{320 \text{ cellules/$\mu$L/an}}
    \end{align*}
    
    \textbf{Comparaison :} Perte physiologique latence $\approx$ 50-80 cellules/$\mu$L/an. Perte observée \textbf{4 à 6 fois plus rapide}.
    \textbf{Interprétation :} Rupture du contrôle immunitaire (épuisement CD8+, destruction architecture ganglionnaire, activation immune chronique). Entrée en \textbf{phase de défaillance immunitaire accélérée (Pré-Sida)}.
    \textbf{Conclusion :} \textbf{Urgence thérapeutique absolue.} Le patient a perdu en 1 an ce qu'il aurait dû perdre en 4-6 ans. Risque imminent d'IO (CD4 $<$ 200). ARV doit démarrer \textbf{cette semaine}.
    
\end{enumerate}

\noindent\textbf{Tâche D : Justification Ordonnance}
\begin{enumerate}
    \item \textbf{Analyse TLD (TDF + 3TC + DTG) - Première Ligne (L1) OMS/PNLS :}
    \begin{center}
    \begin{tabular}{|l|l|l|l|}\hline
    \textbf{Molécule} & \textbf{Classe} & \textbf{Cible Cycle VIH} & \textbf{Ligne OMS} \\ \hline
    TDF (Ténofovir) & INTI (Inhibiteur Nucléotidique Transcriptase Inverse) & \textbf{Transcriptase Inverse} (Incorporation competitive, terminaison chaîne ADN viral) & L1 (Pilier) \\ \hline
    3TC (Lamivudine) & INTI (Inhibiteur Nucléosidique Transcriptase Inverse) & \textbf{Transcriptase Inverse} (Analogue cytidine, terminaison chaîne) & L1 (Pilier) \\ \hline
    DTG (Dolutégravir) & \textbf{II (Inhibiteur Intégrase)} - \emph{INSTI} & \textbf{Intégrase virale} (Blocage transfert brin ADN viral dans ADN hôte) & \textbf{L1 Privilégiée} \\ \hline
    \end{tabular}
    \end{center}
    \textbf{Note :} TDF/3TC = « Backbone » (colonne vertébrale) NRTI. DTG = « Anchor » (ancrage) INSTI. Barrière génétique DTG très haute (mutations multiples nécessaires pour résistance).
    \item \textbf{Choix DTG \& Grossesse Conjointe (12 SA) :}
    \begin{itemize}
        \item \textbf{Sécurité 1er trimestre :} Données \textbf{Tsepamo (Botswana, $>$ 120 000 grossesses)} mise à jour 2021+ : Taux DTN (Défauts Tube Neural) sous DTG periconceptionnel = \textbf{0,3\%} vs 0,1\% sous non-DTG. Différence non significative statistiquement, cliniquement rassurante. OMS 2019 / PNLS Cameroun 2020+ : \textbf{DTG recommandé 1ère ligne TOUTES femmes en âge de procréer y compris 1er trimestre}.
        \item \textbf{Avantages majeurs :} Suppression virale plus rapide (CV $<$ 50 copies/mL à 4-12 sem vs 24 sem EFV/NVP) $\rightarrow$ Protège le fœtus \textbf{plus tôt} (TasP in utero). Meilleure tolérance (moins troubles neuropsychiatriques vs EFV). Barrière génétique haute $\rightarrow$ Moins de résistance transmise à l'enfant en cas d'échec.
        \item \textbf{Supplémentation :} Acide folique 5 mg/j (ici prescrit) + Vitamine D/Calcium si TDF long terme (ostéoporose).
    \end{itemize}
    \item \textbf{Justification Cotrimoxazole (CTX 960 mg/j) :}
    \begin{itemize}
        \item \textbf{Indication formelle :} CD4 = \SI{180}{cellules/\micro\liter} $<$ \textbf{350 cellules/$\mu$L} (Seuil OMS/PNLS).
        \item \textbf{Cibles :} \emph{Pneumocystis jirovecii} (Pneumonie à Pneumocystes - PPJ, 1ère cause décès Sida non traité), \emph{Toxoplasma gondii} (Toxoplasmose cérébrale), Infections bactériennes graves (Pneumocoque, Haemophilus, Salmonella - bactériémies), \emph{Isospora belli}, \textbf{Paludisme} (Efficacité prouvée prévention accès palustres en zone endémique Cameroun).
        \item \textbf{Durée :} Maintenu tant que CD4 $<$ 350 cellules/$\mu$L. Réévaluer à M6/M12 (si CD4 $>$ 350 durable + CV indétectable $\rightarrow$ arrêt possible).
        \item \textbf{Surveillance :} Allergie sulfa (éruption cutanée, syndrome de Lyell/Stevens-Johnson rare), Hyperkaliémie, Leucopénie (ajout AZT), Crétinine (rénal).
    \end{itemize}
\end{enumerate}

\noindent\textbf{Tâche E : Conseil au couple}
\begin{enumerate}
    \item \textbf{PTME (Conjointe Séronégative - Partenaire Séropostif) :}
    \begin{itemize}
        \item \textbf{Pour la conjointe (Mère) :} \textbf{PrEP (Prophylaxie Pré-Exposition) recommandée FORTEMENT} pendant toute la grossesse et l'allaitement (Option B+ adaptée au partenaire). Schéma : \textbf{TDF 300 mg + 3TC 300 mg (ou FTC 200 mg) 1 cp/jour}. Début immédiat (aujourd'hui). Suivi rénal (Créatinine) trimestriel. Dépistage VIH répété à chaque CPN (mensuel) + à l'accouchement + 6 sem post-partum.
        \item \textbf{Pour l'enfant :} \textbf{Pas de prophylaxie ARV systématique à la naissance} (zidovudine/NVP) \textbf{SI} le père est sous ARV efficace (CV indétectable visée) \textbf{ET} la mère sous PrEP \textbf{ET} allaitement exclusif. \textbf{MAIS} : Au Cameroun, par prudence (risque rupture observance père, fenêtre sérologique mère), on prescrit souvent \textbf{NVP (Névirapine) sirop 1,5 mL/kg/j (ou 10-15 mg/j) pendant 6 semaines} à l'enfant (Prophylaxie « exposition au VIH »). \textbf{Décision collégiale médecin/sage-femme.}
    \end{itemize}
    \item \textbf{TasP (M. K.) :}
    \begin{itemize}
        \item \textbf{Concept I = I :} \emph{Indétectable = Intransmissible}. Preuves : PARTNER 1\&2, Opposites Attract, HPTN 052. 0 transmission liée dans couples sérodifférents (hétéro/homo) avec CV $<$ 200 copies/mL (souvent $<$ 50) maintenue $>$ 6 mois.
        \item \textbf{Délai indétectabilité sous TLD :} \textbf{4 à 12 semaines} (DTG + TDF/3TC très puissant). CV $<$ 50 copies/mL attendu à \textbf{Mois 3 (M3)}.
        \item \textbf{Conduite pendant délai (J0 $\rightarrow$ M3) :} \textbf{Préservatif (masculin/féminin) OBLIGATOIRE à chaque rapport.} PrEP conjointe = double protection (Ceinture + Bretelles). Éducation sur « Charge virale $\neq$ Guérison ».
    \end{itemize}
    \item \textbf{Allaitement :}
    \begin{itemize}
        \item \textbf{Mère Séronégative (Cas présent) :} \textbf{Allaitement maternel exclusif 6 mois RECOMMANDÉ} (OMS/PNLS). Risque transmission par lait maternel \textbf{nul} (mère non infectée). PrEP mère = sécurité maximale. Avantages nutritionnels/immunologiques majeurs pour l'enfant (mortalité infanto-juvénile).
        \item \textbf{Si Mère Séropostive (Hypothétique) :} \textbf{Allaitement maternel exclusif 6 mois AUSSI RECOMMANDÉ} au Cameroun (Option B+). Condition : Mère sous ARV efficace (CV indétectable) \textbf{ET} Enfant sous Prophylaxie ARV (NVP 6-12 semaines). Arrêt progressif à 6 mois (sevrage rapide 1 mois). \textbf{Jamais d'allaitement mixte (lait + eau/bouillie) avant 6 mois} (Risque transmission $\times$ 10-20).
    \end{itemize}
\end{enumerate}

\noindent\textbf{Tâche F : Note de synthèse pour l'Infirmier(ère) du CSI}
\begin{center}
\begin{tabularx}{\linewidth}{|X|}\hline
\rowcolor{popblueL}
\textbf{NOTE DE SERVICE / TRANSFERT DOSSIER PATIENT VIH} \\ \hline
\textbf{À :} Infirmier(ère) Référent(e) VIH - CSI Nkolbisson \hfill \textbf{Date :} \today \\
\textbf{De :} Médecin Traité / Chef de Service \hfill \textbf{Urgence :} \textbf{HAUTE (Stade Immuno 4)} \\
\textbf{Patient :} M. K., 32 ans, Maçon, NIP : [Anonymat] \\ \hline
\end{tabularx}
\end{center}

\vspace{0.3cm}
\noindent\textbf{1. DIAGNOSTIC \& STADE}
\begin{itemize}
    \item \textbf{VIH-1 confirmé} (Algo 3 tests + + +).
    \item \textbf{Stade Clinique OMS : 1 (Asymptomatique).}
    \item \textbf{Stade Immunologique OMS : 4 (Sida) : CD4 = 180/$\mu$L (12\%).}
    \item \textbf{Classification CDC : A3.}
    \item \textbf{Charge Virale : 450 000 copies/mL (5,65 log).}
    \item \textbf{Contexte :} Conjointe enceinte 12 SA (Séronégative à CPN1). Anémie modérée (Hb 11 g/dL). Fonction rénale normale (DFG 95).
\end{itemize}

\noindent\textbf{2. TRAITEMENT PRESCRIT (Démarrage \textbf{IMMÉDIAT - J0 ou J1 max})}
\begin{itemize}
    \item \textbf{TLD (TDF 300 / 3TC 300 / DTG 50 mg) : 1 cp/jour, \textbf{LE MATIN} (à jeun ou pas, mais matin impératif pour insomnie DTG).}
    \item \textbf{Cotrimoxazole 960 mg (CTX) : 1 cp/jour} (Prophylaxie IO - CD4 $<$ 350). À prendre au repas (tolérance digestive).
    \item \textbf{Acide Folique 5 mg : 1 cp/jour} (Grossesse conjointe + CTX antifolate).
    \item \textbf{Supplément Fer/Folate :} Pour anémie (Hb 11) $\rightarrow$ Prescrire Fer polymaltose 100 mg 1 cp/j + Acide folique déjà là.
\end{itemize}

\noindent\textbf{3. ÉDUCATION THÉRAPEUTIQUE (ETP) - POINTS CLÉS À RENFORCER À CHAQUE VISITE (J7, J30, M3, M6)}
\begin{enumerate}
    \item \textbf{OBSERVANCE = VIE.} Objectif \textbf{$>$ 95\%} (pas plus de 3 doses oubliées/mois). Oubli = Résistance. « Si oubli $<$ 12h : prendre tout de suite. Si $>$ 12h : sauter, ne pas doubler. »
    \item \textbf{DTG : Prise MATIN.} Risque \textbf{insomnie, cauchemars, anxiété} (fréquent 1-2 premiers mois). Si persistant $>$ 4 sem $\rightarrow$ avis médecin (switch possible mais rare).
    \item \textbf{DTG : Surveillance POIDS \& GLYCÉMIE.} Prise de poids fréquente (moyenne +3 à +5 kg an 1, surtout femme, DTG + TDF/3TC). Conseils hygiéno-diététiques (marche, moins sucre/huile). Glycémie à jeun à M6, M12.
    \item \textbf{TDF : Surveillance RÉNALE \& OSSEUSE.} Créatinine/DFG à M3, M6, puis annuel. Douleur osseuse/ostéoporose rare mais possible long terme.
    \item \textbf{3TC / TDF : Troubles digestifs (nausées, diarrhée) J1-J14.} Prendre au repas si nausées. Ne pas arrêter.
    \item \textbf{COTRIMOXAZOLE :} Prendre au repas. \textbf{Signes alerte (Urgences) :} Éruption cutanée (bulle, muqueuses), Fièvre, Jaunisse, Saignements inhabituels $\rightarrow$ Arrêter CTX + Venir IMMÉDIATEMENT.
    \item \textbf{JAMAIS D'ARRÊT SANS AVIS MÉDECIN.} Risque rebond viral massif, résistance, SRI sévère.
\end{enumerate}

\noindent\textbf{4. PLAN DE SUIVI BIOLOGIQUE (Prescriptions à remettre au patient)}
\begin{center}
\begin{tabularx}{\linewidth}{|l|X|}\hline
\textbf{Moment} & \textbf{Examens demandés} \\ \hline
\textbf{J0 (Aujourd'hui)} & CV (déjà fait), CD4 (fait), Numération Formule Sanguine (NFS), Créatinine/DFG, Glycémie, Hb, Groupe sanguin, Sérologies (VHB, VHC, Syphilis - TPHA/VDRL), Toxoplasmose (IgG/IgM), Grossesse conjointe (Test rapide). \\ \hline
\textbf{J30 (M1)} & NFS (Tolérance AZT/3TC/CTX), Créatinine (TDF), \textbf{Observance (Comptage pilules / Entretien)}. \\ \hline
\textbf{M3 (Mois 3 - CRUCIAL)} & \textbf{CHARGE VIRALE (CV) - Objectif $<$ 50 copies/mL}, NFS, Créatinine, Glycémie. \\ \hline
\textbf{M6 (Mois 6)} & CV (Confirmation indétectabilité), \textbf{CD4 (Remontée attendue $>$ 100-150 gain)}, NFS, Créatinine, Glycémie, Lipidogramme. \\ \hline
\textbf{Puis Annuel} & CV (Si 2 CV indétectables conséccutives), CD4 (Si CV indétectable stable, CD4 optionnel), Créatinine, Glycémie. \\ \hline
\end{tabularx}
\end{center}

\noindent\textbf{5. PRISE EN CHARGE CONJOINTE (PTME / PrEP) - \textbf{ACTION INFIRMIER(ÈRE) PRIORITAIRE CETTE SEMAINE}}
\begin{itemize}
    \item \textbf{Dépistage VIH Conjointe :} Vérifier résultat CPN1 (Négatif). \textbf{Refaire test rapide aujourd'hui} (Fenêtre sérologique).
    \item \textbf{Initiation PrEP (TDF/3TC 1 cp/j) :} \textbf{IMMÉDIATE} si test négatif confirmé. Expliquer : Prise quotidienne, effets indésirables (nausées J1-J7), Suivi Créatinine trimestriel.
    \item \textbf{Dépistages répétés Conjointe :} Mensuel pendant grossesse + Accouchement + 6 sem post-partum + 6 sem post-sevrage.
    \item \textbf{Conseil Allaitement :} Encourager \textbf{Allaitement Maternel Exclusif 6 mois}. Expliquer sécurité (Mère Négative + Père sous ARV + PrEP Mère).
    \item \textbf{Prophylaxie Enfant :} Prévoir sirop NVP (10-15 mg/j) pour 6 semaines à la naissance (Protocole standard exposition VIH Cameroun). Préscrire à l'avance (CPN 8-9 mois).
\end{itemize}

\noindent\textbf{6. PROCHAIN RDV MÉDECIN : \textbf{J7 (Semaine 1)} puis \textbf{M1, M3}.}
\noindent\textbf{Prochain RDV INFIRMIER(ÈRE) (ETP/Observance/PrEP Conjointe) : \textbf{J3, J7, J14, J30, puis Mensuel.}}

\vspace{0.2cm}
\noindent\textbf{Signature Médecin :} \hfill \textbf{Signature Infirmier(ère) (Accusé réception) :}
\end{exoresolu}

\courssub{Bilan de l'activité}
Cette activité d'intégration a permis de mettre en évidence la \textbf{continuité entre la physiopathologie virale (cycle, dynamique CV/CD4), la sémiologie biologique (stadification), la pharmacologie ciblée (ARV, prophylaxie) et la santé publique (PTME, PrEP, TasP)}.

\begin{itemize}
    \item \textbf{Savoirs mobilisés :} Structure VIH, cycle réplication (cibles ARV), dynamique infection (phase aiguë, set point, défaillance), classification OMS/CDC, mécanismes ARV (INTI, II), indications CTX, algorithmes dépistage, stratégies prévention combinée (PrEP, TasP, PTME), particularités DTG/grossesse.
    \item \textbf{Savoir-faire développés :} Interprétation de documents cliniques complexes (sérologie, biologie, courbes modélisées), raisonnement clinique (stadification, urgence), rédaction professionnelle (note de transfert), conseil patient (éducation thérapeutique, prévention), calculs simples (pente CD4, log CV).
    \item \textbf{Savoir-être (Compétences transverses) :} Rigueur scientifique (justifier chaque décision par un seuil, une étude, une guideline), empathie et éthique (confidentialité, non-discrimination, soutien couple), responsabilité (urgence Stade 4, sécurité fœtus), esprit d'équipe (transmission infirmier/médecin/sage-femme).
\end{itemize}

\begin{attention}[Pièges fréquents à éviter pour le BEPC]
\begin{itemize}
    \item \textbf{Confondre Stade Clinique et Stade Immunologique :} Un patient asymptomatique (Clinique 1) peut être Immunologique 4 (Sida). C'est le CD4 qui guide l'urgence CTX/ARV.
    \item \textbf{Oublier la fenêtre sérologique :} Un test anticorps négatif ne prouve pas l'absence d'infection si risque $<$ 3 mois. Penser CV / Ag p24.
    \item \textbf{Dire que DTG est interdit au 1er trimestre :} \textbf{FAUX} depuis 2019 (OMS) et 2020 (PNLS Cameroun). C'est le \textbf{choix de 1ère ligne}.
    \item \textbf{Oublier le Cotrimoxazole à CD4 $<$ 350 :} C'est une prescription \textbf{vitale} (Pneumocystose = urgence vitale).
    \item \textbf{Confondre PrEP (séronégatif) et PEP (urgence post-exposition 72h) et ARV (séropostif).}
    \item \textbf{Dire « Guéri » au lieu de « Contrôlé / Indétectable » :} Le VIH persiste (réservoirs). L'arrêt ARV = Rebond viral certain.
\end{itemize}

\end{attention}

\begin{aretenir}[Message clé pour l'examen BEPC]
Face à un dossier VIH : \textbf{1. Confirmer sérologie (Algo 3 tests). 2. Stadifier (CD4 + Clinique). 3. Prescrire ARV 1ère ligne (TLD) IMMÉDIATEMENT + CTX si CD4 $<$ 350. 4. Éduquer (Observance, Effets indésirables, TasP). 5. Protéger le couple/enfant (PrEP, PTME, Allaitement). 6. Suivre (CV M3, M6).}
\end{aretenir}

\newpage

\coursec{Méthodes et exercices résolus}

\courssub{Méthodes et exercices résolus}

\begin{methode}[Analyser le cycle de réplication du VIH à partir d'un schéma ou d'un texte]
\textbf{Objectif :} Restituer dans l'ordre chronologique les étapes de la multiplication du VIH dans la cellule cible (lymphocyte T CD4+) et identifier le rôle des enzymes virales.
\begin{enumerate}
    \item \textbf{Identifier la cellule cible :} Repérer la présence du récepteur \cle{CD4} et des co-récepteurs (CCR5 ou CXCR4) à la surface du lymphocyte T4.
    \item \textbf{Ordonner les 7 étapes clés :}
    \begin{enumerate}
        \item \textbf{Fixation (Adsorption) :} gp120 (enveloppe) se lie au CD4 puis au co-récepteur.
        \item \textbf{Fusion / Pénétration :} gp41 médie la fusion de l'enveloppe virale avec la membrane cellulaire $\rightarrow$ libération du capside dans le cytoplasme.
        \item \textbf{Transcription inverse :} L'enzyme \cle{transcriptase inverse (TI)} synthétise un ADN simple brin (-) sur le matrice d'ARN viral (+), puis dégrade l'ARN (activité RNase H) et synthétise le brin complémentaire $\rightarrow$ \cle{ADN double brin proviral}.
        \item \textbf{Intégration :} L'\cle{intégrase} insère l'ADN proviral dans le chromosome de la cellule hôte (ADN hôte). Le virus devient \cle{provirus}.
        \item \textbf{Transcription / Expression génique :} L'ADN proviral est transcrit par l'ARN polymérase II de l'hôte en ARNm viraux (génome complet et ARNm épissés).
        \item \textbf{Traduction / Assemblage :} Les ARNm sont traduits au niveau des ribosomes (protéines structurales Gag, Pol, Env). Les protéines et l'ARN génomique s'assemblent au niveau de la membrane plasmique $\rightarrow$ bourgeonnement (particule immature).
        \item \textbf{Maturation :} La \cle{protéase (PR)} clive les polyprotéines Gag et Pol $\rightarrow$ virion mature \textbf{infectieux}.
    \end{enumerate}
    \item \textbf{Repérer les cibles thérapeutiques (ARV) :} Associer chaque classe d'ARV à l'enzyme ou l'étape bloquée (ITI, INI, IP, IF, IAE).
\end{enumerate}
\textbf{Reflexe BEPC :} Utiliser les termes précis : \emph{transcription inverse} (pas « transcription » seul), \emph{ADN proviral}, \emph{intégrase}, \emph{protéase}. Ne pas confondre ARN viral (génome) et ARNm (messager).
\end{methode}

\begin{exoresolu}[Cycle de réplication et cibles des ARV]
\textbf{Énoncé :} Le document ci-dessous représente schématiquement le cycle de réplication du VIH-1 dans un lymphocyte T CD4+.
\begin{center}
\begin{tikzpicture}[scale=0.7, every node/.style={font=\small}]
% Cell membrane
\draw[popblue, very thick] (-5,-3) -- (5,-3) node[midway, below=2mm] {Membrane plasmique (CD4, CCR5)};
\draw[popblue, very thick] (-5,3) -- (5,3) node[midway, above=2mm] {Noyau (ADN hôte)};
% Virus entry
\draw[poporange, thick, ->] (0,-4) -- (0,-3.2) node[midway, right] {1. Fixation/Fusion};
\node[draw, circle, fill=poporange!20, minimum size=8mm] at (0,-4.5) {Virion (ARN, TI, Int, PR)};
% Reverse transcription
\draw[popgreen, thick, ->] (0,-3) -- (0,-1.5) node[midway, right] {2. Transcription inverse};
\node[draw, rectangle, rounded corners, fill=popgreen!10, minimum width=2cm, minimum height=1cm] at (0,-1) {ADN double brin (Provirus)};
% Integration
\draw[poppurple, thick, ->] (0,-0.5) -- (0,0.5) node[midway, right] {3. Intégration};
\node[draw, ellipse, fill=poppurple!10, minimum width=3cm, minimum height=1.5cm] at (0,1.5) {ADN hôte + ADN proviral intégré};
% Transcription/Translation
\draw[poppink, thick, ->] (2,1.5) -- (3.5,1.5) node[midway, above] {4. Transcription};
\node[draw, rectangle, fill=poppink!10, minimum width=2cm, minimum height=0.8cm] at (5,1.5) {ARNm viraux};
\draw[poppink, thick, ->] (5,1.5) -- (5,0) node[midway, right] {5. Traduction};
\node[draw, rectangle, fill=poppink!10, minimum width=2cm, minimum height=0.8cm] at (5,-0.5) {Polyprotéines (Gag, Pol, Env)};
% Assembly/Budding
\draw[popgold, thick, ->] (5,-0.5) -- (2,-0.5) node[midway, above] {6. Assemblage/Bourgeonnement};
\node[draw, circle, fill=popgold!20, minimum size=8mm] at (0,-0.5) {Particule immature};
% Maturation
\draw[popdark, thick, ->] (0,-0.5) -- (0,-1.5) node[midway, left] {7. Maturation (Protéase)};
\node[draw, circle, fill=popdark!20, minimum size=8mm] at (0,-2.2) {Virion mature \textbf{infectieux}};
% ARV targets
\node[draw, dashed, rounded corners, fill=popblue!5, fit=(current bounding box.south west) (current bounding box.north east), label=above right:\textbf{Cibles ARV}] {};
\node[popblue, anchor=north west] at (-4.5, 2.5) {\tiny IF : Inhibiteurs d'Entrée (Maraviroc, Enfuvirtide)};
\node[popblue, anchor=north west] at (-4.5, 2.0) {\tiny ITI : Inhibiteurs Transcriptase Inverse (TDF, 3TC, EFV)};
\node[popblue, anchor=north west] at (-4.5, 1.5) {\tiny INI : Inhibiteurs Intégrase (DTG, RAL)};
\node[popblue, anchor=north west] at (-4.5, 1.0) {\tiny IP : Inhibiteurs Protéase (LPV/r, DRV/r, ATV/r)};
\end{tikzpicture}
\end{center}
\textbf{Questions :}
\begin{enumerate}
    \item Légender les numéros 1 à 7 sur le schéma en utilisant le vocabulaire précis du cours.
    \item Quel est le rôle de l'enzyme \textbf{transcriptase inverse} ? Pourquoi est-elle une cible majeure des ARV ?
    \item Un patient est traité par une trithérapie associant \textbf{TDF + 3TC + DTG}. Préciser la classe thérapeutique de chaque molécule et l'étape du cycle viral qu'elle bloque.
    \item Expliquer pourquoi le VIH est qualifié de \emph{retrovirus} et les conséquences de l'intégration pour la guérison.
\end{enumerate}
\tcblower
\textbf{Solution détaillée :}
\begin{enumerate}
    \item \textbf{Légende du cycle :}
    1. \textbf{Fixation sur CD4/co-récepteur et Fusion} (Pénétration du capside).
    2. \textbf{Transcription inverse} : Synthèse de l'ADN proviral double brin à partir de l'ARN viral par la \textbf{transcriptase inverse (TI)}.
    3. \textbf{Intégration} : Insertion de l'ADN proviral dans le génome de la cellule hôte par l'\textbf{intégrase}. Formation du \textbf{provirus}.
    4. \textbf{Transcription} : Synthèse des ARNm viraux par l'ARN polymérase II de l'hôte.
    5. \textbf{Traduction} : Synthèse des polyprotéines virales (Gag, Pol, Env) aux ribosomes.
    6. \textbf{Assemblage et Bourgeonnement} : Formation de particules immatures à la membrane plasmique.
    7. \textbf{Maturation} : Clivage des polyprotéines par la \textbf{protéase (PR)} $\rightarrow$ \textbf{Virion mature infectieux}.
    \item \textbf{Rôle de la Transcriptase Inverse (TI) :} Elle catalyse la \textbf{transcription inverse} (flux d'information ARN $\rightarrow$ ADN), étape unique aux rétrovirus et absente chez l'hôte (pas de TI cellulaire). Elle synthétise l'ADN double brin proviral à partir du génome ARN simple brin. \textbf{Cible majeure} car son inhibition (par les ITI : TDF, 3TC, EFV, NVP) bloque la formation de l'ADN proviral, empêchant toute intégration et réplication ultérieure. Sa fidélité faible (pas d'activité de relecture) génère des mutations $\rightarrow$ résistance.
    \item \textbf{Analyse de la trithérapie TDF + 3TC + DTG (schéma standard 1ère ligne OMS/Cameroun) :}
    \begin{center}
    \begin{tabularx}{\linewidth}{|l|X|X|}\hline
    \rowcolor{popblueL} \textbf{Molécule} & \textbf{Classe thérapeutique} & \textbf{Étape bloquée / Cible enzymatique} \\ \hline
    \textbf{TDF} (Ténofovir) & ITI-N (Inhibiteur Nucleos(t)idique de la TI) & \textbf{Transcription inverse} : Se lie au site actif de la TI (compétiteur du dATP) $\rightarrow$ terminaison de chaîne. \\ \hline
    \textbf{3TC} (Lamivudine) & ITI-N (Inhibiteur Nucleosidique de la TI) & \textbf{Transcription inverse} : Analogue de la cytidine $\rightarrow$ terminaison de chaîne. Synergie avec TDF. \\ \hline
    \textbf{DTG} (Dolutégravir) & \textbf{INI} (Inhibiteur de l'\textbf{Intégrase}) & \textbf{Intégration} : Bloque le transfert de brin (strand transfer) $\rightarrow$ empêche l'insertion de l'ADN viral dans l'ADN hôte. \\ \hline
    \end{tabularx}
    \end{center}
    \item \textbf{Qualification de rétrovirus :} Le VIH possède un génome \textbf{ARN} et utilise une \textbf{transcriptase inverse} pour synthétiser de l'ADN (flux inverse : ARN $\rightarrow$ ADN), contredisant le dogme central classique (ADN $\rightarrow$ ARN $\rightarrow$ Protéine).
    \textbf{Conséquence de l'intégration :} L'ADN proviral devient partie intégrante du génome de la cellule hôte. Il est \textbf{irréversible} (sauf édition génomique future). La cellule infectée devient une \textbf{réservoir viral} permanent. C'est pourquoi les ARV suppriment la réplication (\textbf{charge virale indétectable}) mais \textbf{ne guérissent pas} l'infection (arrêt du traitement = rebond viral depuis les réservoirs).
\end{enumerate}
\textbf{Conclusion :} La connaissance précise de l'ordre des étapes enzymatiques (TI $\rightarrow$ Intégrase $\rightarrow$ Protéase) permet de comprendre le mécanisme d'action des différentes classes d'ARV et la nécessité d'une association (trithérapie) pour bloquer le cycle à plusieurs points et prévenir la résistance.
\end{exoresolu}

\begin{methode}[Interpréter l'évolution des marqueurs biologiques (CD4, Charge Virale) au cours des phases cliniques]
\textbf{Objectif :} Lier la physiopathologie (destruction des LT4) aux marqueurs biologiques pour identifier la phase de l'infection (Primaire, Asymptomatique, Symptomatique/SIDA) et justifier l'urgence thérapeutique.
\begin{enumerate}
    \item \textbf{Repérer les deux courbes clés sur un graphique (temps en abscisse) :}
    \begin{itemize}
        \item \textbf{Charge Virale (CV) / Viremie (copies/mL) :} Échelle \textbf{logarithmique} (souvent). Pic énorme en phase primaire ($>10^6$), chute brutale (réponse immune), plateau (point de set) en phase asymptomatique, remontée exponentielle en phase SIDA.
        \item \textbf{Taux de Lymphocytes T CD4+ (cellules/$\mu$L ou /mm$^3$) :} Échelle linéaire. Chute transitoire en phase primaire, déclin \textbf{lent et continu} ($\approx$ 50-80 cellules/$\mu$L/an) en phase asymptomatique, chute brutale $< 200$ en phase SIDA.
    \end{itemize}
    \item \textbf{Identifier les 3 phases cliniques (Classification OMS/Cameroun) :}
    \begin{center}
    \begin{tabularx}{\linewidth}{|l|X|X|X|}\hline
    \rowcolor{popblueL} \textbf{Phase} & \textbf{Clinique} & \textbf{CD4 (cellules/$\mu$L)} & \textbf{Charge Virale} \\ \hline
    \textbf{Primaire (Aiguë)} & Syndrome pseudo-grippal / mononucléosique (fièvre, adénopathies, éruption) $\sim$ 2-4 sem. post-infection. Souvent inaperçue. & Chute transitoire puis remontée partielle. & \textbf{Très élevée} ($> 10^6$ cp/mL). Pic de contagiosité. \\ \hline
    \textbf{Asymptomatique (Chronique)} & Aucune symptôme (ou adénopathies généralisées persistantes - AGP). Durée : \textbf{5 à 10 ans} sans traitement. & \textbf{Déclin progressif} ($> 500$ au début $\rightarrow$ $< 350$). & \textbf{Plateau (Set point)} variable ($10^3$ à $10^5$ cp/mL). \\ \hline
    \textbf{Symptomatique / SIDA} & Infections opportunistes (TB, candidose, pneumocystose, toxoplasmose...), cancers (Kaposi, lymphome), amaigrissement sévère. & \textbf{$< 200$} (Seuil définition SIDA). Immunité cellulaire effondrée. & \textbf{Remontée massive} ($> 10^5 - 10^6$ cp/mL). \\ \hline
    \end{tabularx}
    \end{center}
    \item \textbf{Calculer le déclin annuel des CD4 :} $\frac{\text{CD4}_{\text{initial}} - \text{CD4}_{\text{final}}}{\text{Nombre d'années}}$.
    \item \textbf{Seuils décisionnels (Guidelines Cameroun/OMS 2024) :} \textbf{Traitement pour TOUS (Treat All)} d\`es le diagnostic, quelque soit le taux de CD4. Mais CD4 $< 200$ = Urgence (prophylaxie cotrimoxazole + dépistage TB + ARV rapide).
\end{enumerate}
\textbf{Piège :} Ne pas confondre \emph{charge virale} (quantité de virus) et \emph{taux de CD4} (état immunitaire). Une CV indétectable ($< 50$ cp/mL) sous ARV ne signifie pas guérison (réservoirs), mais \textbf{non-transmissibilité sexuelle (I=I : Indétectable = Intransmissible)}.
\end{methode}

\begin{exoresolu}[Analyse de l'évolution immunovirologique d'un patient non traité]
\textbf{Énoncé :} Le graphique ci-dessous montre l'évolution de la charge virale (CV, en copies/mL, échelle log) et du taux de lymphocytes T CD4+ (en cellules/$\mu$L) chez un patient infecté par le VIH-1, non traité, suivi sur 10 ans.
\begin{center}
\begin{tikzpicture}[scale=0.8]
\begin{axis}[
    width=\linewidth, height=8cm,
    xlabel={Temps (années)}, xmin=0, xmax=10,
    ylabel={CD4 (cellules/$\mu$L)}, ymin=0, ymax=1200,
    axis y line*=left, axis x line=bottom,
    xtick={0,1,2,3,4,5,6,7,8,9,10},
    ytick={0,200,350,500,800,1000},
    grid=major, grid style={popline},
    legend style={at={(0.5,1.02)}, anchor=south, legend columns=2, font=\small}
]
% CD4 curve
\addplot[popblue, very thick, smooth, domain=0:10, samples=100] {1000 - 60*x - 10*sin(deg(x*2))};
\addlegendentry{CD4 (cellules/$\mu$L)}
% Viral Load (secondary axis) - simulated via coordinates for log scale feel
\end{axis}
\begin{axis}[
    width=\linewidth, height=8cm,
    xlabel={}, xmin=0, xmax=10,
    ylabel={Charge Virale (copies/mL)}, ymin=1, ymax=1000000,
    axis y line*=right, axis x line=none,
    ymode=log,
    ytick={1,100,1000,10000,100000,1000000},
    yticklabels={$1$,$10^2$,$10^3$,$10^4$,$10^5$,$10^6$},
    legend style={at={(0.5,0.95)}, anchor=north, legend columns=2, font=\small}
]
% Viral Load curve: Peak at 0.25, drop, plateau, rise at end
\addplot[poporange, very thick, smooth, domain=0:10, samples=200] 
    ({x}, {10^(6.5 - 5.5*exp(-(x-0.25)^2/0.05) - 0.3*x + 1.5*max(0,x-7)^2/9)});
\addlegendentry{Charge Virale (copies/mL)}
% Phase lines
\draw[popmuted, dashed] (axis cs:0.5,1) -- (axis cs:0.5,1000000);
\draw[popmuted, dashed] (axis cs:7,1) -- (axis cs:7,1000000);
\node[popmuted, rotate=90, anchor=south] at (axis cs:0.25,500000) {Phase Primaire};
\node[popmuted, rotate=90, anchor=south] at (axis cs:3.5,500000) {Phase Asymptomatique};
\node[popmuted, rotate=90, anchor=south] at (axis cs:8.5,500000) {Phase SIDA};
\end{axis}
\end{tikzpicture}
\end{center}
\textbf{Questions :}
\begin{enumerate}
    \item En vous basant sur les variations des deux courbes, délimiter sur le graphique (années approximatives) les trois phases cliniques de l'infection VIH. Justifier par les valeurs caractéristiques.
    \item Calculer le déclin moyen annuel du taux de CD4 pendant la phase asymptomatique (entre l'an 1 et l'an 7). Exprimer en cellules/$\mu$L/an.
    \item Le patient présente une pneumonie à \emph{Pneumocystis jirovecii} à l'an 9. Quel est le stade clinique ? Quel est le seuil de CD4 associé ? Quelle prophylaxie primaire est indiquée ?
    \item Expliquer le mécanisme immunologique responsable de la chute brutale des CD4 en phase SIDA (rôle de l'activation immune chronique, apoptose, infection directe).
\end{enumerate}
\tcblower
\textbf{Solution détaillée :}
\begin{enumerate}
    \item \textbf{Délimitation des phases :}
    \begin{itemize}
        \item \textbf{Phase Primaire (An 0 à $\sim$ 0,5 an) :} Pic de \textbf{Charge Virale} $\approx 10^6 - 10^7$ copies/mL (très contagieux). Chute transitoire des \textbf{CD4} (ex: de 1000 à $\sim$ 600) suivie d'une remontée partielle (réponse immune CTL).
        \item \textbf{Phase Asymptomatique (An $\sim$ 0,5 à $\sim$ 7 ans) :} \textbf{Charge Virale} se stabilise au \textbf{point de set (set point)} $\approx 10^4 - 10^5$ copies/mL. \textbf{CD4} déclinent \textbf{régulièrement et lentement} (ex: de $\sim$ 800 à $\sim$ 350 cellules/$\mu$L). Patient asymptomatique (ou AGP).
        \item \textbf{Phase SIDA (An $\sim$ 7 à 10 ans) :} \textbf{CD4} chutent \textbf{brutalement} sous le seuil de \textbf{200 cellules/$\mu$L} (ici $\sim$ 50 à l'an 10). \textbf{Charge Virale} repart à la hausse exponentielle ($> 10^5$). Apparition d'infections opportunistes.
    \end{itemize}
    \item \textbf{Calcul du déclin annuel (Phase Asymptomatique) :}
    Lecture graphique approximative : CD4(an 1) $\approx$ 850 cellules/$\mu$L ; CD4(an 7) $\approx$ 350 cellules/$\mu$L.
    \[
    \text{Déclin} = \frac{850 - 350}{7 - 1} = \frac{500}{6} \approx \mathbf{83,3 \text{ cellules/$\mu$L/an}}
    \]
    (Valeur cohérente avec la littérature : 50-100 cellules/$\mu$L/an).
    \item \textbf{Stade clinique :} \textbf{Stade 3 ou 4 (SIDA)} selon classification OMS (infection opportuniste définissant le SIDA).
    \textbf{Seuil CD4 :} \textbf{$< 200$ cellules/$\mu$L} (ici $\ll$ 200).
    \textbf{Prophylaxie primaire indiquée :} \textbf{Cotrimoxazole (Bactrim®) 960 mg/jour} (prévention pneumocystose, toxoplasmose, infections bactériennes, diarrhées). \textbf{+ Dépistage systématique de la Tuberculose (recherche BK, GeneXpert, Rx poumon)}. \textbf{Initiation ARV en urgence (dans les 2 semaines, voire jours si méningite cryptococcique)}.
    \item \textbf{Mécanismes de la chute brutale en phase SIDA :}
    \begin{enumerate}
        \item \textbf{Infection directe et lyse :} Réplication massive du virus (CV élevée) $\rightarrow$ mort cellulaire directe des LT4 infectés.
        \item \textbf{Apoptose (mort programmée) des cellules non infectées (bystander effect) :} Protéines virales (gp120, Tat, Nef) + \textbf{activation immune chronique} (inflammation persistante, cytokines TNF-$\alpha$, Fas/FasL) $\rightarrow$ mort des LT4 « spectateurs ».
        \item \textbf{Défaillance de la régénération :} Destruction de l'architecture ganglionnaire (fibrose) $\rightarrow$ perte du microenvironnement nécessaire à l'homéostasie des lymphocytes (IL-7). Thymus involué chez l'adulte.
        \item \textbf{Épuisement immunitaire :} Expression de marqueurs d'épuisement (PD-1) sur les LT8 et LT4 restants $\rightarrow$ dysfonctionnement.
    \end{enumerate}
\end{enumerate}
\textbf{Conclusion :} L'analyse conjointe CD4/CV permet le stade immunologique et virologique. Le franchissement du seuil de 200 CD4 marque le basculement en phase SIDA, justifiant une prise en charge d'urgence (ARV + prophylaxies).
\end{exoresolu}

\begin{methode}[Résoudre un problème de transmission / prévention : Approche « Combination Prevention »]
\textbf{Objectif :} À partir d'une situation concrète (couple sérodifférent, femme enceinte, exposition accidentelle, population clé), proposer la stratégie de prévention adaptée (biomédicale, comportementale, structurelle).
\begin{enumerate}
    \item \textbf{Identifier le mode de transmission concerné :} Sexuel (vaginal, anal), Sanguin (PIA, scarifications, PTME), Mère-Enfant (grossesse, accouchement, allaitement).
    \item \textbf{Appliquer la prévention combinée (Niveaux) :}
    \begin{itemize}
        \item \textbf{Biomédicale :} \textbf{PrEP} (Prophylaxie Pré-Exposition : TDF/3TC ou TDF/FTC) pour séronégatifs à risque élevé. \textbf{PEP} (Prophylaxie Post-Exposition : TDF/3TC/DTG 28j, $\le$ 72h idéalement $\le$ 24h). \textbf{TasP} (Treatment as Prevention : ARV pour le partenaire séropositif $\rightarrow$ CV indétectable = I=I). \textbf{PTME} (Prévention Transmission Mère-Enfant : ARV mère + AZT/NVP nourrisson + accouchement sécurisé + conseil alimentation).
        \item \textbf{Comportementale :} \textbf{Préservatif} (masculin/féminin) usage correct/systématique. Réduction partenaires. Dépistage régulier. Circoncision masculine médicale (CMM) réduire risque H$\rightarrow$F 60\%.
        \item \textbf{Structurelle :} Lutte contre stigmatisation, éducation sexuelle complète, réduction des risques usagers de drogues (RDR), autonomisation femmes/filles.
    \end{itemize}
    \item \textbf{Calculer l'efficacité / Nombre Nécessaire à Traiter (NNT) si données fournies.}
    \item \textbf{Rédiger la conclusion :} « La stratégie recommandée pour [situation] est [X] car [justification scientifique + contexte camerounais]. »
\end{enumerate}
\textbf{Reflexe BEPC :} Citer \textbf{PrEP}, \textbf{PEP}, \textbf{TasP}, \textbf{PTME}, \textbf{CMM}. Préciser les délais (PEP $< 72$h, PrEP continue). Ne pas oublier le \textbf{dépistage} (porte d'entrée).
\end{methode}

\begin{exoresolu}[Prise en charge d'un couple sérodifférent et PTME au Cameroun]
\textbf{Énoncé :} Marie (28 ans, séronégative) et Paul (32 ans, VIH+ sous ARV depuis 18 mois, Charge Virale indétectable $< 20$ cp/mL depuis 12 mois) souhaitent avoir un enfant. Ils consultent au Centre de Traitement Ambulatoire (CTA) de l'Hôpital de District de Bafoussam.
\begin{enumerate}
    \item Expliquer le concept \textbf{TasP (Treatment as Prevention)} et \textbf{I=I (Indétectable = Intransmissible)}. Paul peut-il transmettre le VIH à Marie par voie sexuelle sans préservatif dans ces conditions ?
    \item Quelle stratégie de prévention biomédicale proposer à Marie \textbf{avant et pendant la période de conception} ? Justifier.
    \item Marie tombe enceinte. Décrire la prise en charge \textbf{PTME (Prévention de la Transmission Mère-Enfant)} selon le protocole Cameroun/OMS (Option B+) : traitement mère, prophylaxie nourrisson, suivi virologique, alimentation.
    \item À la naissance, l'enfant reçoit du \textbf{Névirapine (NVP) sirop} pendant 6 semaines. Expliquer le mécanisme d'action de la NVP et pourquoi on l'utilise en monothérapie prophylactique ici (alors qu'elle est contre-indiquée en monothérapie thérapeutique).
\end{enumerate}
\tcblower
\textbf{Solution détaillée :}
\begin{enumerate}
    \item \textbf{TasP / I=I :} Le principe : Une personne VIH+ sous ARV efficace avec \textbf{Charge Virale (CV) indétectable ($< 50$ ou $< 200$ cp/mL) de façon durable ($\ge$ 6 mois)} ne transmet \textbf{PAS} le VIH par voie sexuelle (\textbf{Preuves : PARTNER, Opposites Attract, HPTN 052}).
    \textbf{Réponse :} \textbf{NON}, Paul ne peut pas transmettre le VIH à Marie par voie sexuelle sans préservatif, car sa CV est indétectable depuis 12 mois ($> 6$ mois) et il est observant. C'est la base scientifique du \textbf{conception naturelle sécurisée} pour couple sérodifférent.
    \item \textbf{Stratégie pour Marie (Pré-conception / Conception) :} \textbf{PrEP (Prophylaxie Pré-Exposition) : TDF/3TC (ou TDF/FTC) 1 cp/jour}.
    \textbf{Justification :} « Ceinture et bretelles ». Bien que TasP soit très efficace, la PrEP offre une protection \textbf{personnelle} à Marie en cas de : « blip » viral transitoire (infection intercurrentes), non-observance ponctuelle de Paul, ou rapports hors couple. Elle est recommandée par l'OMS et le PNLS Cameroun pour les partenaires séronégatifs de personnes VIH+ désirant un enfant. Arrêt PrEP après confirmation grossesse (si CV partenaire reste indétectable) ou poursuite si risque persistant.
    \item \textbf{PTME Option B+ (Standard Cameroun) :}
    \begin{itemize}
        \item \textbf{Mère :} \textbf{ARV à vie (TDF/3TC/DTG 1cp/j)} d\`es le diagnostic (ou poursuite si déjà sous ARV). Objectif : \textbf{CV indétectable ($< 50$ cp/mL) à l'accouchement}. Suivi CV à 34-36 SA et à l'accouchement.
        \item \textbf{Accouchement :} Voie basse si CV indétectable. Césarienne programmée (38 SA) si CV $> 1000$ cp/mL ou inconnue. Perfusion AZT (Zidovudine) mère si CV $> 50$ cp/mL.
        \item \textbf{Nourrisson :} \textbf{Prophylaxie NVP sirop (dose poids) pendant 6 semaines} (si mère CV indétectable $\ge$ 4 sem) OU \textbf{AZT + NVP 6 sem} (si mère CV détectable, accouchement d'urgence, ou mère non traitée). \textbf{Dépistage virologique (PCR ADN/ARN) :} J0 (naissance), 6 sem (sevrage prophylaxie), 9 mois, 18 mois (test sérologique définitif).
        \item \textbf{Alimentation :} \textbf{Allaitement maternel exclusif (AME) 6 mois} recommandé au Cameroun (protection mortalité néonatale $>$ risque résiduel transmission si mère ARV observante CV indétectable). Sevrage progressif 6-12 mois. \textbf{Pas de lait artificiel} sauf conditions AFASS (Acceptable, Faisable, Abordable, Durable, Sûr) réunies.
    \end{itemize}
    \item \textbf{Névirapine (NVP) :} \textbf{ITI-NN (Inhibiteur Non Nucléosidique de la Transcriptase Inverse)}. Se lie au site allostérique de la TI $\rightarrow$ changement conformationnel $\rightarrow$ inhibition non compétitive.
    \textbf{Pourquoi monothérapie prophylactique acceptée ici ?}
    \begin{itemize}
        \item \textbf{Durée courte (6 sem) :} Risque de sélection de mutants résistants (K103N) très faible sur si courte durée vs traitement à vie.
        \item \textbf{Charge virale maternelle indétectable :} Exposition virale du nourrisson quasi nulle. La NVP agit comme « filet de sécurité » contre une exposition résiduelle lors de l'accouchement/allaitement.
        \item \textbf{Coût/Faisabilité :} Sirop pédiatrique, thermostable, peu cher, dose unique hebdo possible (mais 6 sem quotidienne standard).
    \end{itemize}
    \textbf{Attention :} En \textbf{traitement curatif}, la monothérapie NVP est \textbf{formellement contre-indiquée} (résistance rapide en 2-4 semaines $\rightarrow$ échec virologique, perte options futures).
\end{enumerate}
\textbf{Conclusion :} La combinaison TasP (Paul) + PrEP (Marie pré-conception) + PTME Option B+ (Grossesse/Allaitement) permet d'atteindre l'objectif \textbf{« Génération sans Sida »} (élimination transmission mère-enfant $< 5\%$ à 18 mois au Cameroun).
\end{exoresolu}

\begin{methode}[Expliquer le mécanisme d'action des ARV et interpréter une résistance]
\textbf{Objectif :} Classer les ARV par cible enzymatique, décrire leur mode d'action (compétitif/non-compétitif, allostérique) et comprendre la genèse des résistances (barrière génétique).
\begin{enumerate}
    \item \textbf{Connaître les 5 classes majeures (Ordre du cycle viral) :}
    \begin{center}
    \begin{tabularx}{\linewidth}{|l|l|X|}\hline
    \rowcolor{popblueL} \textbf{Classe (Abrév.)} & \textbf{Cible / Étape} & \textbf{Mécanisme / Exemples 1ère ligne Cameroun} \\ \hline
    \textbf{IE / IF} (Inhibiteurs Entrée/Fusion) & Fixation/Fusion (gp120/gp41, CCR5) & Maraviroc (CCR5), Enfuvirtide (gp41 - injectable, 2ème ligne). \\ \hline
    \textbf{ITI-N / NRTI} (Nucléos(t)idiques) & \textbf{Transcriptase Inverse (TI)} & \textbf{Analogues nucléos(t)idiques} : Competiteurs substrats $\rightarrow$ \textbf{Terminaison de chaîne} (manque 3'-OH). \textbf{TDF, 3TC, FTC, AZT, ABC}. \\ \hline
    \textbf{ITI-NN / NNRTI} (Non-Nucléosidiques) & \textbf{Transcriptase Inverse (TI)} & \textbf{Site allostérique} (poche hydrophobicité) $\rightarrow$ Changement conformation $\rightarrow$ Inhibition non-compétitive. \textbf{EFV, NVP, RPV, DOR}. \\ \hline
    \textbf{INI / INSTI} (Intégrase Strand Transfer Inhibitors) & \textbf{Intégrase} & Bloquent \textbf{transfert de brin (Strand Transfer)} $\rightarrow$ Empêchent insertion ADN viral. \textbf{DTG, RAL, BIC, CAB}. \textbf{Barrière génétique HAUTE (DTG, BIC)}. \\ \hline
    \textbf{IP / PI} (Inhibiteurs Protéase) & \textbf{Protéase (PR)} & Inhibiteurs compétitifs (miment substrat) $\rightarrow$ Bloquent clivage Gag/Pol $\rightarrow$ Virions immatures non infectieux. \textbf{LPV/r, DRV/r, ATV/r} (Boostés par Ritonavir/Cobicistat). \\ \hline
    \end{tabularx}
    \end{center}
    \item \textbf{Résistance :} Mutation ponctuelle dans le gène codant la cible ($pol$ : RT, Intégrase, Protéase).
    \begin{itemize}
        \item \textbf{Barrière génétique basse :} 1 mutation suffit (ITI-NN : K103N pour EFV/NVP ; ITI-N : M184V pour 3TC/FTC - mais M184V hypersensibilise à TDF/AZT).
        \item \textbf{Barrière génétique haute :} Plusieurs mutations nécessaires (INI : DTG ; IP boostés : DRV/r).
    \end{itemize}
    \item \textbf{Lecture d'un rapport de génotypage (Résistance) :} Format standard : \textbf{Gène\_Position\_Mutation} (ex: \textbf{RT\_M184V}, \textbf{IN\_R263K}, \textbf{PR\_M46I}). Interprétation : \textbf{S (Sensible)}, \textbf{R (Résistant)}, \textbf{IR (Résistance Intermédiaire)}, \textbf{PL (Potentiel Faible)}. Utiliser algorithme ANRS / Stanford HIVdb.
\end{enumerate}
\textbf{Reflexe BEPC :} Différencier ITI-N (terminaison chaîne, site actif) et ITI-NN (allostérique). Savoir que \textbf{DTG (Dolutégravir)} est l'INI de référence 1ère ligne (barrière haute, peu d'effets secondaires, sans TB). \textbf{Boosting (r/c)} : Ritonavir/Cobicistat inhibent CYP3A4 $\rightarrow$ augmentent taux IP.
\end{methode}

\begin{exoresolu}[Analyse de résistance et choix de 2ème ligne]
\textbf{Énoncé :} Un patient de 35 ans, sous 1ère ligne \textbf{TDF/3TC/EFV} depuis 3 ans, présente un \textbf{échec virologique confirmé} (CV = 45 000 copies/mL à J0 et 52 000 copies/mL à 3 mois après conseil observance). Le génotypage de résistance (Stanford HIVdb) révèle les mutations majeures suivantes :
\textbf{RT : M184V, K103N} \quad \textbf{IN : Aucune} \quad \textbf{PR : Aucune}
\begin{enumerate}
    \item Interpréter l'impact de chaque mutation identifiée sur l'activité des ARV de la 1ère ligne (TDF, 3TC, EFV).
    \item La mutation \textbf{M184V} confère une résistance de haut niveau à la 3TC/FTC. Quel est son effet paradoxal sur la sensibilité à la Ténofovir (TDF) et à la Zidovudine (AZT) ? Expliquer le mécanisme biochimique.
    \item Proposer un schéma de \textbf{2ème ligne standard (OMS/Cameroun)} pour ce patient. Justifier le choix de chaque composant en lien avec le profil de résistance.
    \item Pourquoi le passage en 2ème ligne \textbf{doit s'accompagner d'un renforcement de l'observance} (conseil, groupe de pairs, allègement pilulier) ?
\end{enumerate}
\tcblower
\textbf{Solution détaillée :}
\begin{enumerate}
    \item \textbf{Interprétation des mutations :}
    \begin{itemize}
        \item \textbf{RT K103N :} Mutation classique des \textbf{ITI-NN (EFV, NVP, RPV)}. Confère une \textbf{résistance de haut niveau (R)} à l'Éfavirenz (EFV) et à la Névirapine (NVP). \textbf{EFV est inactif.}
        \item \textbf{RT M184V :} Mutation signature des \textbf{ITI-N (3TC, FTC)}. Confère une \textbf{résistance de haut niveau (R)} à la Lamivudine (3TC) et Emtricitabine (FTC). \textbf{3TC est inactif.}
        \item \textbf{TDF (Ténofovir) :} M184V \textbf{N'EST PAS} une mutation de résistance majeure pour TDF (qui sélectionne plutôt K65R). Cependant, M184V \textbf{hypersensibilise} le virus au TDF (voir Q2). TDF conserve une activité partielle/résiduelle.
        \item \textbf{IN / PR :} Aucune mutation $\rightarrow$ Intégrase et Protéase \textbf{totalement sensibles}.
    \end{itemize}
    \item \textbf{Effet paradoxal de M184V (Hypersensibilisation) :}
    \textbf{Fait :} M184V \textbf{augmente la sensibilité} (diminue l'IC50) au \textbf{TDF} et à l'\textbf{AZT}.
    \textbf{Mécanisme biochimique :} La mutation M184V (Méthionine $\rightarrow$ Valine en position 184 de la RT) rigidifie la poche de liaison des dNTP. Elle \textbf{diminue la discrimination} de la RT entre les analogues nucléotidiques (TDF-DP, AZT-TP) et les dNTP naturels, favorisant leur incorporation. De plus, M184V \textbf{réduit la capacité d'excision (excision activity)} de la RT (mécanisme de résistance à l'AZT par excision ATP-dépendante des mutations TAMs comme M41L, T215Y). $\rightarrow$ Le virus M184V est « piégé » : il incorpore plus facilement TDF/AZT et ne peut pas les enlever.
    \item \textbf{Proposition 2ème ligne standard (OMS/Cameroun 2024) :} \textbf{TDF (ou AZT) + 3TC (ou FTC) + DTG (Dolutégravir 50 mg/j)}.
    \textbf{Justification composant par composant :}
    \begin{itemize}
        \item \textbf{DTG (INI) :} \textbf{Pilier central.} Aucune mutation IN $\rightarrow$ \textbf{Activité pleinement préservée}. Barrière génétique \textbf{très haute} (nécessite 3 mutations majeures pour résistance cliniquement significative). Tolérance excellente. \textbf{Choix n°1.}
        \item \textbf{TDF + 3TC (Backbone NRTI) :} On conserve souvent \textbf{3TC} malgré M184V car : (1) M184V réduit la \textbf{capacité de réplication (fitness cost)} du virus $\rightarrow$ bénéfice clinique/immunologique même si 3TC n'est pas pleinement actif. (2) M184V hypersensibilise au TDF. (3) Simplicité (co-formulé TDF/3TC/DTG = 1 cp/j = \textbf{TLD}).
        \item \textbf{Alternative si TDF contre-indiqué (IRC) :} \textbf{AZT/3TC/DTG} (AZT actif sur M184V, mais plus d'effets secondaires hématologiques).
    \end{itemize}
    \item \textbf{Observance critique en 2ème ligne :}
    \begin{itemize}
        \item La 2ème ligne (IP boostés anciennement, DTG aujourd'hui) est la \textbf{dernière option accessible massivement} au Cameroun (3ème ligne = Darunavir/r + DTG + nouveaux ARV = coûteuse, complexe, rupture de stock fréquente).
        \item Échec 2ème ligne $\rightarrow$ Résistance aux INI (DTG) + NRTI $\rightarrow$ Impasse thérapeutique.
        \item DTG a une \textbf{demi-vie longue} ($\sim$ 14h) et barrière haute, ce qui pardonne \textbf{1 oubli occasionnel}, mais pas une observance $< 90\%$.
        \item \textbf{Action :} Conseil adherence intensif (thérapeute pairs, application rappel, dépistage dépression/alcool, allègement pilulier TLD 1cp/j).
    \end{itemize}
\end{enumerate}
\textbf{Conclusion :} L'analyse du génotypage guide le choix rationnel : on conserve les classes intactes (INI) et on recycle les NRTI avec un effet « fitness cost/hypersensibilisation » (M184V/TDF), le tout dans une co-formulation à barrière haute (TLD) pour sécuriser la durabilité.
\end{exoresolu}

\begin{methode}[Interpréter un algorithme de dépistage VIH (Tests rapides, ELISA, Confirmation, Fenêtre sérologique)]
\textbf{Objectif :} Savoir lire un arbre décisionnel de dépistage (Stratégie 3 tests au Cameroun), calculer la Valeur Prédictive Positive (VPP) / Négative (VPN) et comprendre la fenêtre sérologique.
\begin{enumerate}
    \item \textbf{Connaître la stratégie nationale (3 tests rapides séquentiels) :}
    \begin{enumerate}
        \item \textbf{Test 1 (Dépistage - Haute Sensibilité) :} Ex: Determine HIV-1/2, SD Bioline. Si \textbf{Négatif} $\rightarrow$ \textbf{Négatif définitif} (sauf risque récent $< 6$ sem).
        \item \textbf{Test 2 (Confirmation 1 - Haute Spécificité) :} Si Test 1 Positif. Ex: Uni-Gold, Stat-Pak. Si \textbf{Négatif} $\rightarrow$ \textbf{Discordance} $\rightarrow$ Test 3.
        \item \textbf{Test 3 (Confirmation 2 / Arbitrage) :} Ex: Vikia, Bioline 2.0. \textbf{Positif} $\rightarrow$ \textbf{Positif définitif}. \textbf{Négatif} $\rightarrow$ \textbf{Indéterminé} $\rightarrow$ Contrôle à J+14 (sérologie) ou PCR VIH (charge virale) si suspicion infection récente.
    \end{enumerate}
    \item \textbf{Fenêtre sérologique (Période de latence sérologique) :} Délai entre infection et détection des anticorps.
    \begin{itemize}
        \item \textbf{Tests rapides / ELISA 3ème gén (Ac seul) :} $\sim$ \textbf{3 à 6 semaines} (médiane 22-25 j).
        \item \textbf{Tests 4ème gén (Ag p24 + Ac) / PCR ARN VIH :} $\sim$ \textbf{10-15 jours} (PCR) / \textbf{2-3 semaines} (Ag/Ab).
    \end{itemize}
    \item \textbf{Calcul VPP / VPN (si prévalence $P$, Sensibilité $Se$, Spécificité $Sp$) :}
    \[
    \text{VPP} = \frac{Se \times P}{Se \times P + (1-Sp) \times (1-P)} \quad ; \quad \text{VPN} = \frac{Sp \times (1-P)}{(1-Se) \times P + Sp \times (1-P)}
    \]
    \textbf{Règle :} Faible prévalence $\rightarrow$ VPP faible (faux positifs nombreux) $\rightarrow$ Nécessité confirmation. Haute prévalence $\rightarrow$ VPN faible (faux négatifs possibles si fenêtre).
\end{enumerate}
\textbf{Piège BEPC :} Ne pas confondre \emph{sensibilité} (détecter les vrais +) et \emph{spécificité} (ne pas détecter les faux +). Un test « positif » en dépistage (Test 1) n'est \textbf{JAMAIS} un diagnostic définitif. Dire « Le patient est séropositif » seulement après confirmation (Test 2 + Test 3).
\end{methode}

\begin{exoresolu}[Algorithme de dépistage et fenêtre sérologique]
\textbf{Énoncé :} Dans un centre de dépistage volontaire (CDV) à Yaoundé (prévalence VIH population générale $\approx 2,7\%$), on utilise la stratégie 3 tests rapides : \textbf{Test A (Determine®)}, \textbf{Test B (Uni-Gold®)}, \textbf{Test C (Vikia®)}.
Les performances (constructeur) : Test A : $Se=100\%, Sp=99,5\%$ ; Test B : $Se=100\%, Sp=99,8\%$ ; Test C : $Se=99,9\%, Sp=99,9\%$.
Un jeune homme de 22 ans consulte 10 jours après un rapport sexuel non protégé à risque. Le Test A est \textbf{Positif}, le Test B est \textbf{Négatif}, le Test C est \textbf{Négatif}.
\begin{enumerate}
    \item Quel est le résultat final selon l'algorithme ? Quelle conduite tenir immédiatement ?
    \item Calculer la \textbf{Valeur Prédictive Positive (VPP)} du Test A seul dans cette population (Prévalence $P=0,027$). Interpréter le résultat.
    \item Expliquer pourquoi le Test A peut être positif alors que les Tests B et C sont négatifs, dans le contexte de la \textbf{fenêtre sérologique} et de la \textbf{cinétique des anticorps}.
    \item Quelle est la conduite à tenir \textbf{à 6 semaines} (J+42) si le patient revient ? Quel test demander en priorité ?
\end{enumerate}
\tcblower
\textbf{Solution détaillée :}
\begin{enumerate}
    \item \textbf{Résultat final :} \textbf{Indéterminé / Discordant (A+ / B- / C-)}.
    \textbf{Conduite immédiate :} \textbf{NE PAS ANNONCER LA SÉROPOSITIVITÉ}. Expliquer le résultat discordant. Proposer un \textbf{contrôle sérologique à 14 jours (J+14)} ou réaliser une \textbf{Charge Virale VIH (PCR ARN)} si forte suspicion clinique (syndrome primaire) ou anxiété majeure. Conseiller préservatif / PrEP en attendant.
    \item \textbf{Calcul VPP Test A :}
    $P = 0,027$ ; $Se = 1,00$ ; $Sp = 0,995$.
    \[
    \text{VPP} = \frac{1,00 \times 0,027}{1,00 \times 0,027 + (1 - 0,995) \times (1 - 0,027)} = \frac{0,027}{0,027 + 0,005 \times 0,973} = \frac{0,027}{0,027 + 0,004865} = \frac{0,027}{0,031865} \approx \mathbf{0,847} \text{ (84,7\%)}
    \]
    \textbf{Interprétation :} Même avec un excellent test ($Sp=99,5\%$), dans une population à faible prévalence (2,7\%), \textbf{1 test positif sur 6 est un FAUX POSITIF} (VPP $< 90\%$). D'où l'\textbf{obligation absolue de confirmation} par un 2ème puis 3ème test de spécificité supérieure.
    \item \textbf{Explication Discordance (Fenêtre sérologique) :}
    \begin{itemize}
        \item \textbf{Délai :} 10 jours post-exposition $\ll$ fenêtre sérologique moyenne (3-6 sem pour tests Ac).
        \item \textbf{Cinétique Ac :} Apparition séquentielle : IgM anti-gp41 $\rightarrow$ IgM anti-p24 $\rightarrow$ IgG. Les tests rapides détectent principalement IgG/IgM.
        \item \textbf{Sensibilité différentielle :} Test A (Determine) est optimisé pour une \textbf{Sensibilité maximale (100\%)} $\rightarrow$ détecte les très bas taux d'Ac naissants (IgM précoces) $\rightarrow$ \textbf{Faux positif transitoire (réactivité croisée non spécifique) OU Vrai positif ultra-précoce (séroconversion naissante)}.
        \item Tests B (Uni-Gold) et C (Vikia) ont une \textbf{Spécificité supérieure (99,8-99,9\%)} et un seuil de détection légèrement plus élevé $\rightarrow$ restent \textbf{Négatifs} car taux d'Ac encore sous seuil de détection ou absence d'Ac spécifiques.
    \end{itemize}
    \textbf{Hypothèse la plus probable :} \textbf{Faux positif du Test A} (réactivité non spécifique fréquente en début d'infection ou facteur rhumatoïde) \textbf{OU} tout début de séroconversion (Test A + en premier).
    \item \textbf{Conduite à J+42 (6 semaines) :}
    \begin{itemize}
        \item \textbf{Délai suffisant} pour exclure fenêtre sérologique tests 3ème génération (Ac).
        \item \textbf{Refaire l'algorithme complet 3 tests rapides (A, B, C).}
        \item \textbf{Si A+ / B+ / C+ $\rightarrow$ Positif définitif $\rightarrow$ Prise en charge ARV (Treat All).}
        \item \textbf{Si A- / B- / C- $\rightarrow$ Négatif définitif $\rightarrow$ Rassurer, proposer PrEP si risque persistant.}
        \item \textbf{Priorité :} \textbf{Test 4ème génération (Ag p24 + Ac) ou PCR ARN VIH} si disponible (détection plus précoce, lève l'ambiguïté plus tôt), mais à 6 sem, tests rapides 3ème gen suffisent.
    \end{itemize}
\end{enumerate}
\textbf{Conclusion :} La discordance en phase de fenêtre sérologique est une situation classique. La rigueur algorithmique (3 tests) et le respect du délai de fenêtre évitent les annonces erronées (faux positifs) ou les diagnostics manqués (faux négatifs précoces).
\end{exoresolu}

\begin{methode}[Analyser des données épidémiologiques (Prévalence, Incidence, Cascade de prise en charge 95-95-95)]
\textbf{Objectif :} Calculer et interpréter les indicateurs de suivi de l'épidémie (Prévalence, Incidence, Taux de couverture ARV, Suppression virale) et les objectifs ONUSIDA 95-95-95.
\begin{enumerate}
    \item \textbf{Définitions clés :}
    \begin{itemize}
        \item \textbf{Prévalence (P) :} $\frac{\text{Nombre de cas VIH+ à un instant } t}{\text{Population totale à } t}$. Mesure la \textbf{charge} de la maladie.
        \item \textbf{Incidence (I) :} $\frac{\text{Nouveaux cas sur une période}}{\text{Population saine à risque (personnes-années)}}$. Mesure la \textbf{vitesse} de propagation. Relation approx : $P \approx I \times D$ (Durée moyenne de vie avec VIH).
        \item \textbf{Cascade 95-95-95 (Objectifs 2025/2030) :}
        \begin{enumerate}
            \item 95\% des PVVIH \textbf{connaissent leur statut} (Dépistage).
            \item 95\% des diagnostiqués \textbf{sont sous ARV} (Initiation/Retention).
            \item 95\% des sous ARV \textbf{ont CV indétectable} (Observance/Efficacité).
        \end{enumerate}
        \textbf{Impact global :} $0,95 \times 0,95 \times 0,95 \approx 86\%$ de tous les PVVIH avec CV indétectable $\rightarrow$ Contrôle épidémie (I=I).
    \end{itemize}
    \item \textbf{Calculs types BEPC :}
    \begin{itemize}
        \item Taux de couverture ARV = $\frac{\text{Nbre sous ARV}}{\text{Nbre estimé PVVIH}} \times 100$.
        \item Taux suppression virale = $\frac{\text{Nbre CV indétectable}}{\text{Nbre sous ARV}} \times 100$.
        \item \textbf{Gap (Ecart) :} Identifier où se perd le plus de patients (ex: « 1er 95 » faible = problème dépistage ; « 3ème 95 » faible = problème observance/rupture stock).
    \end{itemize}
    \item \textbf{Contexte Cameroun (Données PNLS/ONUSIDA 2023 approx) :} Prévalence adultes 15-49 ans $\approx 2,7\%$ ($\sim$ 500 000 PVVIH). Couverture ARV $\approx 85\%$ (1er 95 $\approx$ 90\%, 2ème 95 $\approx$ 95\%, 3ème 95 $\approx$ 90\%). Défis : Hommes, Adolescents/Jeunes (AGYW), Populations Clés (HSH, TS, UD), Enfants (diagnostic précoce PCR).
\end{enumerate}
\textbf{Reflexe :} Unités : Prévalence en \%, Incidence en /1000 pers-an. Distinguer \emph{nombre absolu} et \emph{taux}.
\end{methode}

\begin{exoresolu}[Cascade de prise en charge et objectifs 95-95-95 au Cameroun]
\textbf{Énoncé :} Selon le rapport PNLS 2023, sur une population adulte de 15-49 ans de \textbf{10 000 000} d'habitants, on estime \textbf{270 000} PVVIH. Le programme rapporte : \textbf{243 000} PVVIH connaissent leur statut ; \textbf{230 850} sont sous ARV ; \textbf{207 765} sous ARV ont une charge virale indétectable ($< 1000$ cp/mL).
\begin{enumerate}
    \item Calculer la \textbf{prévalence} du VIH chez les adultes 15-49 ans (en \%). 
    \item Calculer les trois « 95 » de la cascade (en \%) pour le Cameroun en 2023. Les objectifs ONUSIDA sont-ils atteints ?
    \item Calculer le \textbf{taux de suppression virale global} (proportion de \textbf{Tous} les PVVIH estimées qui sont viralement supprimées).
    \item Identifier le « maillon faible » de la cascade. Proposer \textbf{deux interventions ciblées} (une biomédicale, une communautaire) pour améliorer ce maillon, en lien avec le contexte camerounais.
\end{enumerate}
\tcblower
\textbf{Solution détaillée :}
\begin{enumerate}
    \item \textbf{Prévalence :}
    \[
    P = \frac{270\,000}{10\,000\,000} \times 100 = \mathbf{2,7\%}
    \]
    \item \textbf{Calcul des 95 :}
    \begin{enumerate}
        \item \textbf{1er 95 (Connaissance statut) :} $\frac{243\,000}{270\,000} \times 100 = \mathbf{90,0\%}$. \textbf{Objectif 95\% NON atteint} (Gap : 27 000 PVVIH ignorent leur statut).
        \item \textbf{2ème 95 (Sous ARV / Diagnostiqués) :} $\frac{230\,850}{243\,000} \times 100 = \mathbf{95,0\%}$. \textbf{Objectif 95\% ATTEINT}.
        \item \textbf{3ème 95 (Suppression virale / Sous ARV) :} $\frac{207\,765}{230\,850} \times 100 = \mathbf{90,0\%}$. \textbf{Objectif 95\% NON atteint} (Gap : 23 085 sous ARV non supprimés).
    \end{enumerate}
    \item \textbf{Taux de suppression virale global (Cascade complète) :}
    \[
    \frac{207\,765}{270\,000} \times 100 = \mathbf{76,9\%}
    \]
    (Equivaut au produit des 3 étapes : $0,90 \times 0,95 \times 0,90 = 0,7695$).
    \textbf{Interprétation :} Seuls $\sim 77\%$ des PVVIH camerounais sont viralement supprimés. $\sim 23\%$ (62 000 personnes) ont une CV détectable $\rightarrow$ Risque transmission, morbidité, mortalité.
    \item \textbf{Maillon faible :} Deux maillons à \textbf{90\%} (1er et 3ème 95). Le \textbf{1er 95 (Dépistage)} est souvent prioritaire car « on ne soigne pas ce qu'on ne connaît pas ». Le \textbf{3ème 95 (Observance/Suppression)} est critique pour l'impact transmission (I=I).
    \textbf{Interventions proposées :}
    \begin{itemize}
        \item \textbf{Pour le 1er 95 (Dépistage) - Biomédicale :} \textbf{Autotest VIH (HIV Self-Testing)} distribué gratuitement dans les pharmacies, universités, lieux de travail, via pairs éducateurs (AGYW, HSH). Levée barrière confidentialité/stigmatisation. Confirmation obligatoire en structure si réactif.
        \item \textbf{Pour le 3ème 95 (Observance) - Communautaire :} \textbf{Groupes de soutien « Adherence Clubs » / Groupes de pairs (Pair Educators)} au niveau des CTA et communautés. Renforcement counseling observance (Outils : boîtes piluliers, applications SMS, soutien psychosocial). \textbf{Modèle de distribution communautaire des ARV (DAC - Differentiated Service Delivery)} : Recharge 3-6 mois hors hôpital pour patients stables (réduit pertes de vue, coûts transport, stigmatisation).
    \end{itemize}
\end{enumerate}
\textbf{Conclusion :} L'analyse de la cascade permet de passer du « nombre de patients sous ARV » à l'« impact épidémiologique ». Le Cameroun performe bien sur l'initiation (2ème 95) mais doit combler les gaps du dépistage (1er 95) et de la rétention/suppression virale (3ème 95) pour atteindre l'élimination du Sida comme menace de santé publique d'ici 2030.
\end{exoresolu}

\begin{attention}[Les 5 erreurs qui coûtent des points au BEPC]
\begin{enumerate}
    \item \textbf{Confondre « Transcription » et « Transcription Inverse » :} Écrire « Le VIH transcrit son ARN en ADN » sans préciser \textbf{« Transcription Inverse »} (catalysée par la \textbf{Transcriptase Inverse}). La transcription classique (ADN $\rightarrow$ ARN) est faite par l'hôte. \textbf{Perte de points :} Vocabulaire imprécis, confusion dogme central.
    \item \textbf{Dire que les ARV « tuent le virus » ou « guérissent le Sida » :} Les ARV sont \textbf{virostatiques} (bloquent la réplication), pas virucides. Ils ne éliminent pas les \textbf{réservoirs} (ADN proviral intégré). Arrêt traitement = Rebond viral. \textbf{Correction :} « Les ARV suppriment la réplication virale (CV indétectable), restaurent l'immunité (CD4), permettent une vie normale, mais ne guérissent pas l'infection. »
    \item \textbf{Inverser les courbes CD4 et Charge Virale en phase SIDA :} Croire que les CD4 remontent ou que la CV baisse en phase SIDA sans traitement. \textbf{Règle :} Phase SIDA = \textbf{CD4 $\downarrow \downarrow$ ($< 200$)} ET \textbf{CV $\uparrow \uparrow$}. L'effondrement immunitaire \textbf{permet} la réplication massive.
    \item \textbf{Oublier la « Fenêtre Sérologique » dans l'interprétation d'un test négatif :} Conclure « Le patient n'est pas infecté » après un test négatif réalisé 10 jours post-risque. \textbf{Obligatoire :} Préciser « Test négatif \textbf{mais non concluant car réalisé dans la fenêtre sérologique}. Contrôle à 6 semaines (ou 3 mois selon tests) indispensable. »
    \item \textbf{Confondre Prophylaxie Pré-Exposition (PrEP) et Post-Exposition (PEP) :} Présenter la PrEP comme un traitement d'urgence après un rapport à risque (c'est la PEP), ou la PEP comme un traitement continu (c'est la PrEP). \textbf{Rappel :} \textbf{PEP} = Urgence ($\le 72$h, 28 jours). \textbf{PrEP} = Continue (quotidien ou à la demande), avant exposition répétée.
\end{enumerate}
\end{attention}

\newpage

\coursec{Exercices}

\courssub{Je vérifie mes connaissances}

\begin{exercice}[QCM : le VIH et son cycle]
Pour chaque question, une seule réponse est exacte. Recopie la lettre de la bonne réponse.
\begin{enumerate}
\item Le VIH est :
\begin{enumerate}
\item une bactérie ;
\item un virus ;
\item un champignon ;
\item un parasite.
\end{enumerate}
\item Le VIH s'attaque principalement :
\begin{enumerate}
\item aux globules rouges ;
\item aux plaquettes ;
\item aux lymphocytes T4 ;
\item aux neurones.
\end{enumerate}
\item L'information génétique du VIH est portée par :
\begin{enumerate}
\item l'ADN ;
\item l'ARN ;
\item les protéines de l'enveloppe ;
\item les lipides.
\end{enumerate}
\item Les médicaments antirétroviraux (ARV) permettent :
\begin{enumerate}
\item de guérir définitivement du Sida ;
\item de bloquer la multiplication du virus ;
\item de tuer tous les virus en une semaine ;
\item de remplacer la vaccination.
\end{enumerate}
\end{enumerate}
\end{exercice}

\begin{exercice}[Vrai ou faux justifié]
Pour chaque affirmation, réponds par Vrai ou Faux, puis justifie ta réponse en une ou deux phrases.
\begin{enumerate}
\item Le VIH pénètre dans le lymphocyte T4 en se fixant sur un récepteur de la membrane cellulaire.
\item On peut attraper le VIH en serrant la main d'une personne séropositive.
\item Une personne séropositive est forcément malade du Sida.
\item Le préservatif est un moyen efficace de prévention de la transmission sexuelle du VIH.
\end{enumerate}
\end{exercice}

\begin{exercice}[Définitions]
Définis en une phrase chacun des termes suivants : \cle{VIH}, \cle{Sida}, \cle{séropositif}, \cle{antirétroviral (ARV)}, \cle{lymphocyte T4}.
\end{exercice}

\begin{exercice}[Ordre des étapes de la multiplication du VIH]
Voici, dans le désordre, les étapes de la multiplication du VIH dans un lymphocyte T4 :
\begin{center}
a) Libération des nouveaux virus \quad b) Fixation du virus sur la cellule \quad c) Fabrication des protéines virales \quad d) Entrée du virus dans la cellule \quad e) Assemblage des nouveaux virus \quad f) Transformation de l'ARN viral en ADN
\end{center}
Remets ces étapes dans l'ordre chronologique et recopie la suite de lettres obtenue.
\end{exercice}

\courssub{Je m'entraîne}

\begin{exercice}[Schéma du cycle viral et action des ARV]
On donne le schéma simplifié du cycle de multiplication du VIH ci-dessous, où les étapes sont numérotées de \textcircled{1} à \textcircled{4}.

\begin{popfigure}
\begin{tikzpicture}[font=\small, >=stealth]
\node[draw=popblue, fill=popblueL, rounded corners, minimum width=3.2cm, minimum height=1cm] (fix) at (0,0) {Fixation et entrée};
\node[draw=poporange, fill=poporangeL, rounded corners, minimum width=3.2cm, minimum height=1cm] (adn) at (4.6,0) {ARN $\rightarrow$ ADN};
\node[draw=popgreen, fill=popgreenL, rounded corners, minimum width=3.2cm, minimum height=1cm] (fab) at (9.2,0) {Fabrication des protéines};
\node[draw=poppurple, fill=poppurpleL, rounded corners, minimum width=3.2cm, minimum height=1cm] (lib) at (13.8,0) {Assemblage et libération};
\draw[->, thick, popdark] (fix) -- (adn) node[midway, above]{\textcircled{1}};
\draw[->, thick, popdark] (adn) -- (fab) node[midway, above]{\textcircled{2}};
\draw[->, thick, popdark] (fab) -- (lib) node[midway, above]{\textcircled{3}};
\node[above=0.15cm of fix]{\textcircled{4}};
\end{tikzpicture}
\legende{Cycle simplifié de multiplication du VIH dans un lymphocyte T4. Les numéros \textcircled{1} à \textcircled{4} désignent des étapes ou des transitions.}
\end{popfigure}

Quatre médicaments ARV agissent chacun à une étape précise :
\begin{itemize}
\item ARV A : empêche la fixation du virus sur le lymphocyte ;
\item ARV B : empêche la transformation de l'ARN viral en ADN ;
\item ARV C : empêche la fabrication des protéines virales ;
\item ARV D : empêche l'assemblage et la libération des nouveaux virus.
\end{itemize}
\begin{enumerate}
\item Associe chaque ARV (A, B, C, D) au numéro de l'étape qu'il bloque.
\item Un malade prend l'ARV B seul. Explique pourquoi le médecin lui prescrit en réalité une association de plusieurs ARV.
\end{enumerate}
\end{exercice}

\begin{exercice}[Lecture d'une courbe : évolution des lymphocytes T4]
Chez une personne infectée par le VIH et non traitée, on a mesuré le nombre de lymphocytes T4 dans le sang à différentes dates :

\begin{center}
\begin{tabularx}{\linewidth}{|l|X|X|X|X|}
\hline
\rowcolor{popblueL} Temps après l'infection (années) & 0 & 2 & 5 & 8 \\
\hline Nombre de T4 (par mm$^3$ de sang) & 1000 & 800 & 450 & 150 \\
\hline
\end{tabularx}
\end{center}

On rappelle : chez une personne saine, le nombre de T4 est compris entre 800 et \num{1200} par mm$^3$ ; on parle de Sida lorsque ce nombre descend en dessous de 200 par mm$^3$.
\begin{enumerate}
\item Trace la courbe du nombre de T4 en fonction du temps (1 cm pour 1 an ; 1 cm pour 100 T4 par mm$^3$).
\item À quelle date le malade entre-t-il dans la phase Sida ? Justifie.
\item Calcule la baisse moyenne du nombre de T4 par an entre la 2\textsuperscript{e} et la 5\textsuperscript{e} année.
\item Explique pourquoi la baisse des T4 rend l'organisme sensible aux maladies dites « opportunistes ».
\end{enumerate}
\end{exercice}

\begin{exercice}[Choisir le bon comportement : situations de la vie courante]
Pour chacune des situations suivantes, indique s'il existe un risque de transmission du VIH et justifie ta réponse.
\begin{enumerate}
\item Partager un repas avec une personne séropositive.
\item Utiliser une lame de rasoir ayant servi à une personne séropositive.
\item Avoir un rapport sexuel sans préservatif avec un partenaire séropositif.
\item Être piqué par un moustique qui a piqué une personne séropositive.
\item Une mère séropositive allaite son bébé sans traitement.
\end{enumerate}
\end{exercice}

\begin{exercice}[Rédaction argumentée]
Un ami te dit : « Je ne prends pas mes ARV le week-end pour faire une pause et reposer mon foie. »
En utilisant tes connaissances sur la multiplication rapide du VIH et sur le rôle des ARV, rédige un texte de 10 à 15 lignes pour expliquer à ton ami pourquoi cette pratique est dangereuse (reprise de la multiplication du virus, baisse des défenses, risque que le médicament devienne inefficace).
\end{exercice}

\courssub{Je me prépare au BEPC}

\begin{exercice}[Type BEPC — Suivi d'un malade au Centre de Santé de Bafoussam]
Au Centre de Santé de Bafoussam, un patient de 30 ans, séropositif, est suivi pendant 5 ans. On mesure chaque année son nombre de lymphocytes T4 et la quantité de virus dans le sang (charge virale). Les résultats sont rassemblés dans le tableau ci-dessous.

\begin{center}
\begin{tabularx}{\linewidth}{|l|X|X|X|X|X|X|}
\hline
\rowcolor{popblueL} Année & 0 & 1 & 2 & 3 & 4 & 5 \\
\hline T4 (par mm$^3$) & 950 & 850 & 700 & 520 & 350 & 180 \\
\hline Charge virale (millions de virus par mL) & 0,05 & 0,1 & 0,4 & 1,2 & 3,5 & 8,0 \\
\hline
\end{tabularx}
\end{center}

\begin{enumerate}
\item Définis les termes VIH et charge virale.
\item Décris l'évolution du nombre de T4 et celle de la charge virale entre l'année 0 et l'année 5.
\item Compare les deux évolutions : que remarques-tu ? Propose une explication en t'aidant du rôle des lymphocytes T4 et du cycle de multiplication du VIH.
\item À partir de quelle année peut-on considérer que le patient est entré en phase Sida (seuil : 200 T4 par mm$^3$) ? Justifie.
\item Le médecin décide de mettre le patient sous ARV à l'année 5. Explique le mode d'action des ARV et le résultat attendu sur la charge virale et le nombre de T4.
\end{enumerate}
\end{exercice}

\begin{exercice}[Type BEPC — Prévention du VIH dans un lycée de Yaoundé]
Dans le cadre de la semaine de la santé, le club de sciences d'un lycée de Yaoundé organise une campagne d'information sur le VIH/Sida. Voici trois affirmations entendues parmi les élèves :
\begin{itemize}
\item Élève 1 : « Le VIH se transmet par la salive et les piqûres de moustiques. »
\item Élève 2 : « Une personne séropositive sous traitement ARV bien suivi peut vivre longtemps. »
\item Élève 3 : « Le dépistage ne sert à rien si l'on se sent en bonne santé. »
\end{itemize}
\begin{enumerate}
\item Pour chaque élève, dis si l'affirmation est correcte ou non, en justifiant scientifiquement.
\item Cite les trois principales voies de transmission du VIH.
\item Cite trois moyens de prévention de la transmission du VIH et explique le fonctionnement de l'un d'entre eux.
\item Une femme enceinte séropositive se rend à l'hôpital. Explique pourquoi un traitement ARV pendant la grossesse protège son bébé.
\item Rédige un slogan de 2 lignes maximum pour la campagne de prévention du lycée.
\end{enumerate}
\end{exercice}

\begin{exercice}[Type BEPC — Problème de synthèse : multiplication du VIH et traitement]
On estime que, chez un malade non traité, le VIH peut produire jusqu'à $10^{10}$ (dix milliards) nouveaux virus par jour dans l'organisme.
\begin{enumerate}
\item Calcule le nombre de virus produits en une semaine (7 jours) chez ce malade non traité. Donne le résultat en écriture scientifique.
\item Rappelle les quatre grandes étapes de la multiplication du VIH dans un lymphocyte T4.
\item En t'aidant de la question 1, explique pourquoi l'infection par le VIH finit, sans traitement, par épuiser les lymphocytes T4.
\item Sous traitement ARV bien suivi, la production de nouveaux virus devient presque nulle. En déduire l'intérêt de prendre les ARV tous les jours, sans interruption.
\item Un malade dit : « Puisque les ARV font baisser le nombre de virus, je suis guéri, je peux arrêter. » Corrige cette erreur en deux ou trois phrases argumentées.
\end{enumerate}
\end{exercice}



\newpage

\coursec{Corrigés des exercices}

\begin{corrige}[Exercice 1 — QCM : le VIH et son cycle]
\textbf{Rappel de l'énoncé :} Choisir la bonne réponse parmi les propositions pour chaque question : nature du VIH, cellule cible, matériel génétique, action des ARV.
\tcblower
\begin{enumerate}
\item \textbf{b)} Le VIH (Virus de l'Immunodéficience Humaine) est un \cle{virus} : il est beaucoup plus petit qu'une bactérie et ne peut se multiplier qu'à l'intérieur d'une cellule vivante.
\item \textbf{c)} Le VIH s'attaque principalement aux \cle{lymphocytes T4} (ou lymphocytes T CD4+), les cellules qui commandent les défenses immunitaires.
\item \textbf{b)} L'information génétique du VIH est portée par l'\cle{ARN} : c'est un virus à ARN, appelé « rétrovirus ».
\item \textbf{b)} Les ARV \textbf{bloquent la multiplication du virus} mais ne l'éliminent pas totalement de l'organisme : ils ne guérissent pas le Sida, ils permettent de vivre longtemps avec le VIH.
\end{enumerate}
\end{corrige}

\begin{corrige}[Exercice 2 — Vrai ou faux justifié]
\textbf{Rappel de l'énoncé :} Dire si chaque affirmation est vraie ou fausse et justifier scientifiquement : mode de pénétration, modes de transmission, distinction séropositif/Sida, efficacité du préservatif.
\tcblower
\begin{enumerate}
\item \textbf{Vrai.} Le virus se fixe d'abord sur un récepteur spécifique (CD4 et co-récepteur CCR5/CXCR4) situé à la surface de la membrane du lymphocyte T4, puis il fusionne avec la membrane pour entrer dans la cellule.
\item \textbf{Faux.} Le VIH ne se transmet pas par les contacts de la vie courante (serrage de mains, salive, sueur, larmes, piqûres de moustiques) ; il se transmet par le sang (objets contaminés, transfusions non sécurisées), les rapports sexuels non protégés et de la mère à l'enfant (grossesse, accouchement, allaitement).
\item \textbf{Faux.} Une personne \cle{séropositive} est porteuse du virus (présence d'anticorps anti-VIH dans le sang) mais peut rester en bonne santé pendant des années (phase asymptomatique) ; le Sida est la phase la plus avancée de la maladie, marquée par l'effondrement des défenses immunitaires.
\item \textbf{Vrai.} Le préservatif (masculin ou féminin) constitue une barrière physique étanche qui empêche le contact entre le virus présent dans les sécrétions sexuelles et la muqueuse du partenaire lors des rapports sexuels.
\end{enumerate}
\textbf{Point de vigilance :} ne pas confondre « séropositif » (porteur du virus, stade asymptomatique possible) et « malade du Sida » (phase avancée avec maladies opportunistes) ; c'est l'erreur la plus fréquente au BEPC.
\end{corrige}

\begin{corrige}[Exercice 3 — Définitions]
\textbf{Rappel de l'énoncé :} Définir précisément les termes suivants : VIH, Sida, Séropositif, Antirétroviral (ARV), Lymphocyte T4.
\tcblower
\begin{itemize}
\item \cle{VIH} : Virus de l'Immunodéficience Humaine, virus à ARN (rétrovirus) qui infecte et détruit les lymphocytes T4, affaiblissant progressivement les défenses de l'organisme.
\item \cle{Sida} : Syndrome d'Immunodéficience Acquise, phase terminale de l'infection par le VIH, où les défenses immunitaires sont effondrées (T4 < 200/mm$^3$) et où apparaissent des maladies opportunistes graves.
\item \cle{Séropositif} : se dit d'une personne dont le sang contient des anticorps anti-VIH (détectés par un test de dépistage), signe qu'elle est infectée par le virus.
\item \cle{Antirétroviral (ARV)} : médicament qui bloque une étape spécifique de la multiplication du VIH (fixation, transcription inverse, intégration, protéase) et ralentit ainsi l'évolution de la maladie.
\item \cle{Lymphocyte T4} : globule blanc (lymphocyte T auxiliaire CD4+) qui dirige et coordonne les réactions de défense de l'organisme (immunité humorale et cellulaire) et qui est la cellule cible principale du VIH.
\end{itemize}
\end{corrige}

\begin{corrige}[Exercice 4 — Ordre des étapes de la multiplication du VIH]
\textbf{Rappel de l'énoncé :} Remettre dans l'ordre chronologique les étapes suivantes : a) Libération, b) Fixation, c) Fabrication protéines, d) Entrée, e) Assemblage, f) Transcription inverse (ARN $\to$ ADN).
\tcblower
L'ordre chronologique est : \textbf{b -- d -- f -- c -- e -- a}.
\begin{enumerate}
\item \textbf{b)} Fixation du virus sur les récepteurs (CD4 et co-récepteur) du lymphocyte T4 ;
\item \textbf{d)} Fusion de l'enveloppe virale avec la membrane cellulaire et entrée du capside (contenant l'ARN viral et les enzymes) dans le cytoplasme ;
\item \textbf{f)} Transformation de l'ARN viral en ADN double brin (grâce à l'enzyme virale \cle{transcriptase inverse}) : c'est la \cle{transcription inverse} ;
\item \textbf{c)} Intégration de l'ADN viral dans le génome de la cellule hôte, puis transcription/translation : fabrication des protéines virales par la machinerie cellulaire ;
\item \textbf{e)} Assemblage des nouveaux virions (capside, ARN, enzymes) au niveau de la membrane cellulaire ;
\item \textbf{a)} Libération des nouveaux virus par bourgeonnement (budding), qui pourront infecter d'autres lymphocytes T4 (la cellule meurt souvent au passage).
\end{enumerate}
\textbf{Point de vigilance :} la fixation précède toujours l'entrée, et la transcription inverse (ARN $\to$ ADN) doit avoir lieu avant la fabrication des protéines virales car le génome viral doit être sous forme d'ADN pour être lu par la cellule.
\end{corrige}

\begin{corrige}[Exercice 5 — Schéma du cycle viral et action des ARV]
\textbf{Rappel de l'énoncé :} À partir d'un schéma schématisant le cycle du VIH (numéroté 1 à 4 sur les flèches/étapes), identifier l'étape bloquée par quatre ARV (A, B, C, D) et expliquer l'intérêt de la trithérapie.
\tcblower
\begin{enumerate}
\item Selon la numérotation du schéma fourni : l'ARV A bloque l'étape \textcircled{4} (fixation et entrée) ; l'ARV B bloque l'étape \textcircled{1} (transcription inverse : ARN $\to$ ADN) ; l'ARV C bloque l'étape \textcircled{2} (fabrication des protéines / maturation) ; l'ARV D bloque l'étape \textcircled{3} (assemblage et libération / bourgeonnement).
\item Le VIH se multiplie très vite (environ $10^{10}$ virus/jour) et mute fréquemment. S'il n'est attaqué que par un seul médicament (monothérapie), des mutants résistants sont rapidement sélectionnés. En associant plusieurs ARV (\cle{trithérapie}) qui bloquent des étapes \emph{différentes} du cycle, on empêche beaucoup plus efficacement la multiplication du virus et on réduit drastiquement le risque qu'un virus devienne simultanément résistant à tous les médicaments du traitement.
\end{enumerate}
\textbf{Point de vigilance :} bien lire la position des numéros sur le schéma fourni : \textcircled{4} est au-dessus de la première case (fixation et entrée), les numéros \textcircled{1}, \textcircled{2}, \textcircled{3} sont sur les flèches de transition entre les cases. La numérotation peut varier selon les manuels ; il faut se fier au document de l'examen.
\end{corrige}

\begin{corrige}[Exercice 6 — Lecture d'une courbe : évolution des lymphocytes T4]
\textbf{Rappel de l'énoncé :} Tracer la courbe d'évolution du taux de lymphocytes T4 (données : année 0 : 1000 ; an 2 : 800 ; an 5 : 450 ; an 8 : 150 ; échelle 1 cm/an, 1 cm/100 T4/mm$^3$). Déterminer quand le seuil de Sida (200 T4/mm$^3$) est franchi. Calculer la vitesse moyenne de baisse entre an 2 et an 5. Expliquer les conséquences.
\tcblower
\begin{enumerate}
\item On place les points $(0\,;\,1000)$, $(2\,;\,800)$, $(5\,;\,450)$ et $(8\,;\,150)$ en respectant l'échelle (1 cm pour 1 an en abscisse ; 1 cm pour 100 T4 par mm$^3$ en ordonnée), puis on les relie par une courbe lisse : la courbe est décroissante, de plus en plus rapide (concave vers le bas).
\item À la 8\textsuperscript{e} année, le nombre de T4 est de 150 par mm$^3$, donc inférieur au seuil de 200 par mm$^3$ : le malade est entré en phase Sida entre la 5\textsuperscript{e} et la 8\textsuperscript{e} année, et le Sida est biologiquement constaté à l'année 8.
\item Entre la 2\textsuperscript{e} et la 5\textsuperscript{e} année, la baisse totale est de $800 - 450 = 350$ T4 par mm$^3$ en 3 ans, soit une vitesse moyenne de baisse de $\dfrac{350}{3} \approx 117$ T4 par mm$^3$ et par an.
\item Les lymphocytes T4 dirigent les défenses immunitaires (activation des lymphocytes B, des lymphocytes T8, des macrophages). Quand leur nombre s'effondre (< 200/mm$^3$), l'organisme ne peut plus lutter contre des microbes habituellement sans gravité : des maladies dites « \cle{opportunistes} » (tuberculose, pneumonie à \emph{Pneumocystis jirovecii}, diarrhées persistantes, candidose, toxoplasmose, etc.) se développent alors.
\end{enumerate}
\textbf{Point de vigilance :} toujours comparer la valeur mesurée au seuil donné dans l'énoncé (200 T4 par mm$^3$) et citer ce seuil explicitement dans la justification. Pour la vitesse, indiquer l'unité (T4/mm$^3$/an).
\end{corrige}

\begin{corrige}[Exercice 7 — Choisir le bon comportement : situations de la vie courante]
\textbf{Rappel de l'énoncé :} Évaluer le risque de transmission du VIH pour 5 situations : boire dans le même verre, partager un rasoir, rapport sexuel non protégé, piqûre de moustique, allaitement par mère séropositive.
\tcblower
\begin{enumerate}
\item \textbf{Pas de risque} : le VIH ne se transmet ni par la salive (concentration virale trop faible, enzymes inhibitrices), ni par la nourriture, ni par les contacts de la vie courante (baiser sur la joue, mains, ustensiles).
\item \textbf{Risque} : la lame de rasoir peut être souillée par du sang contaminé (micro-coupures) ; or le VIH se transmet efficacement par le sang (voie parentérale). Ne jamais partager d'objets piquants/tranchants (rasoirs, aiguilles, ciseaux, brosses à dents).
\item \textbf{Risque élevé} : les rapports sexuels non protégés (vaginaux, anaux, oraux) sont la principale voie de transmission du VIH dans le monde (contact avec sperme, sécrétions vaginales, sang menstruel, lésions muqueuses).
\item \textbf{Pas de risque} : le VIH ne se multiplie pas dans le moustique (pas de récepteurs CD4, digestion du virus) et n'est pas transmis par ses piqûres (le moustique n'injecte pas le sang de la personne précédente, seulement sa salive) ; contrairement au paludisme (plasmodium) ou aux arbovirus.
\item \textbf{Risque} : le VIH peut passer de la mère à l'enfant par le lait maternel (contenance en virus et cellules infectées) ; un traitement ARV efficace de la mère (charge virale indétectable) et un suivi médical (allaitement exclusif court ou lait de substitution si possible) réduisent fortement ce risque (prévention de la transmission mère-enfant, PTME).
\end{enumerate}
\end{corrige}

\begin{corrige}[Exercice 8 — Rédaction argumentée]
\textbf{Rappel de l'énoncé :} Un ami séropositif sous ARV veut arrêter son traitement le week-end car il se sent bien. Rédiger un argumentaire (10-15 lignes) pour le convaincre de prendre son traitement tous les jours.
\tcblower
\textbf{Éléments attendus (10 à 15 lignes) :}
Le VIH se multiplie en permanence dans l'organisme : il peut produire jusqu'à dix milliards ($10^{10}$) de nouveaux virus par jour chez un malade non traité. Les ARV bloquent cette multiplication, mais seulement s'ils sont présents en quantité suffisante dans le sang (concentration plasmatique au-dessus du seuil inhibiteur), ce qui impose une prise quotidienne rigoureuse, sans interruption. Si mon ami arrête ses ARV chaque week-end, la quantité de médicament dans son sang diminue rapidement (demi-vie courte) : le virus recommence alors à se multiplier activement et à détruire ses lymphocytes T4. Ses défenses immunitaires baissent et il s'expose à nouveau aux maladies opportunistes. De plus, pendant ces pauses thérapeutiques, la réplication virale non contrôlée favorise l'apparition de mutations : le virus peut se transformer et devenir \cle{résistant} aux ARV de sa trithérapie ; le traitement risque alors de ne plus être efficace du tout, limitant les options thérapeutiques futures. Il doit donc prendre ses médicaments tous les jours, à heure fixe, et parler de ses inquiétudes (fatigue, effets hépatiques) à son médecin, qui seul peut adapter le traitement (changer de molécules, ajuster les doses).

\textbf{Point de vigilance :} la rédaction doit mobiliser trois idées obligatoires : 1) reprise immédiate de la multiplication du virus à l'arrêt du traitement, 2) destruction des lymphocytes T4 et baisse des défenses immunitaires, 3) risque majeur de sélection de mutants résistants (échec thérapeutique définitif).
\end{corrige}

\begin{corrige}[Exercice 9 — Type BEPC : suivi d'un malade au Centre de Santé de Bafoussam]
\textbf{Rappel de l'énoncé :} Tableau de suivi (Années 0 à 5) : T4 (mm$^3$) et Charge virale (millions/mL). Définir VIH et charge virale. Décrire les évolutions. Comparer. Déterminer entrée en phase Sida (seuil 200 T4). Expliquer action ARV et résultats attendus.
\tcblower
\begin{enumerate}
\item Le \cle{VIH} (Virus de l'Immunodéficience Humaine) est un virus à ARN (rétrovirus) qui infecte et détruit spécifiquement les lymphocytes T4 (CD4+). La \cle{charge virale} est la quantité de particules virales (ARN VIH) présente dans un volume de sang, exprimée ici en millions de copies (ou virus) par mL ; elle reflète l'activité de réplication du virus.
\item Entre l'année 0 et l'année 5, le nombre de lymphocytes T4 diminue régulièrement et fortement, passant de 950 à 180 par mm$^3$ (baisse de $950 - 180 = 770$ T4/mm$^3$). Dans le même temps, la charge virale augmente de manière exponentielle, passant de 0,05 à 8,0 millions de copies par mL.
\item Les deux évolutions sont inversement corrélées : plus la charge virale augmente, plus le nombre de T4 diminue. \textbf{Explication physiopathologique :} le VIH se multiplie \emph{dans} les lymphocytes T4 (cycle de réplication) et les détruit lors de la libération des nouveaux virions (lyse cellulaire) ; plus il y a de virus produits (charge virale élevée), plus les lymphocytes T4 sont infectés et détruits, et moins l'organisme peut assurer sa défense immunitaire.
\item À l'année 5, le nombre de T4 (180 par mm$^3$) est inférieur au seuil diagnostique de 200 par mm$^3$ : le patient est donc entré en phase Sida (stade 3 ou 4 OMS) à l'année 5 (entre la 4\textsuperscript{e} et la 5\textsuperscript{e} année).
\item Les ARV bloquent des étapes clés de la multiplication du VIH (inhibiteurs d'entrée, de la transcriptase inverse, de l'intégrase, de la protéase). \textbf{Résultats attendus sous traitement efficace (observance parfaite) :} la charge virale diminue fortement (idéalement devient indétectable < 50 copies/mL en 3-6 mois), la destruction des T4 s'arrête et leur nombre remonte progressivement (reconstitution immunitaire), ce qui redonne au patient des défenses immunitaires efficaces et fait régresser les maladies opportunistes.
\end{enumerate}
\textbf{Barème indicatif (sur 10) :} définitions 2 pts ; description des deux évolutions avec valeurs chiffrées 2 pts ; comparaison et explication physiopathologique 2 pts ; seuil du Sida justifié par la valeur 2 pts ; mode d'action des ARV et résultats attendus 2 pts.

\textbf{Points de vigilance :} citer les valeurs numériques du tableau dans la description ; la comparaison doit dégager le lien de cause à effet (le virus détruit les T4) et pas seulement constater que les courbes « vont en sens inverse ».
\end{corrige}

\begin{corrige}[Exercice 10 — Type BEPC : prévention du VIH dans un lycée de Yaoundé]
\textbf{Rappel de l'énoncé :} Analyser 3 affirmations d'élèves. Citer 3 voies de transmission. Citer 2 moyens de prévention + fonctionnement. Expliquer protection mère-enfant par ARV. Proposer un slogan.
\tcblower
\begin{enumerate}
\item \textbf{Élève 1 : affirmation fausse.} Le VIH ne se transmet ni par la salive (concentration virale négligeable), ni par les piqûres de moustiques (virus non réplicatif dans le vecteur) ; il se transmet par le sang, les rapports sexuels non protégés et de la mère à l'enfant.
\item \textbf{Élève 2 : affirmation correcte.} Les ARV bloquent la multiplication du virus : bien suivis (observance stricte), ils maintiennent les défenses immunitaires (T4 élevés) et permettent une vie longue et presque normale (espérance de vie proche de la population générale).
\item \textbf{Élève 3 : affirmation fausse.} Une personne infectée peut se sentir en bonne santé pendant des années (phase asymptomatique, phase de latence clinique) tout en pouvant transmettre le virus ; le dépistage permet de connaître son statut, de se faire soigner tôt (meilleur pronostic) et de protéger les autres (prévention).
\item Les trois principales voies de transmission sont : 1) les rapports sexuels non protégés (voie sexuelle, $\approx 80\%$ des cas au Cameroun), 2) le sang contaminé (transfusions non contrôlées, partage d'objets piquants/tranchants : aiguilles, lames, scarifications), 3) la transmission de la mère à l'enfant (pendant la grossesse, l'accouchement, l'allaitement).
\item Moyens de prévention : l'utilisation correcte et systématique du \cle{préservatif} (masculin ou féminin), l'utilisation de matériel stérile à usage unique (aiguilles, seringues, lames de rasoir, matériel de tatouage/piercing) et le \cle{dépistage} volontaire (connaître son statut pour se protéger et protéger les autres). Le préservatif, par exemple, forme une barrière physique étanche (latex ou polyuréthane) qui empêche le virus présent dans les sécrétions sexuelles (sperme, pertes vaginales, sang) d'atteindre la muqueuse du partenaire.
\item Le traitement ARV de la mère (initié le plus tôt possible, idéalement avant la grossesse) fait fortement baisser la quantité de virus dans son sang (charge virale indétectable) et dans ses sécrétions : le risque que le virus passe au bébé pendant la grossesse (transplacentaire), l'accouchement (contact sanguin) ou l'allaitement (lait maternel) devient alors très faible (< 1-2\% avec prise en charge optimale PTME).
\item Exemple de slogan : « Protège-toi, fais-toi dépister : contre le VIH, l'information c'est la prévention ! » ou « Zéro contamination : je me protège, je me teste, je me soigne. »
\end{enumerate}
\textbf{Barème indicatif (sur 10) :} analyse des trois affirmations justifiées 3 pts ; voies de transmission 2 pts ; moyens de prévention et fonctionnement 2 pts ; protection de la mère et du bébé 2 pts ; slogan 1 pt.

\textbf{Point de vigilance :} pour chaque affirmation, la note complète exige à la fois le jugement (correcte/fausse) ET la justification scientifique précise (pas de vague « c'est pas vrai »).
\end{corrige}

\begin{corrige}[Exercice 11 — Type BEPC : problème de synthèse, multiplication du VIH et traitement]
\textbf{Rappel de l'énoncé :} Production virale $10^{10}$/jour. Calcul production 7 jours. Citer 4 étapes cycle. Expliquer épuisement T4. Justifier prise quotidienne ARV. Corriger idée « ARV guérissent / arrêt possible ».
\tcblower
\begin{enumerate}
\item Nombre de virus produits en 7 jours : $10^{10} \times 7 = 7 \times 10^{10}$ virus, soit soixante-dix milliards de nouveaux virus en une semaine.
\item Les quatre grandes étapes du cycle de réplication du VIH dans le lymphocyte T4 sont : 1) \textbf{Fixation et entrée} (reconnaissance CD4/co-récepteur, fusion) ; 2) \textbf{Transcription inverse} (synthèse d'ADN viral à partir de l'ARN par la transcriptase inverse) ; 3) \textbf{Expression du génome viral} (intégration, transcription, traduction : fabrication des protéines virales) ; 4) \textbf{Assemblage et libération} (bourgeonnement, maturation par la protéase).
\item Chaque jour, des milliards de nouveaux virus sont produits et chaque cycle de multiplication aboutit à la destruction (lyse) du lymphocyte T4 infecté. La destruction massive et continue des T4 finit par dépasser la capacité de régénération de l'organisme (thymus, prolifération périphérique) : le stock de T4 s'épuise progressivement et les défenses immunitaires s'effondrent (immunodépression).
\item Les ARV n'agissent que s'ils sont présents en permanence dans le sang à concentration inhibitrice. Les prendre tous les jours sans interruption maintient la réplication virale quasi nulle (charge virale indétectable), protège les lymphocytes T4 de la destruction et empêche la sélection de mutants \cle{résistants} aux médicaments.
\item C'est une erreur grave : les ARV bloquent la multiplication du virus mais ne l'éliminent pas totalement de l'organisme (réservoirs viraux : lymphocytes T4 mémoire, macrophages, système nerveux central, ganglions). Si le malade arrête le traitement, le virus repart des réservoirs, se remet à se multiplier, les T4 sont à nouveau détruits et la maladie progresse vers le Sida. Le traitement ARV doit donc être suivi \textbf{à vie}, avec une observance stricte ($> 95\%$).
\end{enumerate}
\textbf{Barème indicatif (sur 10) :} calcul en écriture scientifique 2 pts ; étapes du cycle 2 pts ; explication de l'épuisement des T4 2 pts ; intérêt de la prise quotidienne 2 pts ; correction de l'idée fausse 2 pts.

\textbf{Points de vigilance :} le résultat du calcul doit être donné en écriture scientifique ($7 \times 10^{10}$) ; ne jamais écrire que les ARV « tuent » ou « éliminent » le virus : ils \textbf{bloquent sa multiplication} (action virustatique, non virucide). Préciser « à vie » pour la durée du traitement.
\end{corrige}

\newpage

\coursec{Fiche bilan}

\begin{popfigure}
\centering
\begin{center}\tcbox[colback=popink,colframe=popink,coltext=white,arc=9pt,boxsep=4pt,left=6pt,right=6pt]{\bfseries\small VIH / Sida}\end{center}
\vspace{-2pt}
\begin{tcbraster}[raster columns=3,raster equal height=rows,raster column skip=6pt,raster row skip=6pt,enhanced,blankest]
\begin{tcolorbox}[enhanced,colback=white,colframe=poporange,boxrule=0.9pt,arc=3pt,title={1. Structure \& Cibles},fonttitle=\bfseries\scriptsize,coltitle=white,colbacktitle=poporange,left=4pt,right=4pt,top=2pt,bottom=2pt,boxsep=1pt,toptitle=1.5pt,bottomtitle=1.5pt,halign=left]
\begin{itemize}[nosep,leftmargin=*,itemsep=1pt,font=\scriptsize]
\item Enveloppe gp120/gp41
\item Capside + ARN + RT
\item Cibles : LT4 (CD4$^+$), Macrophages
\end{itemize}
\end{tcolorbox}
\begin{tcolorbox}[enhanced,colback=white,colframe=popgreen,boxrule=0.9pt,arc=3pt,title={2. Cycle Réplication},fonttitle=\bfseries\scriptsize,coltitle=white,colbacktitle=popgreen,left=4pt,right=4pt,top=2pt,bottom=2pt,boxsep=1pt,toptitle=1.5pt,bottomtitle=1.5pt,halign=left]
\begin{itemize}[nosep,leftmargin=*,itemsep=1pt,font=\scriptsize]
\item Fixation CD4/CCR5
\item RT : ARN $\to$ ADN
\item Intégration (Provirus)
\item Transcription/Traduction
\item Assemblage \& Bourgeonnement
\end{itemize}
\end{tcolorbox}
\begin{tcolorbox}[enhanced,colback=white,colframe=poppurple,boxrule=0.9pt,arc=3pt,title={3. Phases Cliniques},fonttitle=\bfseries\scriptsize,coltitle=white,colbacktitle=poppurple,left=4pt,right=4pt,top=2pt,bottom=2pt,boxsep=1pt,toptitle=1.5pt,bottomtitle=1.5pt,halign=left]
\begin{itemize}[nosep,leftmargin=*,itemsep=1pt,font=\scriptsize]
\item Primaire (fièvre, adénopathies)
\item Latence (asymptomatique)
\item Accélération (OI, amaigrissement)
\item Sida (OI graves, CD4 $<$ 200)
\end{itemize}
\end{tcolorbox}
\begin{tcolorbox}[enhanced,colback=white,colframe=poppink,boxrule=0.9pt,arc=3pt,title={4. Marqueurs Biologiques},fonttitle=\bfseries\scriptsize,coltitle=white,colbacktitle=poppink,left=4pt,right=4pt,top=2pt,bottom=2pt,boxsep=1pt,toptitle=1.5pt,bottomtitle=1.5pt,halign=left]
\begin{itemize}[nosep,leftmargin=*,itemsep=1pt,font=\scriptsize]
\item CD4 $\downarrow$ (immunité)
\item Charge Virale $\uparrow$ (réplication)
\item Sérologie : Ac anti-VIH (ELISA/Western Blot)
\end{itemize}
\end{tcolorbox}
\begin{tcolorbox}[enhanced,colback=white,colframe=popblue,boxrule=0.9pt,arc=3pt,title={5. Prévention \& Dépistage},fonttitle=\bfseries\scriptsize,coltitle=white,colbacktitle=popblue,left=4pt,right=4pt,top=2pt,bottom=2pt,boxsep=1pt,toptitle=1.5pt,bottomtitle=1.5pt,halign=left]
\begin{itemize}[nosep,leftmargin=*,itemsep=1pt,font=\scriptsize]
\item ABC : Abstinence, Fidélité, Préservatif
\item PTME : Mère $\to$ Enfant
\item Dépistage : TROD, Labo (Confidentiel, Gratuit)
\item PEP : Post-Exposition (72h)
\end{itemize}
\end{tcolorbox}
\begin{tcolorbox}[enhanced,colback=white,colframe=popgold,boxrule=0.9pt,arc=3pt,title={6. Traitement ARV (Cameroun)},fonttitle=\bfseries\scriptsize,coltitle=white,colbacktitle=popgold,left=4pt,right=4pt,top=2pt,bottom=2pt,boxsep=1pt,toptitle=1.5pt,bottomtitle=1.5pt,halign=left]
\begin{itemize}[nosep,leftmargin=*,itemsep=1pt,font=\scriptsize]
\item Trithérapie (3 molécules)
\item Cibles : RT, Intégrase, Protéase
\item Objectif : CV indétectable, CD4 $\uparrow$
\item Observance stricte = Vie longue
\end{itemize}
\end{tcolorbox}
\end{tcbraster}

\legende{Carte mentale récapitulative de la leçon sur le VIH/Sida (Programme 3ème SVTEEHB).}
\end{popfigure}

\begin{bilan}

\courssub{Définitions Clés}
\begin{itemize}
    \item \cle{VIH} (Virus de l'Immunodéficience Humaine) : Rétrovirus à ARN, enveloppé, de la famille des \emph{Lentivirus}. Deux types : VIH-1 (mondial, virulent) et VIH-2 (Afrique de l'Ouest, moins virulent).
    \item \cle{Sida} (Syndrome d'Immunodéficience Acquise) : Stade terminal de l'infection VIH non traitée, défini par une numération des LT4 $< 200$/mm$^3$ ou la survenue d'infections opportunistes (IO) graves.
    \item \cle{LT4 (CD4$^+$)} : Lymphocytes T auxiliaires, \textbf{cellules cibles principales} du VIH (récepteur CD4 + co-récepteurs CCR5/CXCR4). Leur chute entraîne l'effondrement de l'immunité cellulaire.
    \item \cle{Charge Virale (CV)} : Quantité d'ARN viral par mL de plasma. Marqueur de l'activité de réplication virale. Objectif ARV : \textbf{indétectable} ($< 50$ copies/mL).
    \item \cle{Séroconversion} : Apparition des anticorps anti-VIH dans le sang (délai fenêtre : 3 à 12 semaines post-infection).
    \item \cle{ARV} (Antirétroviraux) : Médicaments bloquant la réplication virale. Utilisés en \textbf{trithérapie} (association de 3 molécules de 2 classes minimum).
\end{itemize}

\courssub{Mécanismes \& Étapes à Connaître par Cœur}
\begin{center}
\begin{tabularx}{\linewidth}{|l|X|}\hline
\textbf{Étape Cycle VIH} & \textbf{Description \& Cible Thérapeutique} \\ \hline
1. \textbf{Fixation / Entrée} & gp120 se lie au CD4 puis co-récepteur (CCR5/CXCR4) $\to$ fusion membranes. \textit{Cible : Inhibiteurs d'entrée (ex. Maraviroc, Enfuvirtide).} \\ \hline
2. \textbf{Transcription Inverse} & \cle{Transcriptase Inverse (RT)} : ARN viral $(+)$ $\to$ ADN double brin (ADN proviral). \textit{Cible : INRT (TDF, 3TC, AZT) \& INNRT (EFV, NVP, DTG).} \\ \hline
3. \textbf{Intégration} & \cle{Intégrase} : Insère l'ADN proviral dans l'ADN hôte (chromosome). \textit{Cible : II (DTG, RAL, BIC).} \\ \hline
4. \textbf{Transcription / Traduction} & Machinery cellulaire $\to$ ARNm viraux $\to$ Protéines polyprotéines (Gag, Pol, Env). \\ \hline
5. \textbf{Maturation / Assemblage} & \cle{Protéase} : Clive les polyprotéines en protéines fonctionnelles. \textit{Cible : IP (LPV/r, ATV/r, DRV/r).} \\ \hline
6. \textbf{Bourgeonnement} & Nouveau virion infectieux sort de la cellule (lyse ou bourgeonnement). \\ \hline
\end{tabularx}
\end{center}

\courssub{Dynamique Infection : Phases \& Marqueurs (Repères Chiffrés Ordres de Grandeur)}
\begin{center}
\begin{tabularx}{\linewidth}{|c|C|C|C|C|}\hline
\textbf{Phase} & \textbf{Durée approx.} & \textbf{CD4 (cellules/mm$^3$)} & \textbf{Charge Virale} & \textbf{Clinique} \\ \hline
Primaire & 2--6 sem & Chute transitoire & Pic très élevé ($> 10^6$) & Syndrome pseudo-grippal / mononucléosique \\ \hline
Latence (Asymptomatique) & Années (5--10 ans) & Déclin lent ($\sim 50$/an) & Plateau (point de consigne) & Aucune / Adénopathies généralisées persistantes \\ \hline
Accélération (Symptomatique) & Mois--Années & $200$--$500$ & Remontée & Infections bactériennes récidivantes, zona, amaigrissement \\ \hline
\textbf{Sida (Stade 4 OMS)} & Mois & \textbf{$< 200$} & Très élevée & \textbf{IO graves :} Pneumocystose, Toxoplasmose, Tuberculose, Cryptococcose, Lymphome \\ \hline
\end{tabularx}
\end{center}

\courssub{Méthodes Express}
\begin{methode}[Interpréter un résultat de sérologie VIH]
\begin{enumerate}
    \item \textbf{TROD / Autotest positif} $\Rightarrow$ \textbf{Séropositif probable}. \textbf{Obligatoire} : Confirmation au laboratoire (ELISA + Western Blot ou 2 ELISA différents).
    \item \textbf{Négatif} + \textbf{Prise de risque $< 3$ mois} $\Rightarrow$ \textbf{Refaire le test} après délai fenêtre (3 mois). PEP si $< 72$h.
    \item \textbf{Négatif} + \textbf{Prise de risque $> 3$ mois} $\Rightarrow$ \textbf{Séronégatif} (non infecté).
\end{enumerate}
\end{methode}

\begin{methode}[Expliquer l'action des ARV (Trithérapie standard 1ère ligne Cameroun : TDF + 3TC + DTG)]
\begin{enumerate}
    \item \textbf{TDF / 3TC (INRT)} : « Fausses briques » (analogues nucléosidiques) $\to$ bloquent la Transcriptase Inverse $\to$ arrêt synthèse ADN viral.
    \item \textbf{DTG (II - Inhibiteur Intégrase)} : Bloque l'Intégrase $\to$ empêche l'insertion de l'ADN viral dans le génome hôte.
    \item \textbf{Résultat synergique} : Réplication bloquée à 2 étapes $\to$ Charge Virale $\downarrow\downarrow$ (indétectable à 6 mois) $\to$ CD4 $\uparrow\uparrow$ (reconstitution immunitaire).
\end{enumerate}
\end{methode}

\begin{methode}[Argumenter la prévention combinée au Cameroun]
\begin{enumerate}
    \item \textbf{Biomédicale} : Dépistage universel (ANC, TB, STI), PTME (Option B+ : ARV vie toute grossesse/allaitement), PEP (viol, accident exposition sanguine), PrEP (populations clés).
    \item \textbf{Comportementale} : Éducation sexuelle complète (ESC), promotion préservatif (M/F), réduction partenaires, circoncision médicale volontaire (CMV).
    \item \textbf{Structurelle} : Lutte contre stigmatisation (droits humains), maintien à l'école jeunes filles, autonomisation économique.
\end{enumerate}
\end{methode}

\end{bilan}

\begin{aretenir}[Checklist avant le BEPC]
    $\square$ Je sais décrire la structure du VIH (enveloppe gp120/gp41, capside, ARN, enzymes RT/Intégrase/Protéase) et nommer ses cellules cibles (LT4 CD4$^+$, macrophages).\par
    $\square$ Je connais par cœur les 6 étapes du cycle de réplication (Fixation, Transcription Inverse, Intégration, Transcription/Traduction, Maturation, Bourgeonnement) et la classe d'ARV qui bloque chaque étape.\par
    $\square$ Je sais expliquer pourquoi le VIH est un \textbf{rétrovirus} (enzyme Transcriptase Inverse : ARN $\to$ ADN).\par
    $\square$ Je sais tracer et interpréter la courbe dynamique : CD4 $\downarrow$ / Charge Virale $\uparrow$ au cours des 4 phases (Primaire, Latence, Accélération, Sida).\par
    $\square$ Je connais les critères biologiques du Sida : CD4 $< 200$/mm$^3$ ET/OU Infections Opportunistes (IO) définissant le stade 4 OMS.\par
    $\square$ Je sais définir la séroconversion, le délai fenêtre (3 mois) et interpréter un algorithme de dépistage (TROD $\to$ Confirmation Labo).\par
    $\square$ Je connais la stratégie PTME (Option B+) au Cameroun : ARV vie pour toute femme enceinte/allaitante séropositive = élimination transmission Mère-Enfant.\par
    $\square$ Je sais expliquer le principe de la trithérapie (3 molécules, 2 cibles minimum) et l'importance vitale de l'\textbf{observance} ($> 95\%$) pour éviter la résistance.\par
    $\square$ Je sais citer les 3 piliers de la prévention combinée : Biomédicale (Dépistage, ARV, PrEP/PEP), Comportementale (ABC, CMV), Structurelle (Droits, Éducation).\par
    $\square$ Je sais rédiger une conclusion argumentée : « Grâce aux ARV, le VIH devient une maladie chronique contrôlée ; une personne sous traitement efficace (CV indétectable) ne transmet plus le virus (I=I : Indétectable = Intransmissible). »\par
    $\square$ Je connais les numéros utiles au Cameroun : \textbf{117} (Urgences), \textbf{1510} (Info VIH/Sida - CNLS), structures de prise en charge (PEC) décentralisées (Hôpitaux de District, CPEC).
\end{aretenir}

\findecours

\end{document}
